\section{Base lattices I: closing the member construction}\label{sec:pf:baselat1}

\providecommand{\baselatoneS}{\mathcal S}
\providecommand{\baselatoneP}{\mathfrak p}
\providecommand{\baselatoneO}{\mathcal O}
\providecommand{\baselatoneLam}{\Lambda}
\providecommand{\baselatonenv}{\operatorname{nv}}
\providecommand{\baselatoneun}{T^{\mathrm{untw}}}

\subsection{Introduction}\label{pf:baselat1:intro}

This section proves \cref{pf:baselat1:main} below. It closes the \emph{member construction} of the
development: from two norm-choice extensions $N$ and $N_g$ (the ``pair''), the construction of levels
three and four of column witnesses, and a list of published results, it produces the member-existence
statement $\sfname{MemberE}$ and the bundle-existence statement $\sfname{ExBL}$ of
\cref{sec:families} for the parameter attached to $N$. The theorem has exactly one user,
\lref{GlobalTheory.Mem.Top.Bl.exists_baseLattice_of_prime_alt}, which applies it to the ``origin''
extension and its refinement by the goodness condition $\mathrm{Gd}_0$; through that result it enters the
proof of \cref{thm:member-bundles}. Neither a lattice nor a prime occurs in the statement; the formal name
of the theorem refers to its role in the development.

\paragraph{The statement.} Member existence $\sfname{MemberE}(\baselatoneS)$ for a slot family
$\baselatoneS$ is a conjunction of four statements about \emph{cores}, records of construction data
attached to a base $B$, a frame $f$, a twisting datum $T$ and a choice of stages
(\cref{pf:baselat1:defs}):
\begin{enumerate}[label=(M\arabic*)]
\item bases with prescribed coefficient fields exist, eventually in the scale (\emph{base condition});
\item eventually in the scale and for all sufficiently large untwisted twisting data, a core built for
  $\baselatoneS$ exists (\emph{first step});
\item eventually, every built core over untwisted data satisfies the \emph{export condition}
  $\mathcal{X}$, a statement about Bloch--Kato Selmer groups of twists of the Galois
  representations attached to the core (\emph{export});
\item eventually, every exported built core can be moved, by enlarging the base and retwisting with the
  same set of primes, to a core which is built, good and twisted at all places (\emph{second pass}).
\end{enumerate}
Bundle existence $\sfname{ExBL}$ is quoted from \lref{GlobalTheory.T.geometric_parameter_uniform_alt}.

\paragraph{The proof.} (M3) is a theorem in Galois cohomology (\cref{pf:baselat1:export}). For every
twisted place one needs a finite set $Q$ of auxiliary split primes and a character $\psi$ of a ray class
group of conductor $Q$, with values in $(\baselatoneO/\ell^{k})^\times$, such that for \emph{every} power
$\psi^a$ with $\ell\nmid a$ the Bloch--Kato Selmer groups of the corresponding twists have normalised
length at most $\nu+\varepsilon$, the theta values attached to $\psi^a$ have controlled valuation, and the
tame ramification of $\psi^a$ is small. The development constructs $\psi$ as follows. Over the group ring
$\baselatoneLam_M=\baselatoneO_0[(\Z/\ell^M)^{2R}]$, which receives the Galois group through $2R$
characters of order $\ell^M$ ramified at auxiliary primes chosen by Chebotarev's theorem, the Selmer
modules of the universal twist are presented by a square matrix $\Phi$; the length of the Selmer group of a
specialisation is the valuation of $\det\Phi$ at the specialisation, and $\det\Phi$ is the image of a
polynomial $F$ in $2R$ variables having a coefficient of bounded degree and valuation at most
$\nu\, v(\ell)$. The character is then obtained by evaluating $F$ at a
point $(\zeta^{\ell^{e_q}}-1)_q$, $\zeta$ a primitive $\ell^k$-th root of unity, chosen by a Newton polygon
argument (\cref{pf:baselat1:ssec:roots}); this part is proved in full. Of the Selmer-theoretic inputs (existence
of the presentation, the determinant formula and the coefficient bounds), the coefficient bound for positive
exponent is proved in \cref{pf:baselat1:sssec:boundproofs}; the others are quoted.

(M1), (M2) and (M4) are formal consequences (\cref{pf:baselat1:steps2}) of the following construction
steps, each of which is an ``eventual'' statement over column witnesses: level three for starting witnesses
(a hypothesis), level four from level three (\cref{pf:baselat1:has4}), ambient, division and theta data at
level four (\cref{pf:baselat1:ssec:stages}), and the assembly of a core from these data together with a
lattice layer, a kernel datum and a gauge family of Haar measures (\cref{pf:baselat1:core}). For (M4) the
same construction is run for $N_g$ on the witnesses produced by the given second pass, and cores built for
$N_g$ are good for $N$ (\cref{pf:baselat1:goodE}).

\paragraph{Organisation.} \Cref{pf:baselat1:ssec:framework} recalls the framework of eventual statements,
steps, passes, cores and slots. \Cref{pf:baselat1:ssec:calculus} proves the calculus of steps and passes
and the reduction of $\sfname{MemberE}$ to construction steps. \Cref{pf:baselat1:ssec:stages} treats the
construction steps and proves (M1), (M2), (M4). \Cref{pf:baselat1:ssec:export} proves (M3), and
\cref{pf:baselat1:ssec:roots} the evaluation theorem at roots of unity used there.
\Cref{pf:baselat1:ssec:deferred} proves the main construction results quoted in the two preceding parts: gauge
families, lattice layers, division and theta data, spectral objects, Hida settings and spectral data with gap, the
inputs of the twist theorem and the coefficient bound for positive exponent. \Cref{pf:baselat1:ssec:main} proves the
main theorem, and \cref{pf:baselat1:ssec:remarks} collects remarks
on the formal statements.

Throughout we use short notation for the notions of the formal development \cite{LeanSource2026}; the formal
names are given in the tags at the end of the definitions and statements. We unfold the definitions as far as
the arguments require. Many of them are large records whose fields are only used through the statements
quoted here; we do not describe such fields further.

%======================================================================================================
\subsection{The framework}\label{pf:baselat1:ssec:framework}

\subsubsection{Scale contexts and eventual statements}

A \emph{scale context} $K$ consists of a scale $K.s$, with real parameters $x=K.s.x>1$
and $\eps=K.s.\varepsilon$ and derived scale functions such as $X''$, a standard constant record $K.C$, two
reals $c_D=K.c_D$ and $c_h=K.c_h$, and arithmetic data (a pre-character datum with level data) subject to
compatibility conditions; $\mathrm{good}_K(C_5,c_6,C_b)$ is a goodness condition depending on three real
constants. A scale context lies over certain \emph{pre-bases} $P$, written $K.\mathrm{over}(P)$. None of
these notions is unfolded in this section.

\begin{definition}[Eventual statements]\label{pf:baselat1:def:evg}
For a predicate $p$ on scale contexts, $\mathrm{Ev}(p)$ means: there is $\eps_0>0$ such that for all
reals $\eps,C_8,C_5,c_6,C_b$ with $0<\eps\le\eps_0$, $\eps<1/52$ and $C_8\ge1$ there is a transported
constant $K_0$ such that for every standard constant record $C$ with $C.C_{\mathrm{pde}}=C_8$ which is
$K_0$-large there is $c_1$ such that for all $c_D\ge c_1$ there is $c_2$ such that for all $c_h\ge c_2$
there is $x_0$ with: $p(K)$ holds for every scale context $K$ with $K.C=C$, $K.c_D=c_D$, $K.c_h=c_h$,
$K.s.\eps=\eps$, $K.s.x\ge x_0$ and $\mathrm{good}_K(C_5,c_6,C_b)$. The variant $\mathrm{Ev}^{\circ}(p)$ omits the
constants $C_5,c_6,C_b$ and the goodness condition.
\formal{GlobalTheory.Mem.Top.EvG}\formal{GlobalTheory.Mem.Top.Ev}
\end{definition}

\begin{lemma}\label{pf:baselat1:evg-inter}
\begin{enumerate}[label=(\roman*)]
\item If $\mathrm{Ev}(p)$ and $\mathrm{Ev}(q)$, then $\mathrm{Ev}(p\wedge q)$.
\item If $\mathrm{Ev}(p)$ and $p(K)\Rightarrow q(K)$ for every scale context $K$ satisfying
  $\mathrm{good}_K(C_5,c_6,C_b)$ for some constants, then $\mathrm{Ev}(q)$.
\end{enumerate}
\formallabel{GlobalTheory.Mem.Top.EvG.inter}\formallabel{GlobalTheory.Mem.Top.EvG.imp}
\end{lemma}

\begin{proof}
(i) Take the minimum of the two thresholds $\eps_0$, the supremum of the two transported constants (a record
$C$ which is large for the supremum is large for both), and the maxima of the two thresholds $c_1$, of the two
thresholds $c_2$ and of the two thresholds $x_0$. (ii) is immediate from the definition, the quantifiers
being the same.
\end{proof}

If $p(K)$ holds for every scale context then $\mathrm{Ev}(p)$
(\lref{GlobalTheory.Mem.Top.GroupRetwisting.evG_all}).

\subsubsection{Bases, twisting data and large data}

A \emph{pre-base} $P$ fixes ranks $P.R$, a CM field $E$ with auxiliary primes $g_0,g_1,d$
(the \emph{host}), hermitian spaces and words with level primes. A \emph{base} $B$ over a
pre-base $B.\mathrm{pre}$ adds a coefficient field $\mathit{LL}$ with an embedding
$\iota_B\colon\mathit{LL}\to\C$ and subfields $K_s\subseteq\mathit{LL}_0\subseteq\mathit{LL}$; we write
$B.\mathit{LLC}=\iota_B(\mathit{LL})$, $B.K_s^{\C}=\iota_B(K_s)$ and $B.F^{\mathrm{fix}}=\iota_B(\mathit{LL}_0)$
(the ``fixed coefficient field''). A \emph{frame} $f$ of a pre-base is a Hecke-character frame of its host.
For bases $B,B'$, $B'$ \emph{enlarges} $B$, $\mathrm{Enl}(B,B')$, if $B'.\mathrm{pre}=B.\mathrm{pre}$,
$\iota_{B'}(K_s')=\iota_B(K_s)$, $\iota_{B'}(\mathit{LL}_0')=\iota_B(\mathit{LL}_0)$ and
$\iota_B(\mathit{LL})\subseteq\iota_{B'}(\mathit{LL}')$; then also $B'.F^{\mathrm{fix}}=B.F^{\mathrm{fix}}$ and
$B'.K_s^{\C}=B.K_s^{\C}$.

A \emph{twisting datum} $T$ over $B$ (\cref{def:families:twist}) consists of a finite set
$T.\mathrm{pl}$ of rational primes and, for each place $w$ of $\mathit{LL}$ with residue characteristic
$\ell_w$, a finite set $T.Q_w$ of primes of $\OO_E$ not above $\ell_w$, an exponent $T.k_w\in\N$ and a
Hecke character $T.\psi_w$ of $E$ of order dividing $\ell_w^{T.k_w}$ (subject to further conditions); it
determines a set $T.\mathrm{set}(f)$ of Hecke characters of $E$ for every frame $f$. The \emph{untwisted}
datum $\baselatoneun(B,\mathrm{pl})$ has prime set $\mathrm{pl}$, $Q_w=\emptyset$, $k_w=0$ and
$\psi_w=1$; $T$ is \emph{untwisted} if $T=\baselatoneun(B,T.\mathrm{pl})$. $T$ \emph{avoids}
$A\colon\N\to\{\text{finite subsets of }\N\}$ if no $v\in T.Q_w$ has absolute norm in $A(\ell_w)$. $T$ is
\emph{forced} for $K$, $\mathrm{Fd}(K,B,T)$, if every prime $p$ with $p\in\{g_0,g_1,d\}$ or
$p\le X''(n)$ (with $n$ a rank of $B$ and $X''$ the scale function of $K$) belongs to $T.\mathrm{pl}$.

\begin{definition}[Large data]\label{pf:baselat1:def:large}
For a scale context $K$ and a predicate $p(B,f,T)$ on bases, frames and twisting data,
$\mathrm{Lg}_2(K,p)$ means: there are functions $S_1,P_1,A,S_2$ assigning to a pre-base, a frame and
one or two subfields of $\C$ (and, for $S_2$, a set of Hecke characters) a finite set of algebraic numbers
($S_1$, $S_2$), a finite set of natural numbers ($P_1$), and a function $\N\to\{\text{finite subsets of
}\N\}$ ($A$), such that $p(B,f,T)$ holds for all $B,f,T$ with $K.\mathrm{over}(B.\mathrm{pre})$,
$S_1(B.\mathrm{pre},f,B.K_s^{\C})\subseteq B.F^{\mathrm{fix}}$,
$S_2(B.\mathrm{pre},f,B.K_s^{\C},B.F^{\mathrm{fix}},T.\mathrm{set}(f))\subseteq B.\mathit{LLC}$,
$P_1(B.\mathrm{pre},f,B.K_s^{\C},B.F^{\mathrm{fix}})\subseteq T.\mathrm{pl}$, and $T$ avoiding
$A(B.\mathrm{pre},f,B.K_s^{\C},B.F^{\mathrm{fix}})$. We put
$\sfname{EvLgG}(p)=\mathrm{Ev}(K\mapsto\mathrm{Lg}_2(K,p_K))$ for predicates $p(K,B,f,T)$:
\emph{eventually in the scale, $p$ holds for all sufficiently large data}.
\formal{GlobalTheory.Mem.Top.large₂}\formal{GlobalTheory.Mem.Top.EvLgG}
\end{definition}

\begin{lemma}\label{pf:baselat1:large}
Let $K$ be a scale context. (i) If $\mathrm{Lg}_2(K,p)$ and $\mathrm{Lg}_2(K,q)$ then
$\mathrm{Lg}_2(K,p\wedge q)$. (ii) $\mathrm{Lg}_2(K,(B,f,T)\mapsto\mathrm{Fd}(K,B,T))$.
\formallabel{GlobalTheory.Mem.Top.large₂.inter}\formallabel{GlobalTheory.Mem.Top.GroupRetwisting.lg_forced}
\end{lemma}

\begin{proof}
(i) Take the unions of the two families of finite sets $S_1$, $S_2$, $P_1$ and the pointwise unions of the
two functions $A$. (ii) Take $S_1=S_2=\emptyset$, $A=\emptyset$ and
$P_1(P)=\{g_0,g_1,d\}\cup\{0,1,\dots,\lfloor X''(P.R.n)\rfloor\}$; every prime $p$ with $p\in\{g_0,g_1,d\}$
or $p\le X''(n)$ lies in $P_1(P)$.
\end{proof}

\subsubsection{Steps, passes and the base condition}

A \emph{base predicate} is a predicate $\kappa(B,f)$; an \emph{at-predicate} is a predicate
$p(K,B,f,T)$; a \emph{witness predicate} is a predicate $p(K,B,f,T,w)$ in which $w$ is a column witness
for $(B,f,T)$ (see below). The base predicate used throughout is $\kappa_0$: $\kappa_0(B,f)$ holds if the
pre-base of $B$ admits a field datum $\delta$ for which the field $K_s$ of $B$ is exact with respect to
$(\delta,f)$, and the pre-base of $B$ is pinned (these two conditions are not unfolded here).

\begin{definition}[Steps, passes, base condition]\label{pf:baselat1:def:step}
Let $\kappa$ be a base predicate and $p,q,r$ at-predicates.
\begin{enumerate}[label=(\alph*)]
\item $\mathrm{Step}(\kappa,p,q)$ is $\sfname{EvLgG}$ of
  $(K,B,f,T)\mapsto[\kappa(B,f)\Rightarrow\mathrm{Fd}(K,B,T)\Rightarrow p(K,B,f,T)\Rightarrow
  q(K,B,f,T)]$.
\item For witness predicates $p,q$:
  \[\mathrm{Step}^{\mathrm w}(p,q)=\mathrm{Step}\bigl(\kappa_0,\mathrm{true},(K,B,f,T)\mapsto[\forall w:\
  p(K,B,f,T,w)\Rightarrow q(K,B,f,T,w)]\bigr).\]
\item For a slot $Q$ (\cref{pf:baselat1:def:core}), $\mathrm{Pass}(\kappa,Q,r)$ is $\sfname{EvLgG}$ of the
  following statement about $(K,B,f,T)$: if $\kappa(B,f)$ and $T$ is untwisted, then for all stages
  $C,W,V,S,N,G$ over $B$ and every core $c$ over them with $Q(c)$ and
  $\mathcal{X}(\mathrm{asm}(B,C,W,V,S,N,G),c.O.\mathrm{nch},c.O.x_d)$, every
  $A\colon\N\to\{\text{finite subsets of }\N\}$ and every assignment $S_2$ of finite sets of algebraic numbers
  to (pre-base, frame, set of Hecke characters), there are a base $B'$ with $\mathrm{Enl}(B,B')$, the frame
  $f'=f$ of the same pre-base, and a twisting datum $T'$ over $B'$ with $T'.\mathrm{pl}=T.\mathrm{pl}$,
  $T'$ avoiding $A$, $S_2(B'.\mathrm{pre},f',T'.\mathrm{set}(f'))\subseteq B'.\mathit{LLC}$, and
  $r(K,B',f',T')$.
\item $\mathrm{BC}(\kappa)$ means: (ex) $\mathrm{Ev}$ of the statement ``there are a pre-base $P$ and
  a frame $f_0$ of $P$ with $K.\mathrm{over}(P)$ such that for all assignments $S_1,S_2$ of finite sets of
  algebraic numbers to one, resp.\ two, subfields of $\C$ there are a base $B$ over $P$ and the frame $f=f_0$
  with $\kappa(B,f)$, $S_1(B.K_s^{\C})\subseteq B.F^{\mathrm{fix}}$ and
  $S_2(B.K_s^{\C},B.F^{\mathrm{fix}})\subseteq B.\mathit{LLC}$''; and (enl) $\kappa$ is stable under
  enlargement: $\mathrm{Enl}(B,B')$ and $\kappa(B,f)$ imply $\kappa(B',f)$.
\end{enumerate}
\formal{GlobalTheory.Mem.Top.GroupRetwisting.Step}\formal{GlobalTheory.Mem.Top.Bl.StepW}%
\formal{GlobalTheory.Mem.Top.GroupRetwisting.Pass}\formal{GlobalTheory.Mem.Top.GroupRetwisting.BaseCond}%
\formal{GlobalTheory.Mem.Top.Bl.KPx}
\end{definition}

\begin{lemma}[Calculus of steps]\label{pf:baselat1:stepcalc}
Let $\kappa$ be a base predicate, $p,q,r$ at-predicates and $p',q',r'$ witness predicates.
\begin{enumerate}[label=(\roman*)]
\item If $\kappa(B,f)$, $\mathrm{Fd}(K,B,T)$ and $p(K,B,f,T)$ imply $q(K,B,f,T)$ whenever
  $K.\mathrm{over}(B.\mathrm{pre})$, then $\mathrm{Step}(\kappa,p,q)$.
\item $\mathrm{Step}(\kappa,p,q)$ and $\mathrm{Step}(\kappa,p,r)$ imply $\mathrm{Step}(\kappa,p,q\wedge r)$
  \textup{(\lref{GlobalTheory.Mem.Top.GroupRetwisting.Step.and})}.
\item $\mathrm{Step}(\kappa,p,q)$ and $\mathrm{Step}(\kappa,q,r)$ imply $\mathrm{Step}(\kappa,p,r)$
  \textup{(\lref{GlobalTheory.Mem.Top.GroupRetwisting.Step.trans})}.
\item If $p'(w)\Rightarrow q'(w)$ for all $K,B,f,T,w$ with $K.\mathrm{over}(B.\mathrm{pre})$ and
  $\mathrm{Fd}(K,B,T)$, then $\mathrm{Step}^{\mathrm w}(p',q')$.
\item $\mathrm{Step}^{\mathrm w}(p',q')$ and $\mathrm{Step}^{\mathrm w}(p',r')$ imply $\mathrm{Step}^{\mathrm w}(p',q'\wedge r')$
  \textup{(\lref{GlobalTheory.Mem.Top.Bl.StepW.and})}; $\mathrm{Step}^{\mathrm w}(p',q')$ and
  $\mathrm{Step}^{\mathrm w}(q',r')$ imply $\mathrm{Step}^{\mathrm w}(p',r')$ \textup{(\lref{GlobalTheory.Mem.Top.Bl.StepW.trans})}.
\end{enumerate}
\formallabel{GlobalTheory.Mem.Top.GroupRetwisting.Step.of_all}\formallabel{GlobalTheory.Mem.Top.Bl.StepW.of_all}
\end{lemma}

\begin{proof}
(i) By \lref{GlobalTheory.Mem.Top.GroupRetwisting.evG_all} it suffices that
$\mathrm{Lg}_2(K,\cdot)$ holds for every $K$ for the implication in question; this is the case with all
exceptional sets empty, because the implication holds for every $(B,f,T)$ over $K$. (iv) is (i) applied to
the at-predicate $\forall w,\ p'(w)\Rightarrow q'(w)$. Statements (ii), (iii) and (v) are proved in the cited
places by \cref{pf:baselat1:evg-inter} and \cref{pf:baselat1:large}(i): the conjunction of two eventual
large statements is again one, and the implications compose pointwise.
\end{proof}

%------------------------------------------------------------------------------------------------------
\subsubsection{Stages, cores and slots}\label{pf:baselat1:defs}

\begin{definition}[Stages, cores, slots]\label{pf:baselat1:def:core}
\begin{enumerate}[label=(\alph*)]
\item A \emph{stage tuple} $X=(B,C,W,V,S,N,G)$ consists of a base $B$ and, built on it, column data
  $C\in\mathrm{Cols}(B)$, a row rest $W$, a division stage $V$, an $S$-stage $S$, a norm stage $N$ and a
  $GJ$-stage $G$. Its \emph{assembled datum} $\mathrm{asm}(B,C,W,V,S,N,G)$ is a global datum
  (\cref{sec:families}).
\item A \emph{core} $c\in\mathrm{Core}(K,X,f,T)$ is a record over $(K,X,f,T)$; among its components are a
  set of places $c.W_t$, a character tuple $c.\xi$, and an \emph{object} $c.O$ whose components include a
  norm choice $c.O.\mathrm{nch}$, theta data $c.O.x_d$, a base-level row side $c.O.p_0$, division data
  $c.O.X$, kernel data and a gauge family $c.O.\mathrm{gauge}$ (\cref{pf:baselat1:def:gauge}).
\item A \emph{slot} is a predicate $Q(K,X,f,T,c)$ on cores. For a slot $Q$, $\mathrm{Ex}_{\mathrm{core}}(Q)(K,B,f,T)$
  means that there are stages $C,W,V,S,N,G$ over $B$ and a core $c$ over $K$, the stage tuple
  $(B,C,W,V,S,N,G)$, $f$ and $T$ with $Q(c)$.
\item A \emph{column witness} $w\in\mathcal{W}(B,f,T)$ is a record of column input data for $(B,f,T)$
  whose frame is $f$ and whose characters range over $T.\mathrm{set}(f)$. It has a type $w.\mathrm{Lt}$ of
  \emph{levels}; each level $Z$ determines stages $Z.W,Z.V,Z.S,Z.N,Z.G$ over the column data $w.C$, hence
  a stage tuple $w.\mathrm{st}(Z)$ and a global datum $w.D(Z)=\mathrm{asm}(w.\mathrm{st}(Z))$.
\item A \emph{core predicate} is a predicate $P(K,B,f,T,w,Z,c)$, where $c$ is a core over
  $(K,w.\mathrm{st}(Z),f,T)$. The \emph{position} of $P$ is the slot $\mathrm{pos}(P)$ given by
  \begin{multline*}\mathrm{pos}(P)(K,X,f,T,c)\ :\Longleftrightarrow\ \exists\,w,Z,c':\\
  P(K,X.B,f,T,w,Z,c')\ \text{and}\ (X.C,(X.W,\dots,X.G),c)=(w.C,Z,c')\end{multline*}
  (an equality of elements of a dependent sum). We put $\mathrm{Ex}_{Z}(P')(w)\Leftrightarrow\exists Z,\
  P'(w,Z)$ and $\mathrm{Ex}_{Z,c}(P)(w)\Leftrightarrow\exists Z,c,\ P(w,Z,c)$.
\end{enumerate}
\formal{GlobalTheory.Mem.Top.V4.Core}\formal{GlobalTheory.Mem.Top.GroupRetwisting.HasCore}%
\formal{GlobalTheory.Mem.Top.Bl.ColW}\formal{GlobalTheory.Mem.Top.Bl.positionOf}
\end{definition}

\begin{lemma}\label{pf:baselat1:posintro}
If $P(K,B,f,T,w,Z,c)$ then $\mathrm{pos}(P)(K,w.\mathrm{st}(Z),f,T,c)$. Conversely, if
$P(K,B,f,T,w,Z,c)\Rightarrow R(K,w.\mathrm{st}(Z),f,T,c)$ for all arguments, then
$\mathrm{pos}(P)\Rightarrow R$ \textup{(\lref{GlobalTheory.Mem.Top.Bl.positionOf_ind})}.
\formallabel{GlobalTheory.Mem.Top.Bl.positionOf_intro}
\end{lemma}

\begin{proof}
Take $(w,Z,c)$ itself as witness; the required equality holds by definition of $w.\mathrm{st}(Z)$. The
converse is transport along the equality of dependent pairs.
\end{proof}

\begin{definition}[Slot families; built and good cores]\label{pf:baselat1:def:built}
A \emph{slot family} $\baselatoneS=(R_x,N_x,S_x,O_x)$ consists of four type-valued requirements: on row data
($R_x$), on norm-choice data ($N_x$), on spectral data ($S_x$) and on objects ($O_x$). A core $c$ at level $Z$
of $w$ is \emph{built} for $\baselatoneS$ if there is a record $b\in\mathcal{B}(\baselatoneS,K,w,Z,c)$,
consisting of: a column pin of $w$; coherence $\mathrm{Coh}(B,T)$; a group identification; column data
$\mathrm{col}$ at split primes; row data $t\in\mathcal{R}(K,w,Z,c.W_t,c.O.p_0)$ with an element of
$R_x(t)$; the equalities $c.\xi=w.I.\xi$, $c.O.\mathrm{lev}=\mathrm{lev}_0(\mathrm{col})$ and
$c.O.\mathrm{err}=\mathrm{err}(\mathrm{col})$; a norm-choice datum $n\in\mathcal{N}(w,Z,c.W_t,c.O.\mathrm{nch})$
with an element of $N_x(t,c.O.\mathrm{nch},n,c.O.x_d)$; an element of $S_x(t,\mathrm{col})$; a kernel family
compatible with the kernel data of $c.O$; and kernel, division, theta and object records
$\mathrm{Rec}_{\ker}$, $\mathrm{Rec}_{\mathrm{div}}$, $\mathrm{Rec}_{\theta}$, $\mathrm{Rec}_{\mathrm{ob}}$, together with an element of $O_x(c.O,t)$.
Let $\mathrm{Ent}(b.\mathrm{ker})$ be the condition that the kernel entries of $b$ are the canonical ones. We put
\begin{align*}
\mathrm{Bu}(\baselatoneS)(w,Z,c)&:\Longleftrightarrow\exists\,b\in\mathcal{B}(\baselatoneS,K,w,Z,c):\
  \mathrm{Ent}(b.\mathrm{ker}),\\
\mathrm{Gd}(\baselatoneS)(w,Z,c)&:\Longleftrightarrow\mathrm{Bu}(\baselatoneS)(w,Z,c)\wedge
  \mathrm{Gd}_0(w.D(Z),c.O.\mathrm{nch},c.O.x_d),
\end{align*}
and $\widehat{\mathrm{Bu}}(\baselatoneS)=\mathrm{pos}(\mathrm{Bu}(\baselatoneS))$,
$\widehat{\mathrm{Gd}}(\baselatoneS)=\mathrm{pos}(\mathrm{Gd}(\baselatoneS))$ (\emph{built}, resp.\ \emph{good}
cores). Here the \emph{goodness condition} $\mathrm{Gd}_0(D,\mathrm{nch},x_d)$ means: (tw) there is a twist
input $I=(E,N_0,S)$ for $D$ with $E=\mathrm{tw}(x_d)$ (the twist datum of $x_d$) which is good at all
twisted places for $\mathrm{nch}$ (condition $\mathrm{GA}(D,\mathrm{nch},I)$: at every place $w$ with
$\mu_w\ne\mu_0$ the input satisfies the full rank and basis conditions at the index $i_{\mathrm{nch}}(w)$,
and every power of the character $\mathrm{nch}.\psi'_w$ prime to $\ell$ is a good twist at
$(\mathrm{nch}.Q(w),\mathrm{nch}.k'(w))$, a condition of the same kind as those of \cref{pf:baselat1:def:goodtw}
which is not unfolded here) and satisfies a class condition; and (sep) for places $v,v'$ with
$\mu_v\neq\mu_0$, $\mu_{v'}\ne\mu_0$ and $i_{\mathrm{nch}}(v)\neq i_{\mathrm{nch}}(v')$ the sets
$\mathrm{nch}.Q(v)$ and $\mathrm{nch}.Q(v')$ are disjoint. Finally, $c$ is \emph{twisted at all places} if
\[
\mathrm{TA}(D,\mathrm{nch},x_d)\ :\Longleftrightarrow\ \mathrm{TC}(D,\mathrm{nch})\wedge
\exists\,N_0,S:\ \mathrm{GA}\bigl(D,\mathrm{nch},(\mathrm{tw}(x_d),N_0,S)\bigr),
\]
where $\mathrm{TC}$ is a compatibility of $\mathrm{nch}$ with complex conjugation.
\formal{GlobalTheory.Mem.Top.Bl.Slots}\formal{GlobalTheory.Mem.Top.Bl.Built}\formal{GlobalTheory.Mem.Top.Bl.BuiltE}%
\formal{GlobalTheory.Mem.Top.Bl.GoodE}\formal{GlobalTheory.Mem.Top.Bl.GoodO}\formal{GlobalTheory.EX.DQ.TwAllAt}
\end{definition}

\begin{definition}[Levels and construction steps]\label{pf:baselat1:def:levels}
Let $\baselatoneS$ be a slot family.
\begin{enumerate}[label=(\alph*)]
\item Level-three data $l\in\mathrm{Lv}_3(\baselatoneS,K,w,Z)$ consist of level-two data (among them a set of
  places $l.W_t$, a base-level row side $l.p_0$, a group identification, coherence and row data
  $l.\mathrm{what}\in\mathcal{R}(K,w,Z,l.W_t,l.p_0)$ with an element of $R_x$), a norm choice $l.\mathrm{nch}$,
  a norm-choice datum $l.n\in\mathcal{N}$, theta data $l.x_d$, the condition that for $k\in\{0,1\}$ and every
  place $v$ of the relevant type the level characters satisfy $\chi_T(k,v,j)=x^{2(j+1)}$ ($0\le j<m$) for an
  element $x$ of order $g_k-1$ in $(\Z/g_k\Z)^\times$, and an element of $N_x$. Level-four data
  $l\in\mathrm{Lv}_4(\baselatoneS,K,w,Z)$ add column data at split primes $l.\mathrm{col}$, real ties, the
  equalities $Z.S.\mathrm{dsig}(\tau)=d^{\mathrm{true}}(B.R,B.\mathrm{hc}(\tau))$ for all $\tau$ and
  $Z.S.s.\Delta=\Delta_v(K.s,K.c_D)$, and an element of $S_x$.
\item $\mathrm{H}_3(\baselatoneS)(w)\Leftrightarrow\exists Z:\ \mathrm{Lv}_3(\baselatoneS,K,w,Z)\ne\emptyset$;
  $\mathrm{H}_4$ likewise with $\mathrm{Lv}_4$.
\item $\mathrm{Pin}(w)$ means that $w$ has a column pin, and $\mathrm{Pin}\wedge p'$ is the conjunction;
  $w$ is \emph{starting}, $\mathrm{St}(w)$, if $\mathrm{Pin}(w)$, $\mathrm{Supp}(w)$ and $\mathrm{Coh}(B,T)$ hold, where
  $\mathrm{Supp}(w)$ means that $T$ is untwisted or the data of $w$ are good (a condition not needed here), and
  $\mathrm{Coh}(B,T)$ means that $T.Q_v,T.k_v,T.\psi_v$ depend only on the valuation ring induced by $v$ on
  $\mathit{LL}_0$ and that $\mathit{LL}$ contains a primitive root of unity of order $\ell_v^{T.k_v}$ for every $v$.
\item $\mathrm{Amb}(\baselatoneS)=\mathrm{Step}^{\mathrm w}\bigl(\mathrm{Pin},\ w\mapsto\forall Z\ \forall
  l\in\mathrm{Lv}_4:\ \exists a,a_0,\ \mathrm{Am}(l.\mathrm{what}.r.Z,a,a_0)\bigr)$, where $\mathrm{Am}$ is the
  condition that $a,a_0$ are ambient forms tied to the reading of the row data.
\item $\mathrm{Div}(\baselatoneS)=\mathrm{Step}^{\mathrm w}(\mathrm{Pin},\ w\mapsto\forall Z\ \forall l\in
  \mathrm{Lv}_4:\ \exists Y\in\mathcal{D}(w.D(Z))$ with $\mathrm{Rec}_{\mathrm{div}}(l.n,Y)\ne\emptyset$,
  $\mathrm{Mj}(l.\mathrm{what}.r.Z,Y)$, $Y.\kappa_R(\sigma)=\kappa^\sharp(w.D(Z))$ for all $\sigma$,
  and $O_x(O,t)\neq\emptyset$ for every object $O$ with $O.X=Y$, $O.p_0=l.p_0$ and every $t$ equal to
  $l.\mathrm{what})$. Here $\mathcal D(D)$ is the type of division data of $D$, $\mathrm{Rec}_{\mathrm{div}}$ the type of
  division records and $\mathrm{Mj}$ a tie between $Y$ and the reading of the row data.
\item $\mathrm{Th}(\baselatoneS)=\mathrm{Step}^{\mathrm w}\bigl(\mathrm{Pin},\ w\mapsto\forall Z\ \forall
  l\in\mathrm{Lv}_4:\ \exists\,\theta\in\mathcal{T}h(w.D(Z)),\ \mathrm{Rec}_{\theta}(l.n,\theta)\ne\emptyset\bigr)$, where
  $\mathcal Th(D)$ is the type of theta data and $\mathrm{Rec}_{\theta}$ the type of theta records.
\end{enumerate}
\formal{GlobalTheory.Mem.Top.Bl.Lv3}\formal{GlobalTheory.Mem.Top.Bl.Lv4}\formal{GlobalTheory.Mem.Top.Bl.Start}%
\formal{GlobalTheory.Mem.Top.Bl.AmbEx}\formal{GlobalTheory.Mem.Top.Bl.DivEx}\formal{GlobalTheory.Mem.Top.Bl.ThEx}
\end{definition}

\begin{definition}[Member existence]\label{pf:baselat1:def:memberE}
For a base predicate $\kappa$ and slots $Q,Q'$:
\begin{enumerate}[label=(\alph*)]
\item $\mathrm{First}(\kappa,Q)=\mathrm{Step}\bigl(\kappa,\ T\text{ untwisted},\ \mathrm{Ex}_{\mathrm{core}}(Q)\bigr)$;
\item $\mathrm{Exp}(\kappa,Q)$ is $\sfname{EvLgG}$ of the statement: if $\kappa(B,f)$ and $T$ is untwisted, then every
  core $c$ over stages $C,W,V,S,N,G$ on $B$ with $Q(c)$ satisfies the export condition
  $\mathcal{X}(\mathrm{asm}(B,C,W,V,S,N,G),c.O.\mathrm{nch},c.O.x_d)$;
\item $\mathrm{Second}(\kappa,Q,Q')=\mathrm{Pass}\bigl(\kappa,Q,\ \mathrm{Ex}_{\mathrm{core}}(c\mapsto Q(c)\wedge
  \mathrm{TA}(X.D,c.O.\mathrm{nch},c.O.x_d)\wedge Q'(c))\bigr)$, where $X.D$ denotes the assembled datum of the stage
  tuple of $c$.
\end{enumerate}
For a slot family $\baselatoneS$,
\[
\sfname{MemberE}(\baselatoneS)=\mathrm{BC}(\kappa_0)\wedge
\mathrm{First}(\kappa_0,\widehat{\mathrm{Bu}}(\baselatoneS))\wedge
\mathrm{Exp}(\kappa_0,\widehat{\mathrm{Bu}}(\baselatoneS))\wedge
\mathrm{Second}(\kappa_0,\widehat{\mathrm{Bu}}(\baselatoneS),\widehat{\mathrm{Gd}}(\baselatoneS)).
\]
The export condition $\mathcal{X}$ is recalled in \cref{pf:baselat1:def:exportAt}.
\formal{GlobalTheory.Mem.Top.GroupRetwisting.FirstG}\formal{GlobalTheory.Mem.Top.GroupRetwisting.ExportF}%
\formal{GlobalTheory.Mem.Top.GroupRetwisting.SecondG}\formal{GlobalTheory.Mem.Top.Bl.MemberE}
\end{definition}

\subsubsection{Parameters and hypothesis packages}

Fix a proof $h_L$ of $\sfname{RecTExists}$ (\cref{def:rectexists}). A \emph{norm-choice extension} $N$
assigns to every configuration $(K,B,f,T,w,Z,W_t,p_0,t,\mathrm{nch},n,x_d)$ (with $t$ row data and $n$ a
norm-choice datum) and every \emph{norm-choice core} $c$ (a record $\mathcal{C}(t,n)$ of local
data at the norm choice, a norm realisation and a selection condition) a type $N(\dots,c)$. A
\emph{parameter} $\baselatoneP=(R,N,S,\Gamma)$ consists of a row position, a norm-choice extension, a spectral
extension and a geometric predicate; its slot family is
\[
\baselatoneP.s=\bigl(R,\ (\dots)\mapsto\textstyle\sum_{c\in\mathcal{C}}N(\dots,c),\
\mathrm{SX}^{\mathrm e}(S),\ \mathrm{Ob}(\Gamma)\bigr),
\]
where $\mathrm{SX}^{\mathrm e}(S)(\dots)=\sum_{a\in\mathrm{SX}(\dots)}S(\dots,a)$ for a fixed type
$\mathrm{SX}$ of spectral data, and $\mathrm{Ob}(\Gamma)(O,t)$ is the type of pairs consisting of a proof of
$\Gamma(D,\mathrm{core}(t),O.X)$ and a proof of the $\lambda$-condition $\mathrm{La}(t.r.Z,O.X)$. We use only
\[
\baselatoneP_N:=\mathfrak{p}(N,\mathbf{1},\mathrm{cl}(\top))
=(R^{\mathrm{std}},N,\mathbf{1},\mathrm{cl}(\top)),\qquad
\baselatoneS_N:=\baselatoneP_N.s,
\]
where $R^{\mathrm{std}}$ is the standard row position (a record of conditions on the row data $t$),
$\mathbf{1}$ is the constant one-element type, $\top$ is the constantly true geometric predicate, and
$\mathrm{cl}(\Gamma)$ is the closure of $\Gamma$ under the division-transport steps
$(D,J,X)\mapsto(\mathrm{div}(D,N,h),\dots)$; in particular $\mathrm{cl}(\top)$ holds for all arguments (its
base clause applies). The \emph{law extension} $N_{\mathrm{law}}$ is the product of the default extension with
a law for the datum $w.D(Z)$.

\begin{definition}[Further conditions on slot families and parameters]\label{pf:baselat1:def:cond}
\begin{enumerate}[label=(\alph*)]
\item The \emph{level-three movement condition} $\mathrm{Mv}_3(\baselatoneS)$: for all $K,B,f,T,w,Z$, all
  $l\in\mathrm{Lv}_3(\baselatoneS,K,w,Z)$, every $S$-stage $S'$ and all choices of submodules $N_u$ and reals
  $h_u$ ($u$ a word) there is $l'\in\mathrm{Lv}_3(\baselatoneS,K,w,\mathrm{up}(Z,S',N,h))$ whose package is the
  update of the package of $l$: $\mathrm{pk}(l')=\mathrm{pk}(l).\mathrm{up}(Z,S',N,h)$.
\item The \emph{position movement condition} $\mathrm{Mv}(\baselatoneS)$ is a condition on the movement of
  objects and witnesses along division transports $Z\mapsto\mathrm{div}(Z,N,h)$; its form is not needed, since
  it is only passed on to the result cited in \cref{pf:baselat1:ssec:main}.
\item The \emph{multiplicity condition} $\mathrm{Mult}(\baselatoneP)$: eventually (in the parameter sense
  $\mathrm{Ev}_{\mathrm{par}}$), every weighted point $q$ of $\baselatoneP.s$ over $K$ satisfies the multiplicity
  equality for its real datum.
\item $\sfname{ExBL}(\kappa,Q,Q',\mathrm{Lk})$: there is a slot $R$ such that (i) eventually in the scale,
  for all sufficiently large data, if $\kappa(B,f)$ then every core with $Q$ satisfies $R$, and (ii)
  eventually (in the parameter sense $\mathrm{Ev}_{\mathrm{par}}$), for every bundle point $\beta$ whose core satisfies
  $Q'\wedge R$ and every $c\in\R$ there is a bundle $b$ with $\mathrm{Lk}(K,\beta,c,b)$;
  $\mathrm{Lk}(\baselatoneP)$ is the linkage through weighted points of \cref{sec:families}.
\end{enumerate}
\formal{GlobalTheory.Mem.SRe.LevelThreeCondition}\formal{GlobalTheory.Mem.Top.Bl.PositionMv}%
\formal{GlobalTheory.Mem.Top.Bl.MultiplicityCondition}\formal{GlobalTheory.Fc.ExBL}
\end{definition}

\begin{definition}[Hypothesis packages]\label{pf:baselat1:def:packages}
The following four packages consist of results proved elsewhere in the development (they are obtained from
the unitary-group hypotheses of \cref{sec:families}); we refer to their members by the labels below.
\begin{enumerate}[label=(\alph*)]
\item $\mathcal H_A$: (A1) ray class fields; (A2) the principal ideal theorem; (A3) CM elliptic curves with good
  reduction outside a prime; (A4) Harish-Chandra modules; (A5) rapid decay of cusp forms; (A6) the
  Casselman--Wallach theorem; (A7) the Weil representation of unitary dual pairs; (A8) existence of discrete
  series; (A9) Harish-Chandra's finiteness theorem; (A10) discrete series have the infinitesimal character given
  by their Harish-Chandra parameter; (A11) uniqueness of discrete series of $\Uni(p,q)$; (A12)
  finite-dimensionality of rows; (A13) the convergent Siegel--Weil formula for unitary groups; (A14) existence
  of discrete series with prescribed parameter, in the normalisation of volumes; (A15) the formal degree formula;
  (A16) the Siegel mass ratio \cref{prop:families:mass}; (A17) existence and (A18) uniqueness up to unitary
  equivalence of discrete series of $\Uni(n)$ with given Harish-Chandra parameter; (A19) existence of discrete
  series with given parameter for all group data of signature $(p,q)$; (A20) global Artin reciprocity; (A21)
  Weber's ray class fields.
\item $\mathcal H_E$: (E1) existence of degree-one primes with prescribed Frobenius modulo an open normal
  subgroup, avoiding a finite set (Chebotarev); (E2) Poitou--Tate for totally complex fields; (E3)
  cohomological dimension two; (E4) tame local theory; (E5) local finiteness; (E6) local Tate duality; (E7) the
  local Euler characteristic formula; (E8) global finiteness; (E9) the global Euler characteristic formula; (E10)
  unramified orthogonality; (E11) archimedean signature matching; (E12) the Greenberg--Wiles formula; (E13) Artin
  characters of ray class groups; (E14) Chebotarev in degree one; (E15) toric denominators; (E16) rationality and
  (E17) integrality of twist series; (E18) Hecke integrality; (E19) the Hecke transform and (E20) the Hecke
  reading tie; (E21) automorphic forms on $\GL_n$; (E22) Galois representations attached to polarized regular
  algebraic cuspidal representations of $\GL_a$ over CM fields; (E23) de Rham representations are potentially
  semistable; (E24) local--global compatibility at the places above $\ell$ for the Galois representations
  attached to conjugate self-dual regular algebraic cuspidal representations; (E25) the local Bloch--Kato
  statements $\mathrm{BK}_{\mathrm{loc}}(D)$ for all global data $D$.
\item $\mathcal H_M$: (H1) the discrete series cover; (H2) the theta-condition theorems (among them Green's
  description of cuspidal characters); (H3) partial Hecke $L$-functions; (H4), (H5) multiplicity equalities
  at pinned and at starting witnesses.
\item $\mathcal H_L(h_L)$: (L1) the third-layer package of published inputs; (L2) the series package; (L3)
  archimedean comparison, (L4) parameter and (L5) point statements for $h_L$; (L6) the statement that at
  every layer point the row sides have good cuspidal parameters.
\end{enumerate}
\formal{GlobalTheory.Mem.Top.Bl.HypAAlt}\formal{GlobalTheory.Mem.Top.Bl.EllipticHypothesis}%
\formal{GlobalTheory.Mem.Top.Bl.MultiplicityCoverHypotheses}\formal{GlobalTheory.Mem.Top.Bl.LocalLanglandsHypothesis}
\end{definition}

%======================================================================================================
\subsection{Steps, passes and member existence}\label{pf:baselat1:ssec:calculus}

\begin{lemma}\label{pf:baselat1:stephascore}
Let $\alpha$ be an at-predicate, $p$ a witness predicate and $P$ a core predicate. If
$\mathrm{Step}(\kappa_0,\alpha,(K,B,f,T)\mapsto\exists w,\ p(K,B,f,T,w))$ and
$\mathrm{Step}^{\mathrm w}(p,\mathrm{Ex}_{Z,c}(P))$, then $\mathrm{Step}(\kappa_0,\alpha,\mathrm{Ex}_{\mathrm{core}}(\mathrm{pos}(P)))$.
\formallabel{GlobalTheory.Mem.Top.Bl.step_hasCore}
\end{lemma}

\begin{proof}
By \cref{pf:baselat1:stepcalc}(i), $\mathrm{Step}(\kappa_0,\alpha,\mathrm{true})$; composing with the
second hypothesis (\cref{pf:baselat1:stepcalc}(iii)) gives $\mathrm{Step}(\kappa_0,\alpha,\forall w,\
p(w)\Rightarrow\mathrm{Ex}_{Z,c}(P)(w))$, and \cref{pf:baselat1:stepcalc}(ii) combines this with the first
hypothesis. It remains to see (\cref{pf:baselat1:stepcalc}(i),(iii)) that $\exists w\,p(w)$ together with
$\forall w\,(p(w)\Rightarrow\mathrm{Ex}_{Z,c}(P)(w))$ implies $\mathrm{Ex}_{\mathrm{core}}(\mathrm{pos}(P))$: choose
$w$ with $p(w)$ and then $Z,c$ with $P(w,Z,c)$; the stages $w.C,Z.W,Z.V,Z.S,Z.N,Z.G$ and the core $c$
witness $\mathrm{Ex}_{\mathrm{core}}(\mathrm{pos}(P))$ by \cref{pf:baselat1:posintro}.
\end{proof}

\begin{lemma}\label{pf:baselat1:stepWbuilt}
Let $\baselatoneS$ be a slot family and $p$ a witness predicate with $p\Rightarrow\mathrm{Pin}$. If
\[
\mathrm{Step}^{\mathrm w}(\mathrm{Pin}\wedge(\mathrm{H}_3(\baselatoneS)),\mathrm{H}_4(\baselatoneS)),\quad
\mathrm{Step}^{\mathrm w}(\mathrm{Pin}\wedge(\mathrm{H}_4(\baselatoneS)),\mathrm{Ex}_{Z,c}(\mathrm{Bu}(\baselatoneS))),\quad
\mathrm{Step}^{\mathrm w}(p,\mathrm{H}_3(\baselatoneS)),
\]
then $\mathrm{Step}^{\mathrm w}(p,\mathrm{Ex}_{Z,c}(\mathrm{Bu}(\baselatoneS)))$.
\formallabel{GlobalTheory.Mem.Top.Bl.MemberEx2.stepW_builtE}
\end{lemma}

\begin{proof}
By \cref{pf:baselat1:stepcalc}(iv), $\mathrm{Step}^{\mathrm w}(p,\mathrm{Pin})$. Combining it with the third
hypothesis (\cref{pf:baselat1:stepcalc}(v)) gives $\mathrm{Step}^{\mathrm w}(p,\mathrm{Pin}\wedge(\mathrm{H}_3))$, hence
by the first hypothesis $\mathrm{Step}^{\mathrm w}(p,\mathrm{H}_4)$; combining again with $\mathrm{Step}^{\mathrm w}(p,\mathrm{Pin})$
gives $\mathrm{Step}^{\mathrm w}(p,\mathrm{Pin}\wedge(\mathrm{H}_4))$, and the second hypothesis concludes.
\end{proof}

\begin{proposition}[Composition of a pass with a step]\label{pf:baselat1:passstep}
Let $\kappa$ be a base predicate with $\mathrm{BC}(\kappa)$, $Q$ a slot and $r,r'$ at-predicates. If
$\mathrm{Pass}(\kappa,Q,r)$ and $\mathrm{Step}(\kappa,r,r')$, then $\mathrm{Pass}(\kappa,Q,r')$.
\formallabel{GlobalTheory.Mem.Top.GroupRetwisting.pass_step}
\end{proposition}

\begin{proof}
By \cref{pf:baselat1:evg-inter} it suffices to fix a scale context $K$ for which both
$\mathrm{Lg}_2(K,h_s)$ and $\mathrm{Lg}_2(K,h_r)$ hold, where $h_s(B,f,T)$ is the implication
$\kappa(B,f)\Rightarrow\mathrm{Fd}(K,B,T)\Rightarrow r(K,B,f,T)\Rightarrow r'(K,B,f,T)$ and $h_r(B,f,T)$
is the statement inside $\mathrm{Pass}(\kappa,Q,r)$, and to prove $\mathrm{Lg}_2$ of the statement inside
$\mathrm{Pass}(\kappa,Q,r')$.

By \cref{pf:baselat1:large}, $\mathrm{Lg}_2(K,h_s\wedge\mathrm{Fd})$, hence $\mathrm{Lg}_2(K,h_s')$
with $h_s'(B,f,T)$ the implication $\kappa(B,f)\Rightarrow r(K,B,f,T)\Rightarrow r'(K,B,f,T)$. Let
$(S_1,P_1,A_0,S_{2,0})$ be exceptional data for $h_s'$ as in \cref{pf:baselat1:def:large}. The statement
$h_b(B,f,T)$: ``$S_1(B.\mathrm{pre},f,B.K_s^{\C})\subseteq B.F^{\mathrm{fix}}$ and
$P_1(B.\mathrm{pre},f,B.K_s^{\C},B.F^{\mathrm{fix}})\subseteq T.\mathrm{pl}$'' satisfies
$\mathrm{Lg}_2(K,h_b)$ with the exceptional data $(S_1,P_1,\emptyset,S_{2,0})$, since these two
inclusions are among the hypotheses in \cref{pf:baselat1:def:large}. By \cref{pf:baselat1:large}(i) it
suffices to show that $h_r\wedge h_b$ implies the statement inside $\mathrm{Pass}(\kappa,Q,r')$ for all
$(B,f,T)$ over $K$.

So let $(B,f,T)$ be over $K$ with $h_r(B,f,T)$ and $h_b(B,f,T)$; assume $\kappa(B,f)$, $T$ untwisted, and let
$c$ be a core over stages on $B$ with $Q(c)$ and the export condition; let $A$ be an exclusion function and
$S_2$ an assignment of finite sets of algebraic numbers. Apply $h_r(B,f,T)$ to $c$, to the exclusion function
$l\mapsto A(l)\cup A_0(B.\mathrm{pre},f,B.K_s^{\C},B.F^{\mathrm{fix}})(l)$ and to the assignment
$(P,g,\Xi)\mapsto S_2(P,g,\Xi)\cup S_{2,0}(P,g,B.K_s^{\C},B.F^{\mathrm{fix}},\Xi)$ (finite sets of algebraic
numbers). We obtain $B'$ with $\mathrm{Enl}(B,B')$, the frame $f'=f$, and $T'$ with
$T'.\mathrm{pl}=T.\mathrm{pl}$, avoiding both exclusion functions, such that both assignments evaluated at
$(B'.\mathrm{pre},f',T'.\mathrm{set}(f'))$ lie in $B'.\mathit{LLC}$, and $r(K,B',f',T')$.

Since $B'$ enlarges $B$, $B'.\mathrm{pre}=B.\mathrm{pre}$, $B'.K_s^{\C}=B.K_s^{\C}$ and
$B'.F^{\mathrm{fix}}=B.F^{\mathrm{fix}}$; hence $K$ lies over $B'.\mathrm{pre}$, and the exceptional data
$S_1,P_1,A_0,S_{2,0}$ take the same values at the data of $(B',f')$ as at those of $(B,f)$. Consequently
$S_1(\dots)\subseteq B'.F^{\mathrm{fix}}$ and $P_1(\dots)\subseteq T.\mathrm{pl}=T'.\mathrm{pl}$ (by
$h_b(B,f,T)$), $S_{2,0}(\dots,T'.\mathrm{set}(f'))\subseteq B'.\mathit{LLC}$, and $T'$ avoids
$A_0(\dots)$. Thus $h_s'(B',f',T')$ holds. Moreover $\kappa(B',f')$ holds by the enlargement clause of
$\mathrm{BC}(\kappa)$. Hence $r'(K,B',f',T')$, and $(B',f',T')$ has all properties required in
$\mathrm{Pass}(\kappa,Q,r')$.
\end{proof}

\begin{proposition}\label{pf:baselat1:baseCondOf}
Let $\kappa$ be a base predicate satisfying the following \emph{exactness condition} $\mathrm{KC}(\kappa)$: (ex)
eventually in the scale (in the sense $\mathrm{Ev}^{\circ}$ of \cref{pf:baselat1:def:evg}), for every base
$B_0$ over $K$, every frame $f_0$ of $B_0.\mathrm{pre}$ and all assignments $S_1,S_2$ of finite sets of
algebraic numbers, there is a base $B$ with $B.\mathrm{pre}=B_0.\mathrm{pre}$ and $\kappa(B,f_0)$,
$S_1(B.K_s^{\C})\subseteq B.F^{\mathrm{fix}}$, $S_2(B.K_s^{\C},B.F^{\mathrm{fix}})\subseteq B.\mathit{LLC}$; and
(enl) $\kappa$ is stable under enlargement. Let $\pi(P,f)$ be a predicate on pre-bases and frames such that
$\mathrm{Ev}$ of ``there are $P,f$ with $K.\mathrm{over}(P)$ and $\pi(P,f)$'' holds. Then
$\mathrm{BC}((B,f)\mapsto\kappa(B,f)\wedge\pi(B.\mathrm{pre},f))$.
\formallabel{GlobalTheory.Mem.Top.GroupRetwisting.baseCond_of}
\end{proposition}

\begin{proof}
An $\mathrm{Ev}^{\circ}$-statement implies the corresponding $\mathrm{Ev}$-statement
(\lref{GlobalTheory.Mem.Top.evG_of_ev}); by \cref{pf:baselat1:evg-inter} we may fix $K$ for which (ex)
holds and $P,f_0$ with $K.\mathrm{over}(P)$ and $\pi(P,f_0)$. There is a base $B_0$ over $P$
(\lref{GlobalTheory.Mem.BaseExt.exists_base}: for every pre-base and finite sets of algebraic numbers
there is a base over it whose fields contain them; here the sets are empty). Given $S_1,S_2$, (ex) applied to
$B_0,f_0,S_1,S_2$ gives $B$ over $P$ with $\kappa(B,f_0)$ and the required inclusions; also
$\pi(B.\mathrm{pre},f_0)=\pi(P,f_0)$ holds. The enlargement clause follows from (enl), since $\pi$ only
depends on the pre-base, which is unchanged by enlargement, and on the frame.
\end{proof}

\begin{proposition}[Base condition]\label{pf:baselat1:baseCond}
Assume (A1), (A2), (A3) and
(A16). Then $\mathrm{BC}(\kappa_0)$.
\formallabel{GlobalTheory.Mem.Top.Bl.KPxCond.baseCond_KPx}
\end{proposition}

\begin{proof}
Apply \cref{pf:baselat1:baseCondOf} with $\kappa$ the first conjunct of $\kappa_0$ (exactness of $K_s$ for some
field datum) and $\pi$ the second (the pre-base is pinned). The exactness condition for $\kappa$ is
\lref{GlobalTheory.Mem.KsExactCond.ksCond_exact}: from (A1)--(A3) one constructs a field
datum $\delta$ for the pre-base (from ray class fields and CM elliptic curves), a base over it on which
$K_s$ is exact for $(\delta,f_0)$ (from ray class fields and the principal ideal theorem), then a base with the
same $K_s$ whose fixed field contains the prescribed numbers, and an enlargement of it containing the second
set of prescribed numbers; exactness of $K_s$ only depends on the pre-base and on the image of $K_s$, so it is
stable under enlargement. For the second hypothesis of \cref{pf:baselat1:baseCondOf}: by
\lref{GlobalTheory.Mem.Top.ev_preX_frame}, which uses (A16), eventually every
pre-base over $K$ has a frame; and for every $K$ there is a pinned pre-base $P$ over $K$,
namely the coordinate pre-base of the pre-character datum of $K$
(\lref{GlobalTheory.Mem.Top.Bl.KPxCond.exists_prePin}). Choosing a frame of this $P$ gives the required
$\mathrm{Ev}$-statement.
\end{proof}

\begin{proposition}[Starting witnesses]\label{pf:baselat1:stepstart}
Assume (A16). Then
\[\mathrm{Step}\bigl(\kappa_0,\ T\text{ untwisted},\ (K,B,f,T)\mapsto\exists w\in\mathcal{W}(B,f,T):\
\mathrm{St}(K,B,f,T,w)\bigr).\]
\formallabel{GlobalTheory.Mem.Top.Bl.MemberEx2.step_start}
\end{proposition}

\begin{proof}
By \lref{GlobalTheory.Mem.Top.Bl.PinW.step_pinned},
$\mathrm{Step}(\kappa_0,\mathrm{true},\exists w\ \mathrm{Pin}(w))$: eventually, for all sufficiently
large data satisfying $\kappa_0$, there is a pinned column witness (it is assembled from a pre-datum and
frame given by \lref{GlobalTheory.Mem.Top.ev_preX_frame}, a canonical choice of twist indices, column
input data, and a selection of split places of half size). Compose it with $\mathrm{Step}(\kappa_0,
T\text{ untwisted},\mathrm{true})$ and combine with the trivial step $T$ untwisted $\Rightarrow$ $T$ untwisted
(\cref{pf:baselat1:stepcalc}). If $T$ is untwisted and $w$ is pinned, then $\mathrm{Supp}(w)$ holds (its first
alternative) and $\mathrm{Coh}(B,T)$ holds: all $T.Q_v$, $T.k_v$, $T.\psi_v$ are constant, and
$\ell_v^{T.k_v}=1$, so $1\in\mathit{LL}$ is the required root of unity. Hence $\mathrm{St}(w)$.
\end{proof}

\begin{theorem}[Member existence from construction steps]\label{pf:baselat1:steps2}
Assume (A1), (A2), (A3) and
(A16). Let $\baselatoneS,\baselatoneS'$ be slot families such that:
\begin{enumerate}[label=(\arabic*)]
\item $\mathrm{Step}^{\mathrm w}(\mathrm{St},\mathrm{H}_3(\baselatoneS))$;
\item $\mathrm{Step}^{\mathrm w}(\mathrm{Pin}\wedge(\mathrm{H}_3(\baselatoneS)),\mathrm{H}_4(\baselatoneS))$,
  $\mathrm{Amb}(\baselatoneS)$, $\mathrm{Div}(\baselatoneS)$ and $\mathrm{Th}(\baselatoneS)$;
\item $\mathrm{Exp}(\kappa_0,\widehat{\mathrm{Bu}}(\baselatoneS))$;
\item $\mathrm{Pass}\bigl(\kappa_0,\widehat{\mathrm{Bu}}(\baselatoneS),(K,B,f,T)\mapsto\exists w,\
  \mathrm{St}(w)\wedge\mathrm{H}_3(\baselatoneS')(w)\bigr)$;
\item the conditions (2) for $\baselatoneS'$;
\item every core built for $\baselatoneS'$ is good for $\baselatoneS$ and twisted at all places:
  $\mathrm{Bu}(\baselatoneS')(w,Z,c)\Rightarrow\mathrm{Gd}(\baselatoneS)(w,Z,c)\wedge
  \mathrm{TA}(w.D(Z),c.O.\mathrm{nch},c.O.x_d)$.
\end{enumerate}
Then $\sfname{MemberE}(\baselatoneS)$.
\formallabel{GlobalTheory.Mem.Top.Bl.memberE_of_steps2}
\formallabel{GlobalTheory.T.exists_member_of_construction_steps_secondPass}
\end{theorem}

\begin{proof}
We prove the four components of \cref{pf:baselat1:def:memberE}.

\emph{(M1)} is \cref{pf:baselat1:baseCond}.

\emph{(M2)}. By \cref{pf:baselat1:core} (with the mass ratio) and (2),
\[\mathrm{Step}^{\mathrm w}\bigl(\mathrm{Pin}\wedge(\mathrm{H}_4(\baselatoneS)),\mathrm{Ex}_{Z,c}(\mathrm{Bu}(\baselatoneS))\bigr).\]
Since $\mathrm{St}\Rightarrow\mathrm{Pin}$, \cref{pf:baselat1:stepWbuilt} with $p=\mathrm{St}$ and (1),
(2) gives $\mathrm{Step}^{\mathrm w}(\mathrm{St},\mathrm{Ex}_{Z,c}(\mathrm{Bu}(\baselatoneS)))$.
\Cref{pf:baselat1:stephascore}, applied with $\alpha=$ ``$T$ untwisted'' and \cref{pf:baselat1:stepstart},
gives
\[\mathrm{Step}\bigl(\kappa_0,T\text{ untwisted},\mathrm{Ex}_{\mathrm{core}}(\widehat{\mathrm{Bu}}(\baselatoneS))\bigr),\]
which is $\mathrm{First}(\kappa_0,\widehat{\mathrm{Bu}}(\baselatoneS))$.

\emph{(M3)} is hypothesis (3).

\emph{(M4)}. Let $p'(w)=\mathrm{St}(w)\wedge\mathrm{H}_3(\baselatoneS')(w)$, so that $p'\Rightarrow
\mathrm{Pin}$ and $\mathrm{Step}^{\mathrm w}(p',\mathrm{H}_3(\baselatoneS'))$ (\cref{pf:baselat1:stepcalc}(iv)). As in
(M2), (5) and \cref{pf:baselat1:core,pf:baselat1:stepWbuilt} give
\[\mathrm{Step}^{\mathrm w}\bigl(p',\mathrm{Ex}_{Z,c}(\mathrm{Bu}(\baselatoneS'))\bigr),\]
and by (6) and \cref{pf:baselat1:stepcalc}(iv),(v)
\[
\mathrm{Step}^{\mathrm w}\bigl(p',\ \mathrm{Ex}_{Z,c}(G)\bigr),\qquad G(w,Z,c)=\mathrm{Gd}(\baselatoneS)(w,Z,c)\wedge
\mathrm{TA}(w.D(Z),c.O.\mathrm{nch},c.O.x_d).
\]
Let $r(K,B,f,T)=\exists w\,p'(w)$. \Cref{pf:baselat1:stephascore} with $\alpha=r$ (and the trivial step
$r\Rightarrow r$) gives $\mathrm{Step}(\kappa_0,r,\mathrm{Ex}_{\mathrm{core}}(\mathrm{pos}(G)))$. Next,
$\mathrm{pos}(G)$ implies the slot $c\mapsto\widehat{\mathrm{Bu}}(\baselatoneS)(c)\wedge\mathrm{TA}(c)
\wedge\widehat{\mathrm{Gd}}(\baselatoneS)(c)$: by \cref{pf:baselat1:posintro} it suffices that $G(w,Z,c)$ implies
$\mathrm{pos}(\mathrm{Bu}(\baselatoneS))(w.\mathrm{st}(Z),c)$, $\mathrm{TA}$ and
$\mathrm{pos}(\mathrm{Gd}(\baselatoneS))(w.\mathrm{st}(Z),c)$, which holds by
\cref{pf:baselat1:posintro} because $\mathrm{Gd}\Rightarrow\mathrm{Bu}$. With
\cref{pf:baselat1:stepcalc}(i),(iii) we obtain $\mathrm{Step}(\kappa_0,r,r'')$, where $r''=\mathrm{Ex}_{\mathrm{core}}(
\widehat{\mathrm{Bu}}(\baselatoneS)\wedge\mathrm{TA}\wedge\widehat{\mathrm{Gd}}(\baselatoneS))$. Finally
\cref{pf:baselat1:passstep}, applied to (M1), to the pass (4) and to this step, gives
$\mathrm{Pass}(\kappa_0,\widehat{\mathrm{Bu}}(\baselatoneS),r'')$, which is
$\mathrm{Second}(\kappa_0,\widehat{\mathrm{Bu}}(\baselatoneS),\widehat{\mathrm{Gd}}(\baselatoneS))$.
\end{proof}

%======================================================================================================
\subsection{The construction steps}\label{pf:baselat1:ssec:stages}

This subsection proves the construction steps used in \cref{pf:baselat1:steps2} for the slot families
$\baselatoneS_N$ of the main theorem, and the closing theorem \cref{pf:baselat1:closing}.

\subsubsection{Gauge families}

Let $F$ be a number field, $E/F$ a quadratic extension with nontrivial automorphism $\sigma$, $\A_F$ the
ad\`eles of $F$, and for $M\in\N$ and an $M\times M$ matrix $J$ over $E$ let $U_J=\Uni_J(\A_F)$ be the group
of $\A_F\otimes_FE$-points of the unitary group of the hermitian form $(x,y)\mapsto{}^t\!\sigma(x)Jy$, and
$U_J(F)$ its subgroup of rational points.

\begin{definition}[Gauge families]\label{pf:baselat1:def:gauge}
A \emph{gauge family} $G$ assigns a measure $G.\mu_{M,J}$ on $U_J$ to every $M$ and
every matrix $J\in M_M(E)$. It is a \emph{gauge}, $\mathrm{Gauge}(G)$, if
\begin{enumerate}[label=(\alph*)]
\item (Haar) $G.\mu_{M,J}$ is a Haar measure (with the finiteness conditions required in the development)
  whenever $J$ is hermitian and nondegenerate;
\item (basis invariance) if $J$ is hermitian nondegenerate and $J'={}^t\!\sigma(P)JP$ with $P\in\GL_M(E)$,
  then for every isomorphism $e\colon U_J\to U_{J'}$ compatible with $P$ (i.e.\ $g\,P=P\,e(g)$ as
  matrices over $\A_F\otimes E$) one has $e_*G.\mu_{M,J}=G.\mu_{M,J'}$, and such an $e$ exists;
\item (step compatibility) let $M+1\ge3$, $J\in M_{M+1}(E)$ hermitian nondegenerate, $x\in E^{M+1}$ with
  $h_J(x,x)=\beta\in F^\times$, $B\in M_{(M+1)\times M}(E)$ with $h_J(x,B_{\cdot j})=0$ for all $j$, and
  $J'={}^t\!\sigma(B)JB$ nondegenerate. Then there is a nontrivial continuous unitary character $\psi$ of $\A_F$ which is
  trivial on $F$, and for every such $\psi$ there is a measure $\nu_X$ on $(\A_F\otimes_FE)^{M+1}$ which is a
  $\psi$-measure for the quadric $h_J(y,y)=\beta$, and for
  every such $\nu_X$ there is an isomorphism $e$ from the stabiliser $U_{J,x}$ of $x$ in $U_J$ onto $U_{J'}$
  compatible with $B$, and for every such $e$ the measures $G.\mu_{M+1,J}$, $\mu'=e^{-1}_*G.\mu_{M,J'}$ and
  $\nu_X$ are \emph{compatible}: for every measurable $\phi\colon U_J\to[0,\infty]$ and
  every section $s$ of the orbit map ($s(y)x=y$ for $y$ in the orbit of $x$),
  $\int\phi\,dG.\mu_{M+1,J}=\int\!\!\int\phi(s(y)h)\,d\mu'(h)\,d\nu_X(y)$.
\end{enumerate}
\formal{Formalism.MainNotions.GlobalTheoryTheory.GenusMassVolumes.GaugeFam}%
\formal{Formalism.MainNotions.GlobalTheoryTheory.GenusMassVolumes.GaugeFam.IsGauge}
\end{definition}

\begin{theorem}[Existence of gauges]\label{pf:baselat1:gauge}
Let $F$ be totally real and $E/F$ a CM extension which is unramified at all finite places (as given by the CM
data of a host). Then there is a gauge family $G$ with $\mathrm{Gauge}(G)$.
\formallabel{GlobalTheory.GaugeEx.exists_isGauge}
\end{theorem}

\begin{proof}
The CM data of the host give the hypotheses of \cref{pf:baselat1:gaugeCM}: $F$ is totally real, $E$ is totally
complex of degree $2$ over $F$, and $\sigma$ induces complex conjugation under every complex embedding. That theorem
gives the claim.
\end{proof}

\subsubsection{Constructed cores}

\begin{theorem}[Constructed cores]\label{pf:baselat1:core}
Assume (A16). Let $\baselatoneS$ be a slot family with $\mathrm{Amb}(\baselatoneS)$,
$\mathrm{Div}(\baselatoneS)$ and $\mathrm{Th}(\baselatoneS)$. Then
$\mathrm{Step}^{\mathrm w}(\mathrm{Pin}\wedge(\mathrm{H}_4(\baselatoneS)),\mathrm{Ex}_{Z,c}(\mathrm{Bu}(\baselatoneS)))$.
\formallabel{GlobalTheory.Mem.Top.Bl.CoreAsm.stepW_core}
\formallabel{GlobalTheory.T.exists_constructedCore_of_all_construction_stages}
\end{theorem}

\begin{proof}
Let $L$ be the statement of \cref{pf:baselat1:latlayer}, which holds by that proposition: it is a
step $\mathrm{Step}^{\mathrm w}(\mathrm{Pin},\cdot)$ asserting that for every level $Z$ of $w$, all column data
$t$ at split primes, every set of places $W_t$, norm choice $\mathrm{nch}$ with a norm-choice datum, and
every $Y\in\mathcal{D}(w.D(Z))$ such that the characters and places satisfy the compatibility required of cores
and $Y.Q_p$ is the set determined by $w$, there are a lattice instance $G_y$, a kernel family $\mathrm{ker}$
satisfying the lattice-layer condition, a kernel record $k_t\in\mathrm{Rec}_{\ker}(t,Y,G_y,\mathrm{ker})$, and
place and matrix data for which the kernel entries of $k_t$ are the canonical ones
(this is the condition $\mathrm{Ent}$). By \cref{pf:baselat1:stepcalc}(iv),(v) and the hypotheses we have
\[
\mathrm{Step}^{\mathrm w}\bigl(\mathrm{Pin}\wedge(\mathrm{H}_4),\ \mathrm{Pin}\wedge(\mathrm{H}_4)\wedge
((\text{ambient}\wedge\text{division})\wedge\text{theta})\wedge L\bigr),
\]
and by \cref{pf:baselat1:stepcalc}(iii),(iv) it remains to show, for fixed $K,B,f,T,w$ over $K$ with
$\mathrm{Fd}(K,B,T)$, that the right-hand side implies $\mathrm{Ex}_{Z,c}(\mathrm{Bu}(\baselatoneS))(w)$.
Take a column pin of $w$, a level $Z$ and $l\in\mathrm{Lv}_4(\baselatoneS,K,w,Z)$. The three construction
steps give ambient forms $a,a_0$ with $\mathrm{Am}(l.\mathrm{what}.r.Z,a,a_0)$, division data $Y$ with a
record $d_v\in\mathrm{Rec}_{\mathrm{div}}(l.n,Y)$, $\mathrm{Mj}$ and $Y.\kappa_R=\kappa^\sharp$, and theta data
$\theta$ with a record $t_t\in\mathrm{Rec}_{\theta}(l.n,\theta)$. The step $L$, applied to $Z$, $l.\mathrm{col}$,
$l.W_t$, $l.\mathrm{nch}$, $l.n$, $Y$ (its two side conditions are supplied by $l$ and $d_v$), gives
$G_y,\mathrm{ker},k_t$ and the place data. By \cref{pf:baselat1:gauge} applied to the CM extension of the host of
$w.D(Z)$ there is a gauge family $G$ with $\mathrm{Gauge}(G)$. Let $O$ be the object with components
the theta data $\theta$, the lattice data of $l.\mathrm{col}$, the twisted forms of $w$ at $Z$, $G_y$, the
error term of $l.\mathrm{col}$, $l.\mathrm{nch}$, $l.x_d$, the level data of $l.\mathrm{col}$, $Y$, $l.p_0$,
$\mathrm{ker}$, $a$, $a_0$ and $G$. By the last clause of $\mathrm{Div}$, $O_x(O,l.\mathrm{what})$ is
inhabited. Let $c$ be the core with object $O$, places $l.W_t$, characters $w.I.\xi$ and the following
components: forcing (hypothesis), the vanishing of the forced content at places outside $l.W_t$
(\lref{GlobalTheory.Mem.Top.Bl.CoreAsm.fval_off}, from the exactness of the places of $l$), the column
layer conditions (\lref{GlobalTheory.T.columnLayerConditions_trueColumns_laterStages}), the lattice
layer from $L$, the row layer (\lref{GlobalTheory.Mem.Top.Bl.CoreAsm.rowLay_of}, from $\mathrm{Mj}$
and $Y.\kappa_R=\kappa^\sharp$), the remaining fields copied from $l$ (stability, real ties,
$\mathrm{dsig}$ and $\Delta$ equalities, the $\eta$-value, the $\chi_T$ condition and the norm-choice ties),
the theta ties of $t_t$, and $\mathrm{Gauge}(G)$. Finally the record $b\in\mathcal{B}(\baselatoneS,K,w,Z,c)$
is formed from the pin, the
fields of $l$, the kernel family $\mathrm{ker}$ with $k_t$, $d_v$, $t_t$, the object record
$(\mathrm{Am},\mathrm{Mj},\kappa_R)$ and the element of $O_x$; the place data of $L$ prove
$\mathrm{Ent}(b.\mathrm{ker})$. Hence $\mathrm{Bu}(\baselatoneS)(w,Z,c)$.
\end{proof}

\begin{remark}\label{pf:baselat1:rem:core}
The hypothesis (A16) of \cref{pf:baselat1:core} is not used in its proof.
\end{remark}

\subsubsection{Ambient, division and theta data}

\begin{proposition}[Ambient data]\label{pf:baselat1:amb}
If $\baselatoneS.R_x=R^{\mathrm{std}}$, then $\mathrm{Amb}(\baselatoneS)$.
\formallabel{GlobalTheory.AmbCrd.ambEx_of_rx}
\end{proposition}

\begin{proof}
By \cref{pf:baselat1:stepcalc}(iv) it suffices to construct, for every level $Z$ and every
$l\in\mathrm{Lv}_4(\baselatoneS,K,w,Z)$, ambient forms $a,a_0$ with $\mathrm{Am}(l.\mathrm{what}.r.Z,a,a_0)$.
Since $R_x=R^{\mathrm{std}}$, the row data $l.\mathrm{what}$ satisfy the standard row conditions, which contain,
for all tuples $\mu$, tied hermitian forms and lattice ties, both for the row side $\sharp$ and for the
base-level row side $l.p_0$, and the identities relating the tuples of the reading of $l.\mathrm{what}$ to the
tuples of $w$. \Lref{GlobalTheory.AmbCrd.tie_level}, which builds an ambient form with identifications
and an ambient tie from such forms and lattice ties, is applied twice (to the row sides $\sharp$ and $l.p_0$),
giving $a$ and $a_0$.
\end{proof}

\begin{proposition}[Division data]\label{pf:baselat1:div}
Assume (A1), (A2) and (A3). For every row
position $R$, norm condition $N$ and spectral position $S$,
$\mathrm{Div}(R,N,S,\mathrm{Ob}(\mathrm{cl}(\top)))$.
\formallabel{GlobalTheory.Mem.Top.Closing.divEx_geoCl}
\end{proposition}

\begin{proof}
By \cref{pf:baselat1:divprov}, $\mathrm{Div}(\baselatoneS)$ holds for every slot family $\baselatoneS$ with
$O_x=\mathrm{Ob}(\top)$.

For the given $\baselatoneS=(R,N,S,\mathrm{Ob}(\mathrm{cl}(\top)))$ apply this to
$(R,N,S,\mathrm{Ob}(\top))$; the two slot families have the same $\mathrm{Lv}_4$, since
$\mathrm{Lv}_4$ does not involve $O_x$. Composing with the step of \cref{pf:baselat1:stepcalc}(iv) which
replaces an element $(g,\lambda)$ of $\mathrm{Ob}(\top)$ by
$(\mathrm{cl}\text{-base}(g),\lambda)$ gives the claim.
\end{proof}

\begin{remark}\label{pf:baselat1:rem:div}
The three hypotheses of \cref{pf:baselat1:div} are not used in its proof
(\cref{pf:baselat1:divprov}).
\end{remark}

\begin{lemma}\label{pf:baselat1:iozero}
If $\baselatoneS.R_x=R^{\mathrm{std}}$, then for all row data $t$ with $R_x(t)$, every tuple $\mu$ which is not
used at level $0$ ($\neg\mathrm{U}_0(w.D(Z),\mu)$) and every $j<m$, the coefficient
$\mathrm{io}_\mu(j)$ of the datum $w.D(Z)$ vanishes.
\formallabel{GlobalTheory.Mem.Top.Closing.io_zero}
\end{lemma}

\begin{proof}
This is part of the field $\mathrm{io}$ of the record $R^{\mathrm{std}}(t)$, applied to $\mu$ and $j$.
\end{proof}

\begin{proposition}[Theta data]\label{pf:baselat1:theta}
Assume (A4)--(A7) and (A17)--(A19). Let $\baselatoneS$ be a slot family such that the conclusion of
\cref{pf:baselat1:iozero} holds for $\baselatoneS$. Then $\mathrm{Th}(\baselatoneS)$.
\formallabel{GlobalTheory.ThetaRecordTrue.thEx_main}
\formallabel{GlobalTheory.T.exists_thetaData_trueRowsColumns_allCompatible}
\end{proposition}

\begin{proof}
By \cref{pf:baselat1:stepcalc}(iv) we fix $w$, a level $Z$ and $l\in\mathrm{Lv}_4$. From the rows
$l.\mathrm{what}$, the norm-choice datum $l.n$ and the vanishing of $\mathrm{io}_\mu$ for unused $\mu$, the
row function data are obtained (\lref{GlobalTheory.ThetaRecordTrue.recFun_of_rows}); they give the row-function
hypothesis $(\mathrm{RF})$ of \cref{pf:baselat1:def:alltied}.
\Cref{pf:baselat1:alltied}, applied to the datum $w.D(Z)$ with its character frame, the characters
$l.n.\mathrm{tp}.\xi$, the set of tuples used at level $0$ and the real ties of $l$,
gives theta data $X\in\mathcal{T}h(w.D(Z))$ with the $\mathrm{io}$-ties, an adelic Weil datum $A$ with
realisation $a$ which is tied to $X$, and, for the norm choice and tuple pins of $l$, a character comparison
$H$. Then $\mathrm{Rec}_{\theta}(l.n,X)$ is inhabited by the record formed from $A$, $a$, the Weil tie
(\lref{GlobalTheory.Mem.Th.weilTiedEq_of_eqTied}), $H$ (converted by
\lref{GlobalTheory.SameChrFin.sameChr_of_chrEq}), the standardness of the additive character of $a$, and
the $\mathrm{io}$-ties.
\end{proof}

\begin{remark}\label{pf:baselat1:rem:theta}
The hypotheses (A6) and (A4) of \cref{pf:baselat1:theta} are
not used in its proof.
\end{remark}

\subsubsection{Level four from level three}

\begin{definition}\label{pf:baselat1:def:satnew}
The following four statements are hypotheses of the passage from level three to level four.
\begin{enumerate}[label=(\alph*)]
\item \emph{Spectral saturation} $(\mathrm{SN})$: for every base $B$, column data $C$, late stages $a$ over
  $(B,C)$ with datum $a.D$, row package $p$ for $a.D$, forms $F_0$, and spectral package
  $P\in\mathrm{SPk}(a.D,F_0)$ with integrable spectral function, if $a.D$ satisfies the gap hypothesis
  $\mathrm{Gap}(a.D)$, the local ties $\mathrm{LT}(a.D)$ and the real ties $\mathrm{RTd}(a.D)$, if its constant
  $D^*$ has the standard value, if $(\mathrm{IR})(a.D)$ holds and $p$ satisfies the row hypothesis, then there is
  a new $S$-stage $Z_s$ such that for all division updates $(N,h)$ and forms $F_1$ on the replaced stages
  which agree with $F_0$, the spectral data at the replaced stages exist.
\item \emph{Real ties} $(\mathrm{RT})$: every datum $D$ with the volume tie $\mathrm{VT}(D)$, the true scale and
  the true signatures $D.\mathrm{dsig}(\tau)=d^{\mathrm{true}}(D.R,D.\mathrm{hc}(\tau))$ is real tied.
\item \emph{Spectral packages} $(\mathrm{SP})$: for every datum $D$ and forms $F$ with $\mathrm{Gap}(D)$ and
  $\mathrm{LT}(D)$ there is $P\in\mathrm{SPk}(D,F)$ whose spectral function is integrable for the archimedean
  measure.
\item \emph{Volumes and idempotents}: the volume--mass identity $(\mathrm{VM})$ for all $D$ with $\mathrm{VT}(D)$,
  and the reality $(\mathrm{IR})(D)$ of the good idempotent (\cref{pf:baselat1:idem}) for all $D$ with the series
  tie $\mathrm{ST}(D)$.
\end{enumerate}
Here a spectral package $\mathrm{SPk}(D,F)$ is a record consisting of a Hida setting $X$ for $D$ with an adelic
bridge isomorphism, a full datum $Y$ over $X$ with finiteness, a standard additive character $\psi$ of $\A_F$ with
compatible local and global measures, transports of the forms $F_u$ into $X$ for all words $u$ with level
matching and column dictionaries (into, bijective, isometric), the identification of the instance of $X$ with
the instance of $D$, and a good operator.
\formal{GlobalTheory.Mem.SRe.UpS.SatNew}\formal{GlobalTheory.Mem.SRe.UpS.RealOf}%
\formal{GlobalTheory.Mem.SRe.UpS.PackEx}\formal{GlobalTheory.Mem.SRe.SPack}
\end{definition}

\begin{theorem}[Level four from level three]\label{pf:baselat1:has4}
Assume (A5) and (A9)--(A16). Let $h_L$ be a
proof of $\sfname{RecTExists}$ and $\baselatoneP$ a parameter with $\baselatoneP.S=\mathbf{1}$ and
$\mathrm{Mv}_3(\baselatoneP.s)$. Then
$\mathrm{Step}^{\mathrm w}(\mathrm{Pin}\wedge(\mathrm{H}_3(\baselatoneP.s)),\mathrm{H}_4(\baselatoneP.s))$.
\formallabel{GlobalTheory.T.has4_of_has3_of_lv3Mv}
\formallabel{GlobalTheory.Mem.SRe.UpS.has4_of_has3_of_sExt_none}
\end{theorem}

\begin{proof}
The four statements of \cref{pf:baselat1:def:satnew} hold:
\begin{itemize}
\item $(\mathrm{SN})$ by \cref{pf:baselat1:satnewthm}, from (A5) and (A9)--(A12);
\item $(\mathrm{RT})$ by \cref{pf:baselat1:realOf} below, from (A14) and (A15);
\item $(\mathrm{SP})$ by \cref{pf:baselat1:spack} below, from (A12) and (A13);
\item $(\mathrm{VM})$ by \lref{GlobalTheory.Mem.VolMass.volumes_eq_masses_of_unfold}, which deduces the
  volume--mass identity from the volume ratio identity \lref{GlobalTheory.VolRatioFromSiegelWeil.volumeRatio}
  (a consequence of (A13)), and $(\mathrm{IR})$ by \cref{pf:baselat1:idem}.
\end{itemize}
Now (\lref{GlobalTheory.Mem.SRe.UpS.pin_has3_to_has4_of_sat_new}) eventually in the scale the gap
hypothesis holds for all data over $K$ (\lref{GlobalTheory.Mem.Top.NodeGap.ev_gapHyp}). Fix a pinned
witness $w$ with a level $Z$ and $l\in\mathrm{Lv}_3(\baselatoneP.s,K,w,Z)$. There is an $S$-stage $S$ over
$(B,w.C,Z.W,Z.G)$ with the true $\Delta$, signatures and scale, and the local ties at the updated stages
(\lref{GlobalTheory.Mem.SProv.exists_sStage_true}). The level-three condition, applied to $l$, $S$ and the
division update $(N_u,h)=((Z.V.\mathrm{div}.N_u).L,\,Z.V.\mathrm{div}.h)$, gives level-three data $l_0$ at the
updated level $\mathrm{up}(Z,S,N,h)$. \Lref{GlobalTheory.Mem.SRe.UpS.exists_late_level_of_gap}, applied
to $\baselatoneP$ (with $\baselatoneP.S=\mathbf{1}$), the four statements above, $l_0$, the gap
hypothesis and the properties of $S$, produces level-four data at a level of $w$; hence $\mathrm{H}_4(w)$.
\end{proof}

\begin{remark}\label{pf:baselat1:rem:has4}
The hypotheses (A16) and $h_L$ of \cref{pf:baselat1:has4} are not used in its proof.
\end{remark}

\begin{lemma}\label{pf:baselat1:realOf}
(A14) and (A15) imply $(\mathrm{RT})$.
\formallabel{GlobalTheory.Mem.SRe.UpS.realOf}
\end{lemma}

\begin{proof}
Let $D$ satisfy the hypotheses of $(\mathrm{RT})$. $\mathrm{RTd}(D)$ consists of four statements: the
scale tie, which holds because the scale recorded in $D$ is the true scale and the true scale is a scale
(\lref{GlobalTheory.Mem.Vol.ScaleExists.trueScale_isScale}); for every $\tau$, the existence of a
discrete series representation with the Harish-Chandra parameter attached to $D.\mathrm{hc}(\tau)$ whose formal
degree, for the Haar measure normalised by the true scale, is $D.\mathrm{dsig}(\tau)$; for every $\tau$, the
statement that every such representation has formal degree $D.\mathrm{dsig}(\tau)$; and $\mathrm{VT}(D)$,
which is a hypothesis. Since $D.\mathrm{dsig}(\tau)=d^{\mathrm{true}}(D.R,D.\mathrm{hc}(\tau))$, the second and
third statements are the two parts of \lref{GlobalTheory.T.exists_discreteSeries_formalDegree_product}
(existence of a discrete series with this parameter and formal degree $d^{\mathrm{true}}$, and uniqueness of the
formal degree), which follows from (A14) and (A15).
\end{proof}

\begin{lemma}\label{pf:baselat1:idem}
For every datum $D$ with $\mathrm{ST}(D)$ one has $(\mathrm{IR})(D)$: the coefficients of
the good idempotent $e_0$ of $D$ are fixed by complex conjugation.
\formallabel{GlobalTheory.Mem.SRe.UpS.goodIdempotent_real_of_harishChandra_series}
\end{lemma}

\begin{proof}
By \lref{GlobalTheory.GoodIdemReal.star_coeff_of_reverse} it suffices to prove, for $k\in\{0,1\}$, every
monic $f\in\F_{g_k}[X]$ with $f(0)\ne0$ and every $x$, that the complex conjugate of the $x$-coefficient of
the series idempotent $e_{\mathrm{ser},k}(f)$ is the $x$-coefficient of $e_{\mathrm{ser},k}(f(0)^{-1}f^{\mathrm{rev}})$,
where $f^{\mathrm{rev}}$ is the reversed polynomial. By $\mathrm{ST}(D)$, $e_{\mathrm{ser},k}$ is the series
idempotent attached to a datum $L$ for $\GL_r(\F_{g_k})$, and for these idempotents the identity is the map
identity \lref{GlobalTheory.GLq.GalSer.eSerC_map} applied to the duality relation between $f$ and
$f(0)^{-1}f^{\mathrm{rev}}$ (\lref{GlobalTheory.SeriesSupport.rel_dual}).
\end{proof}

\begin{proposition}[Spectral packages]\label{pf:baselat1:spack}
(A12) and (A13) imply $(\mathrm{SP})$.
\formallabel{GlobalTheory.Mem.SRe.exists_sPack_integrable_alt}
\end{proposition}

\begin{proof}
Let $D,F$ satisfy $\mathrm{Gap}(D)$ and $\mathrm{LT}(D)$. By \cref{pf:baselat1:specgap} there are a Hida setting $X$ with an
adelic bridge isomorphism, a standard additive character $\psi$, a finite full datum $Y$, transports of the forms,
and all compatibilities listed in \cref{pf:baselat1:def:satnew}, together with integrability of the spectral
function. There is a good operator for $X$ (\lref{GlobalTheory.rowDict_goodOp_nonempty}). These data form
an element of $\mathrm{SPk}(D,F)$.
\end{proof}

\subsubsection{The closing theorem}

\begin{lemma}\label{pf:baselat1:goodE}
Let $N,N_g$ be norm-choice extensions, $S$ a spectral extension and $\Gamma$ a geometric predicate, and assume that
for every configuration and norm-choice core $c$, an element of $N_g(\dots,c)$ yields an element of
$N(\dots,c)$ and $\mathrm{Gd}_0(w.D(Z),\mathrm{nch},x_d)$. Let $\baselatoneS=\mathfrak{p}(N,S,\Gamma).s$
and $\baselatoneS_g=\mathfrak{p}(N_g,S,\Gamma).s$. Then every core built for
$\baselatoneS_g$ is good for $\baselatoneS$ and satisfies $\mathrm{TA}(w.D(Z),c.O.\mathrm{nch},c.O.x_d)$.
\formallabel{GlobalTheory.Mem.Top.Closing.goodE_of_builtE_alt}
\end{lemma}

\begin{proof}
Let $b\in\mathcal{B}(\baselatoneS_g,K,w,Z,c)$ with $\mathrm{Ent}(b.\mathrm{ker})$. Its component
$b.n_x$ is a pair $(c_n,e)$ of a norm-choice core $c_n$ and $e\in N_g(\dots,c_n)$. By hypothesis $e$ yields
$e'\in N(\dots,c_n)$ and $\mathrm{Gd}_0(w.D(Z),c.O.\mathrm{nch},c.O.x_d)$. Replacing $b.n_x$ by $(c_n,e')$
gives a record in $\mathcal{B}(\baselatoneS,K,w,Z,c)$ (the two slot families differ only in $N_x$), with
the same kernel component, so $\mathrm{Gd}(\baselatoneS)(w,Z,c)$. For $\mathrm{TA}$: the
conjugation condition $\mathrm{TC}$ consists of the two conjugation clauses of the choice condition
$\mathrm{Ch}(w.D(Z),c.O.\mathrm{nch})$ contained in $c_n$; and the clause (tw) of $\mathrm{Gd}_0$ gives a
twist input $I=(E,N_0,S)$ with $E=\mathrm{tw}(c.O.x_d)$ and $\mathrm{GA}(w.D(Z),c.O.\mathrm{nch},I)$,
so that $N_0,S$ witness the existential in $\mathrm{TA}$.
\end{proof}

\begin{theorem}[Closing]\label{pf:baselat1:closing}
Assume (A1)--(A7) and (A9)--(A19) (in the formal statement, (A5) occurs twice). Let $h_L$ be a proof of
$\sfname{RecTExists}$ and $N,N_g$ norm-choice extensions, and assume
\begin{enumerate}[label=(\alph*)]
\item $\mathrm{Step}^{\mathrm w}(\mathrm{St},\mathrm{H}_3(\baselatoneS_N))$;
\item $\mathrm{Mv}_3(\baselatoneS_N)$ and $\mathrm{Mv}_3(\baselatoneS_{N_g})$;
\item $\mathrm{Exp}(\kappa_0,\widehat{\mathrm{Bu}}(\baselatoneS_N))$;
\item $\mathrm{Pass}\bigl(\kappa_0,\widehat{\mathrm{Bu}}(\baselatoneS_N),(K,B,f,T)\mapsto\exists w,\
  \mathrm{St}(w)\wedge\mathrm{H}_3(\baselatoneS_{N_g})(w)\bigr)$;
\item for every configuration and norm-choice core $c$, an element of $N_g(\dots,c)$ yields an element of
  $N(\dots,c)$ and $\mathrm{Gd}_0(w.D(Z),\mathrm{nch},x_d)$.
\end{enumerate}
Then $\sfname{MemberE}(\baselatoneS_N)$.
\formallabel{GlobalTheory.Mem.Top.Closing.memberE_closing_alt_alt}
\end{theorem}

\begin{proof}
Apply \cref{pf:baselat1:steps2} with $\baselatoneS=\baselatoneS_N$ and $\baselatoneS'=\baselatoneS_{N_g}$. Its
hypotheses: (1) is (a); (3) is (c); (4) is (d); (6) is \cref{pf:baselat1:goodE} with
$S=\mathbf{1}$ and $\Gamma=\mathrm{cl}(\top)$. For (2) and (5), let
$\baselatoneS''\in\{\baselatoneS_N,\baselatoneS_{N_g}\}$; it is the slot family of a parameter with spectral
extension $\mathbf{1}$, row position $R^{\mathrm{std}}$ and object component
$\mathrm{Ob}(\mathrm{cl}(\top))$. Level four from level three is \cref{pf:baselat1:has4}
(with (b)); $\mathrm{Amb}$ is \cref{pf:baselat1:amb}; $\mathrm{Div}$ is \cref{pf:baselat1:div};
$\mathrm{Th}$ is \cref{pf:baselat1:theta}, whose hypothesis on $\mathrm{io}$ holds by \cref{pf:baselat1:iozero}.
\end{proof}

%======================================================================================================
\subsection{Export: twisting to small Selmer groups}\label{pf:baselat1:ssec:export}

\subsubsection{The export condition}

Let $D$ be a global datum (\cref{sec:families}) with CM field $E$ (of its host), totally real subfield $F$,
algebraic closure $\Omega$ of $E$ with complex conjugation $c$ on $\Omega$, and $G_E=\Gal(\Omega/E)$. The
datum contains twisting data $X=D.\mathrm{ntw}$ and expansion data $Y=D.\mathrm{nexp}$. We recall the components
of $X$ which enter the arguments below. There is a finite index type $X.\mathrm{Tw}$ of \emph{twist indices};
for every $i$:
\begin{itemize}
\item a prime $\ell=\ell_i$, discrete valuation rings $\baselatoneO_0\subseteq\baselatoneO$ (depending on $i$)
  with $\ell$ in the maximal ideal of $\baselatoneO$, and for every $k\in\N$ a discrete valuation ring
  $\baselatoneO_k\supseteq\baselatoneO$ containing a primitive $\ell^k$-th root of unity $\zeta_k$;
\item a finite type $\Sigma(i)$ of \emph{signatures}; for $s\in\Sigma(i)$ a representation
  $\rho_s\colon G_E\to\GL_r(\baselatoneO_0)$, an index $j_s$ and a sign $b_s$, and characters
  $\chi_s=\chi_1(i,j_s,b_s)\colon G_E\to\baselatoneO_0^\times$ and $\delta_i\colon G_E\to\baselatoneO_0^\times$;
  we put $T_s=\chi_s^{-1}\rho_s$;
\item finite sets $S_0(i)$ and $p_D(i)$ of places of $E$; for every finite set $Q$ of places a ray class group
  $\mathrm{Cl}(i,Q)$ with an Artin map $\mathrm{art}_Q\colon G_E\to\mathrm{Cl}(i,Q)$ and restriction maps;
\item a finite type $\mathcal L_i$ of \emph{labels} $\lambda$, with signs $b_\lambda$ and integers
  $s_2(\lambda)$, and \emph{theta values} $\vartheta(Q,\lambda,\delta)\in\operatorname{Frac}\baselatoneO$ for
  $\delta\in\mathrm{Cl}(i,Q)$.
\end{itemize}
For a discrete valuation ring $R$ in which $\ell$ is not a unit and $x\ne0$ in $R$ we write
$\baselatonenv_R(x)=v_R(x)/v_R(\ell)$ for the \emph{normalised valuation}, extended to
$\operatorname{Frac}R$ by $\baselatonenv_R(a/b)=\baselatonenv_R(a)-\baselatonenv_R(b)$. For a homomorphism
$\psi$ of a group into an abelian group and a sign $b$ we write $\psi^{(b)}=\psi$ if $b$ is true and
$\psi^{(b)}=\psi^{-1}$ otherwise.

\begin{definition}[Twisted Selmer groups]\label{pf:baselat1:def:selmer}
Let $i\in X.\mathrm{Tw}$, $s\in\Sigma(i)$, $k\in\N$, $Q$ a finite set of places of $E$ and
$\psi\colon G_E\to\baselatoneO_k^\times$ a continuous character. Let $M_{s,\psi}=(\operatorname{Frac}
\baselatoneO_k/\baselatoneO_k)^r$ with $g\in G_E$ acting by $\psi(g)^{-1}\chi_s(g)^{-1}\rho_s(g)$. The
\emph{Selmer group} $\mathrm{Sel}(i,s,k,Q,\psi)\subseteq H^1(G_E,M_{s,\psi})$ is the submodule of classes whose
restriction to the decomposition group at every finite place $v$ lies in: the Bloch--Kato subgroup
$\mathrm{bk}_v$ (the image of the finite cocycles of the $\ell$-adic representation $T_s$ at $v$, defined with the
$p$-adic Hodge theoretic data of $D$) if $v\mid\ell$; all of $H^1$ if $v\nmid\ell$ and $v\in\mathrm{Spl}(i)\cup Q$,
where $\mathrm{Spl}(i)$ is the set of places above $g_0,g_1$ or in $p_D(i)$ not above $\ell$; the unramified
classes otherwise. Its \emph{normalised length} is
$L_{\mathrm{Sel}}(i,s,k,Q,\psi)=\len_{\baselatoneO_k}(\mathrm{Sel}(i,s,k,Q,\psi))/v_{\baselatoneO_k}(\ell)
\in[0,\infty]$.
\end{definition}

A signature $s$ is \emph{good} if there is a nonzero vector in the corresponding expansion space of $Y$ which is
fixed by the idempotent $e_{\mathrm{Alg}}$ and is an eigenvector of the Hecke operators at the split places
outside $S_0(i)$ with the eigenvalues $h_{\mathrm{Eig}}(i,s,\cdot,\cdot)$; $\Sigma^{\mathrm g}(i)$ is the finite set
of good signatures. Let $\nu(i)\in\N$ be the integer attached by the development to
$\ell$, the $\ell$-adic cyclotomic character of $E$ and the product $tr$ of two ranks of $D$ (its definition is
not needed here).

\begin{definition}[Good twists]\label{pf:baselat1:def:goodtw}
Let $i\in X.\mathrm{Tw}$, $A,Q$ finite sets of places of $E$, $k'\in\N$, $\psi'\colon\mathrm{Cl}(i,Q)\to
\baselatoneO_{k'}^\times$ a homomorphism, $\psi=\psi'\circ\mathrm{art}_Q$, and $\eps'\in\R$.
\begin{enumerate}[label=(\alph*)]
\item $\mathrm{T}_0(i,A,Q,k',\psi')$: every $P\in Q$ is split over $F$, of residue degree one and ramification
  index one over $\Q$; $Q\cap(S_0(i)\cup A)=\emptyset$; $k'\ge1$; $\ker\psi$ is open; $\psi(c g c^{-1})=\psi(g)^{-1}$;
  $\psi$ takes values in the subgroup generated by $\zeta_{k'}$; $\psi$ is ramified exactly at the places of $Q$;
  and $\psi$ is trivial on the decomposition groups at the places above $\ell$.
\item $\mathrm{T}_{\mathrm{Sel}}(i,Q,k',\psi',\eps')$: there are $\eps_s\ge0$ ($s\in\Sigma(i)$) with
  $\sum_{s\in\Sigma^{\mathrm g}(i)}\eps_s\le\eps'$ and $L_{\mathrm{Sel}}(i,s,k',Q,\psi^{(b_s)})\le\nu(i)+\eps_s$ for every
  good $s$.
\item $\mathrm{T}_{\Theta}(N_0,i,Q,k',\psi',\eps')$: for every label $\lambda$ the \emph{theta value}
  \[\Theta_Q(\lambda,\psi')=\sum_{\delta\in\mathrm{Cl}(i,Q)}\vartheta(Q,\lambda,\delta)\,\psi'^{(b_\lambda)}(\delta)
  \in\operatorname{Frac}\baselatoneO_{k'}\]
  is nonzero and
  $-N_0(i)\le\baselatonenv(\Theta_Q(\lambda,\psi'))\le c_0(\lambda)+\eps'$, where
  $c_0(\lambda)=\baselatonenv\bigl(\sum_\delta\vartheta(\emptyset,\lambda,\delta)\bigr)$.
\item $\mathrm{T}_{\mathrm{tm}}(i,Q,k',\psi',\eps')$: $t\cdot\sum_{P\in Q}\mathrm{tv}_P(\psi)\le\eps'$, where $t$ is a rank of
  $D$ and $\mathrm{tv}_P(\psi)$ is the infimum of $\baselatonenv(\psi(\sigma)-1)$ over the inertia elements
  $\sigma$ at $P$ with $\psi(\sigma)\ne1$.
\end{enumerate}
\end{definition}

The full \emph{good twist} condition $\mathrm{Good}(N_0,i,\eps',S,A,Q,k',\psi')$ is the conjunction of
$\mathrm{T}_0(i,\mathrm{levs}(S)\cup A,Q,k',\psi')$ (where $\mathrm{levs}(S)$ is the finite set of levels of
the basis elements $S$), index bounds for the twisted theta lattices with errors summing to at most $\eps'$,
two lattice conditions on the twisted theta lattices, and (b), (c), (d).

\begin{definition}[Export condition]\label{pf:baselat1:def:exportAt}
For $i\in X.\mathrm{Tw}$, $\mathrm{SH}(i)$ is the conjunction of: the true proposition; $\mathrm{SD}(i,s)$
(the dual of $T_s$ twisted by the cyclotomic character is conjugate to $\varepsilon^{t}T_s\circ c$) and
$\mathrm{NJ}(i,s)$ (all Hodge--Tate jumps $a$ of $T_s$ at places above $\ell$ satisfy $a<-1$ or $a>t$) for
all good $s$; $\sum_\delta\vartheta(\emptyset,\lambda,\delta)\ne0$ and $s_2(\lambda)=2s'$ for some
$1\le s'\le N(D.R)+1$, for all $\lambda$; $r\ge100$; and $2r\le g_k$ for $k=0,1$. For a twist datum $E$
(\cref{sec:families}) and $N_0\colon X.\mathrm{Tw}\to\N$, the statement
$\mathrm{GT}(E,N_0,i)$ means: if $\mathrm{SH}(i)$, then for every $\eps'>0$, if
the twisted theta series have full rank ($\mathrm{FR}$) for all $j\in[1,m]$ and signs $b$, then for all
bases $S$ ($\mathrm{Basis}$) and every finite set $A$ of places there are $Q$, $k'$ and
$\psi'\colon\mathrm{Cl}(i,Q)\to\baselatoneO_{k'}^\times$ such that $\mathrm{Good}(N_0,i,\eps',S,A,Q,k',\psi'^a)$
for every $a\in\Z$ prime to $\ell$. A twist index $i$ is a \emph{nonzero index place} if it is attached to a place
$w'$ of the coefficient field with $\mu_{w'}\ne\mu_0$ (the maximal ideal of $\baselatoneO_0$ restricts to the prime
of $\mathit{LL}_0$ at $w'$). Finally, for a norm choice $\mathrm{nch}$ (which assigns a twist index
$i_{\mathrm{nch}}(w)$ to every place $w$) and theta data $x_d$,
\begin{multline*}
\mathcal{X}(D,\mathrm{nch},x_d)\ :\Longleftrightarrow\ \exists\,N_0,S:\ \bigl(\forall i\text{ nonzero
index place}:\ \mathrm{GT}(\mathrm{tw}(x_d),N_0,i)\bigr)\\
\wedge\ \forall w\ (\mu_w\ne\mu_0):\ \mathrm{SH}(i_{\mathrm{nch}}(w))\wedge\mathrm{OK}(\dots),
\end{multline*}
where $\mathrm{OK}((\mathrm{tw}(x_d),N_0,S),i)$ is the conjunction of $\mathrm{FR}$ and
$\mathrm{Basis}(S_{i,j,b})$ for all $j\in[1,m]$ and $b$.
\end{definition}

\subsubsection{Reduction to the twist theorem}

\begin{theorem}[Export]\label{pf:baselat1:export}
Assume (E1)--(E25), (A5) and (A20), let $h_L$ be a proof of $\sfname{RecTExists}$, and assume (L1), (L2) and
(L6) for $h_L$. Let $N$ be a norm-choice extension such that every element of $N(\dots,c)$ yields an element of
$N_{\mathrm{law}}(\dots,c)$, and let $S$, $\Gamma$ be arbitrary. Then
$\mathrm{Exp}(\kappa_0,\widehat{\mathrm{Bu}}(\mathfrak{p}(N,S,\Gamma).s))$.
\formallabel{GlobalTheory.T.automorphy_of_hypotheses_alt_alt}
\formallabel{GlobalTheory.Mem.Top.ExportTrue.exportF_prmR_alt_alt}
\formallabel{GlobalTheory.ExportFirst.exportF_isBuiltE_alt_alt}
\end{theorem}

\begin{proof}
Write $\baselatoneP=\mathfrak{p}(N,S,\Gamma)$. By
\lref{GlobalTheory.T.nonempty_twOf_of_extensionWithLaw} (from (L1), (L2), $h_L$ and (L6)), eventually in
the scale and for all sufficiently large data: if $\kappa_0(B,f)$ and $T$ is untwisted, then for every
configuration $(w,Z,W_t,p_0,t,\mathrm{nch},n,x_d)$ and norm-choice core $c$, an element of
$N_{\mathrm{law}}(\dots,c)$, a proof of the standard row conditions for $t$ and a column pin of $w$ yield an element
of the \emph{twisted real join} $\mathrm{TRJ}(h_L,w.D(Z),\mathrm{nch},x_d,\dots)$, a record whose components are
described below.
By the hypothesis on $N$, the same holds with $N(\dots,c)$ in place of $N_{\mathrm{law}}(\dots,c)$
(\cref{pf:baselat1:evg-inter}(ii)). We show that this statement implies $\mathrm{Exp}$ pointwise. Let
$(B,f,T)$ satisfy it, $\kappa_0(B,f)$, $T$ untwisted, and let $c$ be a core over stages on $B$ with
$\widehat{\mathrm{Bu}}(\baselatoneP.s)(c)$: there are $w,Z,c'$ with $\mathrm{Bu}(\baselatoneP.s)(w,Z,c')$
such that the stages of $c$ and $c$ itself equal $w.C$, $Z$ and $c'$. Transporting along this equality, it
suffices to prove $\mathcal{X}(w.D(Z),c'.O.\mathrm{nch},c'.O.x_d)$. Let $b$ be the record of
$\mathcal{B}$. Its component $b.n_x$ consists of a norm-choice core $c_n$ and an element of $N(\dots,c_n)$;
together with the row conditions $b.r_x$ and the pin $b.\mathrm{pin}$, the statement above gives an
element $j$ of the twisted real join. Its components are: a real side $X_r$ of $D=w.D(Z)$; a series side $Y_s$ over
$X_r$ for the twist datum $\mathrm{tw}(x_d)$; Galois-side data $G=j.\mathrm{gal}\in\mathcal{G}(D)$; the
field condition $\mathrm{FO}(X_r)$; truth of the character groups $j.\mathrm{cg}$; and, since $T$ is
untwisted, the base-tuple condition $\mathrm{BT}(Y_s,i_{\mathrm{nch}}(w))$ at every place $w$ with
$\mu_w\ne\mu_0$. By \lref{UpperTheory.GalS.ghOf'_alt}, which derives the Hodge condition from the truth of the
character groups, $h_L$, (E21)--(E24) and the label truth of $G$, we have the Hodge condition
$\mathrm{Hdg}(D,G.\mathrm{wd})$. The core $c'$ contains the norm-choice tie $\mathrm{NT}(D,\mathrm{nch})$.
\Cref{pf:baselat1:seriesreal} now gives the claim.
\end{proof}

\begin{proposition}[Series realisation]\label{pf:baselat1:seriesreal}
Assume (E1)--(E20), (A5) and (A20). Let $D$ be a global datum with
norm choice $\mathrm{nch}$ and theta data $x_d$, $X_r$ a real side of $D$, $Y_s$ a series side over $X_r$ for
$\mathrm{tw}(x_d)$, and $G\in\mathcal{G}(D)$; assume $\mathrm{FO}(X_r)$, $\mathrm{Hdg}(D,G.\mathrm{wd})$,
$\mathrm{BK}_{\mathrm{loc}}(D)$, $\mathrm{NT}(D,\mathrm{nch})$, and $\mathrm{BT}(Y_s,i_{\mathrm{nch}}(w))$ for all $w$
with $\mu_w\ne\mu_0$. Then $\mathcal{X}(D,\mathrm{nch},x_d)$.
\formallabel{GlobalTheory.T.series_realization_of_published_alt}
\formallabel{GlobalTheory.ExportClose0.exportAt_bundles_alt}
\formallabel{GlobalTheory.T.realization_of_published_alt}
\formallabel{GlobalTheory.ExportClose.exportAt_of_bundles_alt}
\end{proposition}

\begin{proof}
From (A5), (A20), (E16), (E17) and
$\mathrm{FO}(X_r)$ one obtains the rationality of the rows of the series side at every nonzero index place
($\mathrm{RR}$) and the coefficient bound $\mathrm{XB}$
(\lref{GlobalTheory.T.thetaSeries_rationality_and_coefficientBound}); from (A5),
(E18), (E19) and (E20), the Hecke operators preserve the
sup-lattices of the real side ($\mathrm{HS}$, \lref{GlobalTheory.T.heckeOperator_preserves_supLattice_realSide});
and $\mathrm{FO}(X_r)$ contains that $\mathit{LL}_0$ is Galois over $\Q$ and contains all roots of unity of
order dividing $g_k(g_k-1)$, $k=0,1$.

\emph{Inputs.} By \cref{pf:baselat1:inputs} (from (E14),
(A20), (E15), $\mathrm{RR}$, $\mathrm{XB}$, $\mathrm{HS}$,
$\mathrm{Hdg}$ and the roots of unity) there are $N_0$ and a record $I\in\mathcal{I}(D,\mathrm{tw}(x_d),N_0)$;
its fields are the admissibility statements $\mathrm{adm}_1$, $\mathrm{adm}_2$, $\mathrm{adm}_8$, $\mathrm{adm}_9(N_0)$,
$\mathrm{adm}_{10}$ of the Bloch--Kato reduction, goodness of all signatures, the signature condition, $\mathrm{SH}$
at every nonzero index place, and the series, lattice and Hecke statements used in \cref{pf:baselat1:twistthm}.

\emph{Twisted places.} By \lref{GlobalTheory.Mem.Top.twistedPlaceIndex_mem_firstPair_of_compatibleChoice}
(from $\mathrm{NT}$ and the Galois property of $\mathit{LL}_0$), for every place $w$ with $\mu_w\ne\mu_0$ the
index $i_{\mathrm{nch}}(w)$ is a nonzero index place.

\emph{Bases.} For a nonzero index place $i$, $j\in[1,m]$ and a sign $b$, let $S_{i,j,b}$ be a basis of averaged
series as provided by \lref{GlobalTheory.T.exists_basis_averaged_series_of_class_sums} (from the field
$\mathrm{adm}_{T4}$ of $I$), and $S_{i,j,b}=\emptyset$ otherwise.

We verify $\mathcal{X}(D,\mathrm{nch},x_d)$ with these $N_0$ and $S$. For every nonzero index place $i$,
$\mathrm{GT}(\mathrm{tw}(x_d),N_0,i)$ is \cref{pf:baselat1:twistthm} (with the fields
of $I$ and $\mathrm{BK}_{\mathrm{loc}}(D)$). For $w$ with $\mu_w\ne\mu_0$ and $i=i_{\mathrm{nch}}(w)$:
$\mathrm{SH}(i)$ is a field of $I$; $\mathrm{FR}$ for $j\in[1,m]$ and every $b$ is
\lref{GlobalTheory.T.fullRank_thetaSeries_forms_of_rational_scalars} (from (A5), the base tuple
condition at $w$ and the rationality of rows at $i$); and $\mathrm{Basis}(S_{i,j,b})$ holds by the choice of
$S_{i,j,b}$, since $i$ is a nonzero index place.
\end{proof}

\begin{remark}\label{pf:baselat1:rem:export}
The formal name of \cref{pf:baselat1:export} refers to automorphy, but the statement is the export condition
$\mathrm{Exp}$; no automorphy statement occurs in it. Its conclusion holds for every spectral extension $S$ and
geometric predicate $\Gamma$.
\end{remark}

\subsubsection{The twist theorem}

The construction of good twists is carried out over the group rings of the groups
$\Delta_M^\iota=(\Z/\ell^M)^\iota$. For a finite type $\iota$ let
$\baselatoneLam_{M,\iota}=\baselatoneO_0[\Delta_M^\iota]$ and let
$j_\iota\colon\baselatoneO_0[X_q:q\in\iota]\to\baselatoneLam_{M,\iota}$ be the homomorphism with
$X_q\mapsto[e_q]-1$, $e_q$ the $q$-th basis vector. For characters $d_q\colon G_E\to\Z/\ell^M$ ($q\in\iota$) let
$\kappa_d\colon G_E\to\baselatoneLam_{M,\iota}^\times$, $g\mapsto[(d_q(g))_q]^{-1}$.

\begin{definition}[Auxiliary places and Selmer polynomials]\label{pf:baselat1:def:aux}
Let $i\in X.\mathrm{Tw}$, $\ell=\ell_i$, $M\ge1$.
\begin{enumerate}[label=(\alph*)]
\item For a split place $P$, a character $d\colon G_E\to\Z/\ell^M$ is \emph{good at $P$} if
  $\ker d$ is open, $d(cgc^{-1})=-d(g)$, $d$ is trivial on the decomposition groups above $\ell$ and on the
  inertia groups at all places other than $P$ and its conjugate $\bar P$, and $d$ maps the inertia group at $P$
  and at $\bar P$ onto $\Z/\ell^M$.
\item An \emph{admissible family} $(H,S,d)$ consists of an open normal subgroup $H\subseteq G_E$
  containing the inertia groups at the places of $p_D(i)$, a finite set $S$ of places, and characters $d_P$ good at
  $P$ for every split place $P\notin S$ whose decomposition group lies in $H$.
\item An element $z\in G_E$ is a \emph{tame inertia element at a twist prime} if it lies in an inertia group at
  a place $v\in p_D(i)$,
  $\delta_i(z)^2-1$ is a unit, and for every open normal subgroup $U$ some power $z^n$ with $n$ prime to
  $\ell\cdot\Nm v$ lies in $U$.
\item For a finite set $A$ of places and a family $P\colon\iota\to\{\text{places}\}$ with characters
  $d\colon\iota\to\mathrm{Hom}(G_E,\Z/\ell^M)$, $(P,d)$ is an \emph{auxiliary family of level $M$ avoiding $A$},
  $\mathrm{Aux}(i,M,A,P,d)$, if: every place $v$ equal or
  conjugate to some $P_q$ is split, of residue degree and ramification index one over $\Q$, not in $S_0(i)\cup A$,
  and satisfies $\Nm v\equiv1\pmod{2\ell^M}$; distinct $q$ give non-conjugate places; and $d_q$ is good at $P_q$.
  For $\iota=\{1,\dots,R\}\sqcup\{1,\dots,R\}$, $(P,d)$ is \emph{adapted to $(H,z)$}, $\mathrm{Adp}(i,M,H,z,P,d)$,
  if for every $a$ every Frobenius
  element at $P_{(1,a)}$ is conjugate to $z$ modulo $H$, and $d_{(1,b)}(\mathrm{Frob}_{P_{(2,a)}})=\delta_{ab}$.
  We write $Q(P)$ for the set of places equal or conjugate to some $P_q$.
\item For a signature $s$ and $Q=Q(P)$, a square matrix $\Phi$ over $\baselatoneLam_{M,\iota}$
  \emph{presents the Selmer modules} if for every discrete $\ell$-power torsion
  $\baselatoneLam_{M,\iota}$-module $\mathcal B$ there is an isomorphism, natural in $\mathcal B$, between the kernel of
  $\Phi$ on $\mathcal B^a$ and the Selmer group (\cref{pf:baselat1:def:selmer}, with $\mathcal B$ in place of
  $\operatorname{Frac}\baselatoneO_k/\baselatoneO_k$) of $\mathcal B^r$ with $g$ acting by $\kappa_d(g)T_s(g)$.
\item A polynomial $F\in\baselatoneO_0[X_q:q\in\iota]$ is a \emph{Selmer polynomial} for
  $(s,M,d,Q)$ if for all $k$, all $u\in(\baselatoneO_k^\times)^\iota$ with $u_q^{\ell^M}=1$ and
  $\psi=\prod_q u_q^{d_q}$: $L_{\mathrm{Sel}}(i,s,k,Q,\psi)=v_{\baselatoneO_k}(F((u_q-1)_q))/v_{\baselatoneO_k}(\ell)$.
\end{enumerate}
\formal{GlobalTheory.TwSys.GoodChar}\formal{GlobalTheory.TwSys.AdmFam}%
\formal{GlobalTheory.TwSys.IsTameInertiaElementAtTwistPrime}\formal{GlobalTheory.TwSys.SlopeTestRemainingChoices}%
\formal{GlobalTheory.TwSys.Steps}\formal{GlobalTheory.TwSys.IsQ}\formal{GlobalTheory.TwSys.Computes}%
\formal{GlobalTheory.TwSys.SelPoly}
\end{definition}

\begin{definition}[Galois-side statements]\label{pf:baselat1:def:galside}
The \emph{Galois hypotheses} $\mathrm{GH}(i,s)$ are the conjunction of: $\mathrm{adm}_1$ at $(i,s)$ ($\rho_s$ is
unramified outside $\ell$ and the auxiliary primes $g_0,g_1$; at the places above $\ell$, $\rho_s$ and $T_s$ are de
Rham and potentially crystalline, and crystalline if $\ell\notin\{g_0,g_1\}$; $T_s$ and its dual have no invariants
under the decomposition groups; and the Hodge dimensions at conjugate embeddings add up to $r$); $\mathrm{adm}_2$ at $(i,s)$ (the reduction of
$\rho_s$ spans the matrix algebra, and at the places above $g_k\ne\ell$ wild inertia acts trivially on it and some
inertia element has no eigenvalue equal to $\chi_s$); $\mathrm{SD}(i,s)$; $\mathrm{NJ}(i,s)$; the true
proposition; the local Bloch--Kato statements $\mathrm{adm}_3$; and $r\ge100$.
\begin{enumerate}[label=(\alph*)]
\item $(\mathrm{S}_1)$: for all $i$, $s$ with $\mathrm{GH}(i,s)$, $M\ge1$, families $(P,d)$ with
  $\mathrm{Aux}(i,M,\emptyset,P,d)$ and $Q=Q(P)$, there is a square matrix $\Phi$ over
  $\baselatoneLam_{M,\iota}$ which presents the Selmer modules and has no unit entry.
\item $(\mathrm{S}_2)$: in the situation of (a), if $\Phi$ presents the Selmer modules and
  $j_\iota(F)=\det\Phi$, then $F$ is a Selmer polynomial.
\item $\mathrm{CB}(i,\mu)$: for every set $\mathcal T$ of signatures with $\mathrm{GH}(i,s)$ for $s\in\mathcal T$
  there are $R,D_d,M_0$ such that for all $M\ge M_0$, admissible families $(H,S,d)$, tame inertia elements $z$ at a
  twist prime, open normal subgroups $H'$ and finite sets $A$ there is $P\colon\{1,\dots,R\}\sqcup\{1,\dots,R\}\to
  \{\text{places}\}$ with $\mathrm{Aux}(i,M,A,P,d\circ P)$, $\mathrm{Adp}(i,M,H',z,P,d\circ P)$, and such that for
  all $s\in\mathcal T$, $Q=Q(P)$, all $\Phi$ presenting the Selmer modules and all $F$ with
  $j_\iota(F)=\det\Phi$, some monomial of degree at most $D_d$ has coefficient of valuation at most $\mu$ in $F$.
\item $(\mathrm{S}_3)$: $\mathrm{CB}(i,0)$ whenever $\nu(i)=0$. $(\mathrm{S}_4)$:
  $\mathrm{CB}(i,v_{\baselatoneO_0}(\ell^{\nu(i)}))$ whenever $\nu(i)>0$.
\item $(\mathrm{S}_5)$: for every $i$ and every set $\mathcal T$ of signatures with
  $\mathrm{GH}$ there are $R,D_d$ and $M_0\ge1$ such that for all $M\ge M_0$, admissible families, tame inertia
  elements $z$, open normal $H'$ and finite $A$ there is $P$ with $\mathrm{Aux}(i,M,A,P,d\circ P)$ and
  $\mathrm{Adp}(i,M,H',z,P,d\circ P)$ such that for all $s\in\mathcal T$ and $Q=Q(P)$ there is a
  Selmer polynomial $F$ with a monomial of degree at most $D_d$ whose coefficient has valuation at most
  $v_{\baselatoneO_0}(\ell^{\nu(i)})$.
\end{enumerate}
\formal{GlobalTheory.TwSys.ThmS}\formal{GlobalTheory.TwSys.SelmerLengthsByDeterminant}%
\formal{GlobalTheory.TwSys.OrdBound}\formal{GlobalTheory.TwSys.ResidualBoundOddPrimes}%
\formal{GlobalTheory.TwSys.ThmC}\formal{GlobalTheory.TwSys.GaloisSideTwistedSelmerAuxiliaryPrimes}
\end{definition}

\begin{proposition}\label{pf:baselat1:auxprimes}
$(\mathrm{S}_1)$, $(\mathrm{S}_2)$, $(\mathrm{S}_3)$ and $(\mathrm{S}_4)$ imply
$(\mathrm{S}_5)$.
\formallabel{GlobalTheory.TwSys.galois_twisted_selmer_auxprimes_of_matrix_and_bounds}
\end{proposition}

\begin{proof}
Fix $i$ and $\mathcal T$. We have $\mathrm{CB}(i,v(\ell^{\nu(i)}))$: if $\nu(i)=0$ then
$v(\ell^0)=v(1)=0$ and this is $(\mathrm{S}_3)$; otherwise it is $(\mathrm{S}_4)$. Let $R,D_d,M_0$ be
as in $\mathrm{CB}$ for $\mathcal T$, and replace $M_0$ by $\max(M_0,1)$. Given $M\ge M_0$ and the remaining data,
let $P$ be as in $\mathrm{CB}$. For $s\in\mathcal T$ and $Q=Q(P)$, $(\mathrm{S}_1)$ (applicable since
$M\ge1$ and $\mathrm{Aux}$ for $A$ implies it for $\emptyset$) gives $\Phi$ presenting the Selmer modules. The map
$j_\iota$ is surjective ($\baselatoneLam_{M,\iota}$ is generated by the elements $[e_q]$, and
$[x]=\prod_q(1+([e_q]-1))^{x_q}$), so $\det\Phi=j_\iota(F)$ for some $F$. By
$(\mathrm{S}_2)$ $F$ is a Selmer polynomial, and by the choice of $P$ it has the required
coefficient.
\end{proof}

The three Galois-side inputs are proved in the formal development from the Galois-cohomological hypotheses
of \cref{pf:baselat1:export}; we state them and describe their proofs.

\begin{theorem}[Galois-side inputs]\label{pf:baselat1:galinputs}
\begin{enumerate}[label=(\roman*)]
\item The hypotheses (E2), (E3), (E5),
  (E6), (E7), (E8), (E9) and
  (E11) imply $(\mathrm{S}_1)$ for every datum
  \textup{(\lref{GlobalTheory.T.selmerPresentationMatrix_of_galois_cohomology})}.
\item $(\mathrm{S}_2)$ holds for every datum
  \textup{(\lref{GlobalTheory.SelDetLength.twisted_selmer_lengths_eq_determinant_valuation})}.
\item (E1)--(E12) imply $(\mathrm{S}_3)$
  \textup{(\lref{GlobalTheory.T.residualSelmerBound_of_auxiliaryPrimes_odd})}, and (E1)--(E11) imply
  $(\mathrm{S}_4)$ (\cref{pf:baselat1:boundtwo}).
\end{enumerate}
\end{theorem}

The proofs are as follows. (i) The Selmer functor $\mathcal B\mapsto\mathrm{Sel}(\mathcal B^r)$ on discrete
$\ell$-power torsion $\baselatoneLam_{M,\iota}$-modules is shown to be left exact
(\lref{GlobalTheory.T.selmer_functor_left_exact_torsion_coefficients}), to commute with colimits, and to have
no Galois invariants with torsion coefficients; a left exact functor of this kind whose value on the residue field is
bounded by the residual Selmer group (\lref{GlobalTheory.T.selmerObstruction_le_residualSelmer_simpleExtension},
proved with Poitou--Tate and the local and global Euler characteristic formulas) is the kernel of a square matrix
with non-unit entries (\lref{GlobalTheory.T.torsion_left_exact_functor_kernel_square_matrix}). (ii) For a square
matrix $\Phi$ over a discrete valuation ring, the length of the kernel of $\Phi$ on $(\operatorname{Frac}/\mathcal
O)^a$ equals the valuation of $\det\Phi$ (\lref{GlobalTheory.length_ker_matAct_kmod_eq_addVal_det}); this is
applied to the specialisation of $\Phi$ along $\Delta_M^\iota\to\baselatoneO_k^\times$, $e_q\mapsto u_q$. (iii) The
coefficient bounds are proved by choosing the auxiliary places by Chebotarev's theorem with prescribed Frobenius
elements (\lref{GlobalTheory.T.exists_auxiliary_primes_prescribed_frobenius}) and bounding the residual Selmer
groups: the residual Selmer group and its dual have the same dimension
(\lref{GlobalTheory.T.residual_selmer_finrank_eq_dual}), local cohomology at the auxiliary places vanishes in
the relevant degree, and global reciprocity for the cup product pairing
(\lref{GlobalTheory.ResRecip.reciprocity}) controls the order of the residual determinant. When $\nu(i)>0$
the bound $(\mathrm{S}_4)$ is proved in \cref{pf:baselat1:boundtwo} by an induction in which classes of the dual Selmer
group are killed by introducing new auxiliary primes (\cref{pf:baselat1:newprime}).

\begin{theorem}[Twist theorem]\label{pf:baselat1:twistthm}
Assume (E1)--(E13). Let $D$ be a global datum with twist datum $E$ and $N_0\colon X.\mathrm{Tw}\to\N$, and assume
$\mathrm{BK}_{\mathrm{loc}}(D)$ and the following statements about $D$ (they are the fields of the input record of
\cref{pf:baselat1:seriesreal}): the admissibility conditions $\mathrm{adm}_1$, $\mathrm{adm}_2$, $\mathrm{adm}_8$
(for every good signature there is a character $\eta$ such that $\eta\rho_s$ has prescribed restrictions to the
inertia groups and the characteristic polynomials of Frobenius of $\eta\rho_s$ are the Hecke polynomials),
$\mathrm{adm}_9(N_0)$ (integrality of $\ell^{N_0(i)}\vartheta$) and $\mathrm{adm}_{10}$ (the recursion of theta values
under removal of a pair of conjugate auxiliary primes); goodness of all signatures; the signature condition; the
universal character condition; the range condition on the Hecke indices; the corresponding recursion $\mathrm{UR}$ of
the universal theta elements; the specialisation of twisted series; the condition $\mathrm{adm}_{T4}(E)$; the
type-keeping condition for $E$; the sup-lattice conditions (universal, for the group, full, for the Hecke operators);
commutation of the group action and of the Hecke operators; and semisimplicity of the Hecke action. Then
$\mathrm{GT}(E,N_0,i)$ for every nonzero index place $i$.
\formallabel{GlobalTheory.TwExportTop.twistAllPowers_of_galoisCohomology_classField}
\formallabel{GlobalTheory.TwExportSup.twistTheorem_allPowers_of_selmer_supLattice}
\end{theorem}

\begin{proof}
By \cref{pf:baselat1:galinputs} we have $(\mathrm{S}_1)$, $(\mathrm{S}_2)$,
$(\mathrm{S}_3)$ and $(\mathrm{S}_4)$ for the Bloch--Kato data of $D$, hence
$(\mathrm{S}_5)$ by \cref{pf:baselat1:auxprimes}. The \emph{admissibility package}
$\mathcal{A}(N_0)=\mathrm{adm}_1\wedge\mathrm{adm}_2\wedge\mathrm{adm}_3\wedge\textup{(E11)}
\wedge\mathrm{adm}_9(N_0)\wedge\mathrm{adm}_{10}$ holds: $\mathrm{adm}_3$ is the first part of
\lref{GlobalTheory.BlochKato₂.duality_dim_of_localTheorem}, which derives the local Bloch--Kato core
statements from $\mathrm{BK}_{\mathrm{loc}}(D)$, and the other components are hypotheses. The \emph{coordinate tie}
$\mathrm{CT}$ (for every nonzero index place $i$ and $M\ge1$ there is an admissible family $(H,S,d)$ whose
characters $d_P$ are the coordinates of the universal character $\delta_U(i,Q,M)$ on conjugation-closed sets $Q$)
follows from the universal character condition and the existence of an admissible family of ray class characters
(\lref{GlobalTheory.UnivChar.delUTie_of_isUnivChar}, \lref{GlobalTheory.AdmChars.exists_admissible_family},
the latter using (E13)). From $\mathrm{UR}$, the Hecke laws (commutation of the Hecke operators and
$h_{\mathrm{Eig}}(i,s,0,P)=1$, which follow from the commutation and range conditions), goodness of all signatures
and $\mathrm{adm}_8$ one obtains the index expansion $(\mathrm{E}_{\mathrm{ind}})$ of \cref{pf:baselat1:def:ordhyp}
(\lref{GlobalTheory.TwistDet.propABCD_of_rel}). By \cref{pf:baselat1:choice} we obtain the conclusion
$\mathrm{AIB}(N_0)$ of that theorem. Further, the isotypic parts of the twisted forms form an internal direct
sum (\lref{GlobalTheory.T.admInt_of_sup_laws}, from the sup-lattice, type-keeping and commutation
conditions), the group algebra acts integrally (\lref{GlobalTheory.T.admALG_of_commG_var}), and good
lattice points exist (\lref{GlobalTheory.T.admT3_of_hecke_semisimple}, from these two statements, the
signature condition and semisimplicity). Finally \lref{GlobalTheory.exportAll_of_six_inputs}, i.e.\
\lref{GlobalTheory.ExportFromInputs.twistModule_all_character_powers_of_admissible}, turns
$\mathrm{AIB}(N_0)$, these three statements, $\mathrm{adm}_{T4}(E)$ and the specialisation of
twisted series into $\mathrm{GT}(E,N_0,i)$ for every nonzero index place $i$: the character
$\psi'$ of \cref{pf:baselat1:choice} is used unchanged, and the index bounds and lattice conditions of
$\mathrm{Good}$ are deduced from the bounds $\mathrm{T}_{\mathrm{lat}}$ for the lattices spanned by the bases $S$.
\end{proof}

\begin{remark}\label{pf:baselat1:rem:gw}
The hypothesis (E12) is passed to \lref{GlobalTheory.T.residualSelmerBound_of_auxiliaryPrimes_odd},
whose proof (\lref{GlobalTheory.ResSel.residual_bound_of_parts}) does not use it; it is not used elsewhere in the
proof of \cref{pf:baselat1:twistthm}. Likewise \lref{GlobalTheory.SelFun.oneMatrix_of_galoisCohomology} does not
use (E1), (E4) and (E10).
\end{remark}

\subsubsection{Choice of the twisting character}

For a nonzero index place $i$, a family $\mathcal F$ of lattices of twisted forms (lattices
$B_{j,b}$ and ``balls'' $\mathrm{Ball}_{j,b}$ in the form spaces of $Y$, operators $A_{j,b}$ and finite sets
$S_{j,b}$ of data) satisfying the lattice hypotheses $\mathrm{AB}$ (finite generation and Hecke stability of the
balls and lattices, compatibility of $A_{j,b}$ with the Hecke operators, and the condition that the images of the
theta elements of $S_{j,b}$ form a basis of the space spanned by $B_{j,b}$), the statement $\mathrm{T}_{\mathrm{lat}}(\mathcal
F,Q,k',\psi',\eps')$ means: there are $M$ and a character $\chi$ of $(\Z/\ell^M)^Q$ with
$\psi'\circ\mathrm{art}_Q=\chi\circ\delta_U(i,Q,M)$, and $\eps_{j,b}\ge0$ with $\sum_{j\in[1,m],b}\eps_{j,b}\le\eps'$,
such that for all $j\in[1,m]$, $b$ and every $\baselatoneO_{k'}$-lattice $T_\psi$ in
$\baselatoneO_{k'}\otimes B_{j,b}$ containing the images under $A_{j,b}$ of the specialisations at $\chi^{(b)}$ of
the universal theta elements of $S_{j,b}$, the normalised index satisfies
$\operatorname{ind}_{\mathrm n}(\baselatoneO_{k'}\otimes B_{j,b},T_\psi)\le\operatorname{ind}_{\mathrm n}(B_{j,b},\Theta_{j,b})+\eps_{j,b}$.

\begin{definition}[Polynomial expansions]\label{pf:baselat1:def:ordhyp}
Let $\baselatoneO$ be a discrete valuation ring, $\ell$ a prime, $M,R,n_{\max},v\in\N$. A polynomial
$f\in\baselatoneO[X_q:q\in\{1,\dots,R\}\sqcup\{1,\dots,R\}]$ satisfies $\mathrm{OH}(M,R,n_{\max},v)$ if there are
polynomials $D_d,g$ with
\begin{enumerate}[label=(\roman*)]
\item $f-D_dg$ lies in the ideal generated by the $X_{(2,a)}$ and the $(1+X_{(1,a)})^{\ell^M}-1$;
\item $v(g(0))\le v$;
\item modulo the maximal ideal, $D_d\equiv u\prod_aX_{(1,a)}^{n_a}$ modulo the ideal generated by the
  $X_q^{\ell^M}$, for a polynomial $u$ with $u(0)\ne0$ and exponents $n_a\le n_{\max}$.
\end{enumerate}
$(\mathrm{E}_{\Theta})(N_0)$ and $(\mathrm{E}_{\mathrm{ind}})$ are the statements that the theta values $\ell^{N_0(i)}\Theta_Q(\lambda,\psi')$,
respectively the index bounds for the lattices of $\mathrm{T}_{\mathrm{lat}}$, are values $f((u_q-1)_q)$ of polynomials $f$
satisfying $\mathrm{OH}$ with $n_{\max}=2$, $v=v_E(\lambda)$, $v_E(\lambda)/v(\ell)=c_0(\lambda)+N_0(i)$,
respectively with constants $n_B,v_B$ and an excluded finite set $A_B$, for all auxiliary families satisfying
$\mathrm{Aux}$ and $\mathrm{Adp}$ with respect to an open normal subgroup $H_E(\lambda)$, respectively $H_B$,
and all $\psi'$ with $\psi'\circ\mathrm{art}_Q=\prod_qu_q^{d_q}$ (for $(\mathrm{E}_{\mathrm{ind}})$: whenever the value is nonzero, the
index is at most the index of $\Theta$ plus $\baselatonenv(f(u-1))-v_B/v(\ell)$).
\end{definition}

\begin{lemma}\label{pf:baselat1:limitpow}
Let $\Gamma$ be a compact topological group, $\sigma\in\Gamma$ and $m\in\N$. There is $z\in\Gamma$ which lies in every
closed subgroup containing $\sigma$, such that for every open subgroup $U$ one has $z\equiv\sigma^{m^k}$ modulo $U$
for some $k$, and for every open normal subgroup $U$ one has $z^n\in U$ for some $n$ prime to $m$.
\formallabel{GlobalTheory.TwSys.InertiaElt.exists_limit_pow}
\end{lemma}

\begin{proof}
Let $z$ be a cluster point of the sequence $(\sigma^{m^k})_k$ (compactness). Closed subgroups containing $\sigma$ contain
all $\sigma^{m^k}$, hence $z$. For open $U$ the neighbourhood $zU$ of $z$ contains $\sigma^{m^k}$ for infinitely many
$k$. If $U$ is also normal, $\Gamma/U$ is finite; write the order of the image of $\sigma$ as $t'K$ with $t'$ prime to
$m$ and $K$ dividing a power $m^{k_0}$. Choose $k\ge k_0$ with $z\equiv\sigma^{m^k}$ modulo $U$; then the image of
$z^{t'}=\sigma^{m^kt'}$ in $\Gamma/U$ is trivial, since $t'K$ divides $m^kt'$. Take $n=t'$.
\end{proof}

\begin{lemma}\label{pf:baselat1:tame}
For every twist index $i$ there is a tame inertia element at a twist prime.
\formallabel{GlobalTheory.TwSys.exists_tameInertiaElementAtAuxiliaryPrime}
\end{lemma}

\begin{proof}
The set $p_D(i)$ is the set of places above one of two auxiliary primes of the host, chosen different from $\ell$;
let $v\in p_D(i)$, and choose a place of $\Omega$ above $v$ with inertia group $I_v$. Let $p\in\{3,5\}$ with
$p\ne\ell$. The image $\delta_i(I_v)$ has order $15$ (a property of $\delta_i$), so it contains an element
$\delta_i(\sigma)$ of order $p$, $\sigma\in I_v$. Let $z$ be as in \cref{pf:baselat1:limitpow} for $\sigma$ and
$m=\ell\Nm v$. Then $z\in I_v$ (a closed subgroup), the condition on powers of $z$ in the definition of tame inertia
elements holds, and, the kernel of $\delta_i$ being open, $\delta_i(z)=\delta_i(\sigma)^{(\ell\Nm v)^{k}}$ for some
$k$. Since $\Nm v\equiv1\pmod{15}$ (a property of $\delta_i$) and $p\ne\ell$, $p$ is prime to $2(\ell\Nm v)^k$, so
$\delta_i(z)^2$ is a primitive $p$-th root of unity. Hence $\delta_i(z)^2-1$ divides $p=\prod_{j=1}^{p-1}(1-\zeta_p^j)$ in
$\baselatoneO_0$; as $p$ and $\ell$ are coprime integers and $\ell$ is not a unit in the local ring $\baselatoneO_0$,
$p$ is a unit, and so is $\delta_i(z)^2-1$.
\end{proof}

\begin{lemma}\label{pf:baselat1:isQ}
For every finite family $P\colon\iota\to\{\text{places of }E\}$ the set $Q(P)$ is finite with
$|Q(P)|\le2|\iota|$, and it is closed under conjugation of places.
\formallabel{GlobalTheory.TwistChoice.exists_isQ}\formallabel{GlobalTheory.TwistChoice.conjCl_of_isQ}
\end{lemma}

\begin{proof}
A place has at most one conjugate place; let $c_q$ be the conjugate of $P_q$ if it exists. Then
$Q(P)=\{P_q\}\cup\{c_q\}$ has the required properties.
\end{proof}

\begin{lemma}\label{pf:baselat1:psiok}
Let $M,k\ge1$, $A_0$ a finite set of places, $(P,d)$ a family with $\mathrm{Aux}(i,M,A_0,P,d)$, $Q=Q(P)$,
and $e\colon\iota\to\N$, $\eps>0$ with $e_q<k$ and $\ell^{e_q}/\varphi(\ell^k)\le\eps$ for all $q$.
Let $\zeta\in\baselatoneO_k^\times$ be a primitive $\ell^k$-th root of unity in the subgroup generated by $\zeta_k$, and
$\psi'\colon\mathrm{Cl}(i,Q)\to\baselatoneO_k^\times$ a homomorphism such that
$\psi'\circ\mathrm{art}_Q=\chi\circ\delta_U(i,Q,M)$ for a character $\chi$ and
$\psi'(\mathrm{art}_Q(g))=\prod_q(\zeta^{\ell^{e_q}})^{d_q(g)}$ for all $g$. Then $\mathrm{T}_0(i,A_0,Q,k,\psi')$ and
$\sum_{v\in Q}\mathrm{tv}_v(\psi'\circ\mathrm{art}_Q)\le|Q|\,\eps$.
\formallabel{GlobalTheory.TwistChoiceAll.psiOk_partD_root}
\end{lemma}

\begin{proof}
Since $e_q<k$, $\zeta^{\ell^{e_q}}$ is a primitive $\ell^{k-e_q}$-th root of unity, so it is $\ne1$ and
$\baselatonenv(\zeta^{\ell^{e_q}}-1)=1/\varphi(\ell^{k-e_q})=\ell^{e_q}/\varphi(\ell^k)\le\eps$
(\cref{pf:baselat1:rootval}). The conditions of $\mathrm{T}_0$: the places of $Q$ are split, of degree and
ramification index one, and avoid $S_0(i)\cup A_0$ by $\mathrm{Aux}$; $k\ge1$; the kernel of
$\psi'\circ\mathrm{art}_Q$ contains the open kernel of $\mathrm{art}_Q$; anti-invariance under $c$ holds for
$\mathrm{art}_Q$; the values are products of powers of $\zeta$; if $\psi$ is ramified at $v$ then $v\in Q$, because
$\psi'\circ\mathrm{art}_Q$ is trivial on inertia outside $Q$ (property of $\mathrm{art}_Q$); conversely for $v\in Q$,
equal or conjugate to $P_q$, since $d_q$ maps inertia at $v$ onto $\Z/\ell^M$ and the $d_{q'}$, $q'\ne q$, are
trivial there, there is an inertia element $\sigma$ with $\psi(\sigma)=\zeta^{\ell^{e_q}}\ne1$; and $\psi$ is trivial
above $\ell$ because $\delta_U(i,Q,M)$ is. Finally, for $v\in Q$ the element $\sigma$ just found gives
$\mathrm{tv}_v(\psi)\le\baselatonenv(\zeta^{\ell^{e_q}}-1)\le\eps$, and we sum over $Q$.
\end{proof}

\begin{theorem}[Choice of the twisting character]\label{pf:baselat1:choice}
Let $D$ be a global datum and $N_0\colon X.\mathrm{Tw}\to\N$. Assume $\mathrm{CT}$, $(\mathrm{E}_{\mathrm{ind}})$,
$(\mathrm{S}_5)$ and $\mathcal{A}(N_0)$. Then
$\mathrm{AIB}(N_0)$: for every nonzero index place $i$ with $\mathrm{SH}(i)$, every
$\eps'>0$, every family $\mathcal F$ satisfying $\mathrm{AB}$ for all $j\in[1,m]$ and $b$, and every finite set $A_0$
of places, there are $Q$, $k'$ and $\psi'\colon\mathrm{Cl}(i,Q)\to\baselatoneO_{k'}^\times$ such that for every
$a\in\Z$ prime to $\ell$
\begin{multline*}
\mathrm{T}_0(i,A_0,Q,k',\psi'^a)\wedge\mathrm{T}_{\mathrm{Sel}}(i,Q,k',\psi'^a,\eps')\wedge\mathrm{T}_{\Theta}(N_0,i,Q,k',\psi'^a,\eps')\\
\wedge\mathrm{T}_{\mathrm{tm}}(i,Q,k',\psi'^a,\eps')\wedge\mathrm{T}_{\mathrm{lat}}(\mathcal F,Q,k',\psi'^a,\eps').
\end{multline*}
\formallabel{GlobalTheory.TwistChoiceAll.admissibleIndexBoundAllPowers_of_conditions}
\end{theorem}

\begin{proof}
Fix $i$, $\eps'$, $\mathcal F$, $A_0$. \emph{Step 1 (auxiliary data).} Let $z$ be a tame inertia element at a twist
prime (\cref{pf:baselat1:tame}). Every good signature $s$ satisfies $\mathrm{GH}(i,s)$: $\mathrm{adm}_1$,
$\mathrm{adm}_2$, $\mathrm{adm}_3$ are in $\mathcal{A}$, $\mathrm{SD}$, $\mathrm{NJ}$ and $r\ge100$ are in
$\mathrm{SH}(i)$. Let $R,D_d,M_0\ge1$ be the constants of
$(\mathrm{S}_5)$ for $\mathcal T=\Sigma^{\mathrm g}(i)$. For every label $\lambda$,
$(\mathrm{E}_{\Theta})$ (\lref{GlobalTheory.ToricFactor.propE}, from $\mathrm{adm}_9$, $\mathrm{adm}_{10}$, the condition
$s_2(\lambda)=2s'$ and $\sum_\delta\vartheta(\emptyset,\lambda,\delta)\ne0$ in $\mathrm{SH}$, and $z$) gives an open
normal subgroup $H_E(\lambda)$ and $v_E(\lambda)$; for every $y=(j,b)$ with $j\in[1,m]$, $(\mathrm{E}_{\mathrm{ind}})$ gives
$H_B(y)$, a finite set $A_B(y)$ and constants $v_B(y),n_B(y)$. Let $H=\bigcap_\lambda H_E(\lambda)\cap\bigcap_yH_B(y)$, an
open normal subgroup (finite intersection), and $A=A_0\cup\bigcup_yA_B(y)$. For every $M\ge M_0$, $\mathrm{CT}$ gives an
admissible family $(H_0,S_0,d^M)$ whose characters are the coordinates of $\delta_U(i,\cdot,M)$, and
$(\mathrm{S}_5)$ (for this family, $z$, $H$ and $A$) gives $P^M$ with
$\mathrm{Aux}(i,M,A,P^M,d^M\circ P^M)$ and $\mathrm{Adp}(i,M,H,z,P^M,d^M\circ P^M)$. Let $Q^M$ be as in
\cref{pf:baselat1:isQ}, $|Q^M|\le4R$. For $s\in\Sigma^{\mathrm g}(i)$ let $F^M_s$ be the Selmer polynomial with a coefficient of
degree $\le D_d$ and valuation $\le\mu:=\nu(i)\,v_{\baselatoneO_0}(\ell)$. Since $\mathrm{Aux}$ and $\mathrm{Adp}$ are
monotone in $A$ and in $H$, $(\mathrm{E}_{\Theta})$ and $(\mathrm{E}_{\mathrm{ind}})$ apply to $(P^M,d^M,Q^M)$ and give polynomials
$f^M_\lambda$ and $f^M_y$ satisfying $\mathrm{OH}$ with the bounds $(2,v_E(\lambda))$ and $(n_B(y),v_B(y))$, hence
with $n_{\max}=2+\sum_yn_B(y)$.

\emph{Step 2 (evaluation).} Let $C=|\Sigma^{\mathrm g}(i)|+2m+4tR+1$ and $\eps=\eps'/C$; note $c\,\eps\le\eps'$ for every $c\le C$.
Apply \cref{pf:baselat1:orderEval} (in the form \cref{pf:baselat1:existsLevel}) to the polynomials $F^M_s$ (indexed by
$I=\Sigma^{\mathrm g}(i)$, with bound $\mu$) and $f^M_\lambda,f^M_y$ (indexed by $J=\mathcal L_i\sqcup([1,m]\times\{\pm\})$) and to
$\eps$. We obtain $M\ge M_0$, $1\le k\le M$ and $e\colon\iota\to\N$ with $e_q<k$ and $\ell^{e_q}/\varphi(\ell^k)\le\eps$ such
that for every primitive $\ell^k$-th root of unity $\zeta\in\baselatoneO_k$, writing $u_q=\zeta^{\ell^{e_q}}$, all values
$F^M_s(u-1)$, $f^M_\lambda(u-1)$, $f^M_y(u-1)$ are nonzero with
\[
\baselatonenv(F^M_s(u-1))\le\nu(i)+\eps,\quad \baselatonenv(f^M_\lambda(u-1))\le v_E(\lambda)/v(\ell)+\eps,\quad
\baselatonenv(f^M_y(u-1))\le v_B(y)/v(\ell)+\eps.
\]
Write $P=P^M$, $Q=Q^M$, $d=d^M$. Let $\chi_0$ be the character of $(\Z/\ell^M)^Q$ with
$\chi_0(x)=\prod_q(\zeta_k^{\ell^{e_q}})^{x(P_q)}$; it is well defined since $(\zeta_k^{\ell^{e_q}})^{\ell^M}=1$ ($k\le M$).
The character $\delta_U(i,Q,M)$ factors through $\mathrm{art}_Q$ (\lref{GlobalTheory.TwistChoice.exists_factor}, by the
universal property of $\mathrm{art}_Q$, as $Q\cap S_0(i)=\emptyset$): $\delta_U=\phi\circ\mathrm{art}_Q$. Put $k'=k$ and
$\psi'=\chi_0\circ\phi$.

\emph{Step 3 (verification).} Let $a$ be prime to $\ell$. Then $\zeta=\zeta_k^a$ is a primitive $\ell^k$-th root of unity in
$\langle\zeta_k\rangle$, $\chi=\chi_0^a$ satisfies $\chi(x)=\prod_q(\zeta^{\ell^{e_q}})^{x(P_q)}$ and
$\psi'^a\circ\mathrm{art}_Q=\chi\circ\delta_U(i,Q,M)$; since the coordinates of $\delta_U$ are the $d_q$ ($\mathrm{CT}$, $Q$
being conjugation closed), $\psi'^a(\mathrm{art}_Q(g))=\prod_q(\zeta^{\ell^{e_q}})^{d_q(g)}$. The same holds with $\zeta^{-1}$
in place of $\zeta$ for $(\psi'^a)^{-1}$, and $\zeta^{-1}$ is again primitive. Now:
\begin{itemize}
\item $\mathrm{T}_0$ and the bound $\sum_{v\in Q}\mathrm{tv}_v\le|Q|\eps$ follow from \cref{pf:baselat1:psiok} (with
  $\mathrm{Aux}$ for $A_0\subseteq A$).
\item $\mathrm{T}_{\mathrm{Sel}}$: take $\eps_s=\eps$; then $\sum_{s\in\Sigma^{\mathrm g}(i)}\eps_s=|\Sigma^{\mathrm g}(i)|\eps\le\eps'$. For good $s$,
  since $F^M_s$ is a Selmer polynomial and $(\psi'^a\circ\mathrm{art}_Q)^{(b_s)}=\prod_qu_q^{d_q}$ with $u_q=(\zeta^{\pm1})^{\ell^{e_q}}$,
  $L_{\mathrm{Sel}}=\baselatonenv(F^M_s(u-1))\le\nu(i)+\eps$.
\item $\mathrm{T}_{\Theta}$: by $(\mathrm{E}_{\Theta})$, $f^M_\lambda(u-1)=\ell^{N_0(i)}\Theta_Q(\lambda,\psi'^a)$ for the appropriate sign.
  As $f^M_\lambda(u-1)\ne0$, $\Theta_Q(\lambda,\psi'^a)\ne0$ and
  $\baselatonenv(\Theta_Q)=\baselatonenv(f^M_\lambda(u-1))-N_0(i)$ (\lref{GlobalTheory.TwistChoice.vK_bounds}); this is
  $\ge-N_0(i)$ since $\baselatonenv\ge0$ on $\baselatoneO_k$, and $\le v_E(\lambda)/v(\ell)+\eps-N_0(i)=c_0(\lambda)+\eps\le
  c_0(\lambda)+\eps'$.
\item $\mathrm{T}_{\mathrm{tm}}$: $t\sum_{v\in Q}\mathrm{tv}_v\le t|Q|\eps\le4tR\,\eps\le\eps'$.
\item $\mathrm{T}_{\mathrm{lat}}$: take $M$, $\chi$ as above and $\eps_{j,b}=\eps$, so that $\sum\eps_{j,b}=2m\eps\le\eps'$; for every
  lattice $T_\psi$ as in $\mathrm{T}_{\mathrm{lat}}$, $(\mathrm{E}_{\mathrm{ind}})$ gives
  $\operatorname{ind}_{\mathrm n}(\baselatoneO_k\otimes B_{j,b},T_\psi)\le\operatorname{ind}_{\mathrm n}(B_{j,b},\Theta_{j,b})+\baselatonenv(f^M_y(u-1))-v_B(y)/v(\ell)
  \le\operatorname{ind}_{\mathrm n}(B_{j,b},\Theta_{j,b})+\eps$.\qedhere
\end{itemize}
\end{proof}

%======================================================================================================
\subsection{Evaluation of polynomials at roots of unity}\label{pf:baselat1:ssec:roots}

This subsection proves the evaluation theorem used in Step~2 of the proof of \cref{pf:baselat1:choice}. It is
elementary commutative algebra. Throughout, $\baselatoneO$ is a discrete valuation ring with uniformiser $\varpi$
and valuation $v$, and $\ell$ is a prime number with $\ell\ne0$ and $\ell$ not a unit in $\baselatoneO$, so that
$\varpi\mid\ell$; we write $\baselatonenv(x)=v(x)/v(\ell)$. For a finite set $\sigma$ of variables, a monomial
$m\in\N^\sigma$ has degree $|m|=\sum_qm_q$, and for $e\in\N^\sigma$ its \emph{$e$-weight} is
$w_e(m)=\sum_qm_q\ell^{e_q}\ge|m|$. Let $S_e\colon\baselatoneO[X_q:q\in\sigma]\to\baselatoneO[X]$ be the
homomorphism with $S_e(X_q)=(1+X)^{\ell^{e_q}}-1$. For $f\in\baselatoneO[X_q]$ and $\zeta$ in an
$\baselatoneO$-algebra, $S_e(f)(\zeta-1)=f\bigl((\zeta^{\ell^{e_q}}-1)_q\bigr)$.

\subsubsection{Binomial coefficients and substitutions}

\begin{lemma}\label{pf:baselat1:binom}
Let $a,c\in\N$ and $1\le i<\ell^{c+1}$. Then $\ell^{a-c}$ divides $\binom{\ell^a}{i}$ (with $a-c$ truncated at $0$).
\formallabel{GlobalTheory.pow_sub_dvd_choose_prime_pow_of_lt}\formallabel{GlobalTheory.T.prime_pow_sub_dvd_choose_prime_pow_of_lt}
\end{lemma}

\begin{proof}
Write $i=\ell^{v}m$ with $\ell\nmid m$; then $v\le c$ since $\ell^v\le i<\ell^{c+1}$. From
$i\binom{\ell^a}{i}=\ell^a\binom{\ell^a-1}{i-1}$ we get $\ell^a\mid\ell^vm\binom{\ell^a}{i}$, hence (if $v\le a$)
$\ell^{a-v}\mid m\binom{\ell^a}{i}$, and as $m$ is prime to $\ell$, $\ell^{a-v}\mid\binom{\ell^a}{i}$; this is trivial
if $v>a$. Finally $a-c\le a-v$.
\end{proof}

\begin{lemma}\label{pf:baselat1:subst}
\begin{enumerate}[label=(\roman*)]
\item $(1+X)^{\ell^e}-1\equiv X^{\ell^e}\pmod{\varpi}$ for every $e\in\N$; in particular
  $\prod_q S_e(X_q)^{m_q}\equiv X^{w_e(m)}\pmod\varpi$.
\item If $n+T\le a$, the coefficients of $(1+X)^{\ell^a}-1$ in degrees $\le T$ are divisible by $\varpi^n$.
\item The constant term of $S_e(g)$ is $g(0)$, and reduction modulo $\varpi$ commutes with $S_e$.
\item Let $k$ be a field of characteristic $\ell$, $D,u\in k[X_q]$ and $n\in\N^\sigma$ with $u(0)\ne0$,
  $D-u\prod_qX_q^{n_q}$ in the ideal generated by the $X_q^{\ell^M}$, and $w_e(n)<\ell^M$. Then the coefficient of
  $X^{w_e(n)}$ in $S_e(D)$ is $u(0)\neq0$.
\end{enumerate}
\formallabel{GlobalTheory.OrderEval.pow_sub_one_mem_lowIdeal}\formallabel{GlobalTheory.OrderEval.coeff_zero_subst}%
\formallabel{GlobalTheory.OrderEval.map_subst}\formallabel{GlobalTheory.coeff_aeval_one_add_X_pow_sub_one_ne_zero}%
\formallabel{GlobalTheory.T.coeff_aeval_one_add_X_pow_sub_one_ne_zero'}
\end{lemma}

\begin{proof}
(i) The coefficient of $X^i$, $0<i<\ell^e$, is $\binom{\ell^e}{i}$, which is divisible by $\ell$, hence by
$\varpi$. (ii) For $1\le i\le T$ we have $i<\ell^{T+1}$, so $\ell^{a-T}\mid\binom{\ell^a}{i}$ by
\cref{pf:baselat1:binom}, and $a-T\ge n$. (iii) $S_e(X_q)$ has constant term $0$; the second statement is
checked on monomials. (iv) In characteristic $\ell$, $S_e(X_q)=X^{\ell^{e_q}}$, so $S_e(X_q^{\ell^M})$ is divisible
by $X^{\ell^M}$ and the ideal does not contribute to coefficients of degree $<\ell^M$. Moreover
$S_e(u\prod_qX_q^{n_q})=S_e(u)X^{w_e(n)}$, and the constant term of $S_e(u)$ is $u(0)$.
\end{proof}

\begin{lemma}[Coefficient of minimal weight]\label{pf:baselat1:minweight}
Let $F\in\baselatoneO[X_q:q\in\sigma]$, $e\in\N^\sigma$, $m_0$ a monomial, $\mu\in\N$ and $W=w_e(m_0)$. Assume that
$\varpi^\mu$ divides the coefficient $F_m$ for every $m$ with $|m|\le W$, that $\varpi^{\mu+1}\nmid F_{m_0}$, and that
$\varpi^{\mu+1}\mid F_m$ for every $m\ne m_0$ with $w_e(m)\le W$. Then $\varpi^{\mu+1}$ does not divide the coefficient of
$X^W$ in $S_e(F)$.
\formallabel{GlobalTheory.not_dvd_coeff_aeval_oneAddX_pow_sub_one}\formallabel{GlobalTheory.T.not_dvd_coeff_aeval_one_add_X_pow_sub_one}
\end{lemma}

\begin{proof}
The coefficient in question is $\sum_mF_mc_m$, where $c_m$ is the coefficient of $X^W$ in $\prod_qS_e(X_q)^{m_q}$. By
\cref{pf:baselat1:subst}(i), $c_m\equiv1$ or $0\pmod\varpi$ according as $w_e(m)=W$ or not; and since $X$ divides every
$S_e(X_q)$, $c_m=0$ if $|m|>W$. Let $m\ne m_0$. If $w_e(m)\le W$, then $\varpi^{\mu+1}\mid F_m$. If $w_e(m)>W$ and
$|m|\le W$, then $\varpi\mid c_m$ and $\varpi^\mu\mid F_m$. If $|m|>W$ then $c_m=0$. In all cases
$\varpi^{\mu+1}\mid F_mc_m$. Moreover $F_{m_0}c_{m_0}-F_{m_0}=F_{m_0}(c_{m_0}-1)$ is divisible by $\varpi^{\mu}\cdot\varpi$,
because $|m_0|\le w_e(m_0)=W$. Hence the coefficient is congruent to $F_{m_0}$ modulo $\varpi^{\mu+1}$.
\end{proof}

\begin{lemma}[Least weight]\label{pf:baselat1:leastweight}
Let $\rho\in\N$ and $\emptyset\ne S\subseteq\N^\rho$. There is $m\in S$ such that for every $c\in\N^\rho$ with
$\sum_{i'<i}m_{i'}c_{i'}<c_i$ for all $i$, one has $\sum_im_ic_i<\sum_ia_ic_i$ for all $a\in S$, $a\ne m$.
\formallabel{GlobalTheory.exists_exponent_strictly_least_weight}\formallabel{GlobalTheory.T.exists_least_weight_exponent}
\end{lemma}

\begin{proof}
Order $\N^\rho$ lexicographically, comparing first the coordinates of largest index; this is a well-order, and we take
$m$ minimal in $S$. Let $a\in S$, $a\ne m$, and let $k$ be the largest index with $a_k\ne m_k$; minimality of $m$ gives
$m_k<a_k$. Then
\[
\sum_im_ic_i=\sum_{i<k}m_ic_i+m_kc_k+\sum_{i>k}a_ic_i<(m_k+1)c_k+\sum_{i>k}a_ic_i\le\sum_ia_ic_i.\qedhere
\]
\end{proof}

\begin{lemma}[Limit patterns]\label{pf:baselat1:limit}
Let $\sigma$, $I$ and $\mathcal R$ be finite sets and $F\colon\N\times I\times\N^\sigma\to\mathcal R$. There is
$G\colon I\times\N^\sigma\to\mathcal R$ such that for all $N,M_0$ there is $M\ge M_0$ with $F(M,x,m)=G(x,m)$ for all $x$
and all $m$ with $|m|\le N$.
\formallabel{GlobalTheory.exists_limit_family_of_finite_coefficients}\formallabel{GlobalTheory.T.exists_limit_family_of_finite_coefficients'}
\end{lemma}

\begin{proof}
Give $\mathcal R$ the discrete topology. The product space $\mathcal R^{I\times\N^\sigma}$ is compact, so the sequence
$M\mapsto F(M,\cdot,\cdot)$ has a cluster point $G$. For given $N$ the set of functions agreeing with $G$ at the finitely
many arguments $(x,m)$ with $|m|\le N$ is a neighbourhood of $G$, so it contains $F(M,\cdot,\cdot)$ for infinitely many
$M$.
\end{proof}

\subsubsection{Low coefficients after substitution}

\begin{proposition}\label{pf:baselat1:lowlimit}
Let $\sigma$ be a finite set enumerated as $\sigma=\{\tau_0,\dots,\tau_{\rho-1}\}$, $I$ a finite set, $D,M_0\in\N$,
$\mu\colon I\to\N$, and $F_M(x)\in\baselatoneO[X_q:q\in\sigma]$ for $M\in\N$, $x\in I$, such that for $M\ge M_0$ and every
$x$ some monomial of degree $\le D$ has coefficient in $F_M(x)$ not divisible by $\varpi^{\mu(x)+1}$. Then there is $a\in\N$
such that for every $e\in\N^\sigma$ with $e_{\tau_{i'}}+a\le e_{\tau_i}$ for all $i'<i$ there is $W\in\N$ with the
following property: for every $M_1$ there is $M\ge M_1$ such that for every $x$ some coefficient of $S_e(F_M(x))$ of degree
$\le W$ is not divisible by $\varpi^{\mu(x)+1}$.
\formallabel{GlobalTheory.OrderEval.exists_low_coeff_of_limit}
\end{proposition}

\begin{proof}
Let $K\ge\mu(x)$ for all $x$. \Cref{pf:baselat1:limit}, applied to the truth values of ``$\varpi^j$ divides the coefficient
of $m$ in $F_M(x)$'' for $j\le K+1$, gives a pattern $G(x,m,j)$ such that for all $N,M_1$ there is $M\ge M_1$ with:
$\varpi^j\mid F_M(x)_m\Leftrightarrow G(x,m,j)$ for all $x$, all $|m|\le N$ and all $j\le K+1$. For every $x$ there are $j$
and $m$ with $\neg G(x,m,j+1)$ (take $j=\mu(x)$, an $M\ge M_0$ agreeing with $G$ up to degree $D$, and the monomial of the
hypothesis); let $\nu(x)\le\mu(x)$ be the least such $j$. By \cref{pf:baselat1:leastweight} (via the enumeration) choose
$m_0(x)$ in $\{m:\neg G(x,m,\nu(x)+1)\}$. Let $B$ bound all exponents of all $m_0(x)$ and put $a=\rho B$.

Let $e$ be as in the statement and $W=\max_xw_e(m_0(x))$. Given $M_1$, choose $M\ge M_1$ agreeing with $G$ up to degree
$W$. Fix $x$. The weights $c_i=\ell^{e_{\tau_i}}$ satisfy $\sum_{i'<i}m_0(x)_{\tau_{i'}}c_{i'}<c_i$, because
$\bigl(\sum_{i'<i}m_0(x)_{\tau_{i'}}\ell^{e_{\tau_{i'}}}\bigr)\ell^{\rho B}\le\rho B\,\ell^{e_{\tau_i}}<\ell^{\rho B}\ell^{e_{\tau_i}}$.
We apply \cref{pf:baselat1:minweight} to $F_M(x)$, $m_0(x)$ and $\nu(x)$, with $W(x)=w_e(m_0(x))\le W$. If $\nu(x)>0$,
minimality of $\nu(x)$ gives $G(x,m,\nu(x))$ for all $m$, hence $\varpi^{\nu(x)}\mid F_M(x)_m$ for $|m|\le W$; if
$\nu(x)=0$ this is trivial. Next $\varpi^{\nu(x)+1}\nmid F_M(x)_{m_0(x)}$ since $|m_0(x)|\le W$. Finally, if $m\ne m_0(x)$,
$w_e(m)\le W(x)$ and $\varpi^{\nu(x)+1}\nmid F_M(x)_m$, then $\neg G(x,m,\nu(x)+1)$ (as $|m|\le w_e(m)\le W$), and
\cref{pf:baselat1:leastweight} gives $w_e(m_0(x))<w_e(m)$, a contradiction. So $\varpi^{\nu(x)+1}$, and a fortiori
$\varpi^{\mu(x)+1}$, does not divide the coefficient of $X^{W(x)}$ in $S_e(F_M(x))$.
\end{proof}

\begin{lemma}\label{pf:baselat1:congmul}
Let $\mathcal D,g,f\in\baselatoneO[[X]]$ and $m_0,v,n_T,M_T\in\N$. Assume that some coefficient $\mathcal D_m$ with
$m\le m_0$ is a unit, $\varpi^{v+1}\nmid g_0$, $v+1\le n_T$, $(v+1)m_0\le M_T$, and that $\varpi^{n_T}$ divides the
coefficients of $f-\mathcal Dg$ in degrees $\le M_T$. Then some coefficient $f_i$ with $i\le(v+1)m_0$ is not divisible by
$\varpi^{v+1}$.
\formallabel{GlobalTheory.exists_low_coeff_not_dvd_of_congr_mul}\formallabel{GlobalTheory.T.exists_low_coeff_not_dvd_of_congr_mul'}
\end{lemma}

\begin{proof}
Suppose the contrary. Then $\varpi^{v+1}$ divides the coefficients of $\mathcal Dg$ in degrees $\le(v+1)m_0$. Let $m_1\le
m_0$ be minimal with $\mathcal D_{m_1}$ a unit, so $\varpi\mid\mathcal D_a$ for $a<m_1$. We show by induction on $j\le v+1$:
$\varpi^j\mid g_i$ whenever $i+jm_1\le(v+1)m_0$. For $j=0$ there is nothing to prove. For $j+1$ we use induction on $i$:
the coefficient of $X^{i+m_1}$ in $\mathcal Dg$ is $\mathcal D_{m_1}g_i+\sum_{a+b=i+m_1,\,a\ne m_1}\mathcal D_ag_b$. If $a<m_1$
then $\varpi\mid\mathcal D_a$ and $\varpi^j\mid g_b$ (as $b+jm_1\le i+(j+1)m_1$); if $a>m_1$ then $b<i$ and
$\varpi^{j+1}\mid g_b$ by the inner induction. Since $i+m_1\le(v+1)m_0$ and $j+1\le v+1$, $\varpi^{j+1}$ divides the whole
coefficient, hence $\mathcal D_{m_1}g_i$, hence $g_i$. For $j=v+1$ and $i=0$ (note $(v+1)m_1\le(v+1)m_0$) we get
$\varpi^{v+1}\mid g_0$, a contradiction.
\end{proof}

\begin{proposition}\label{pf:baselat1:lowfactor}
Let $R,n_{\max},V,s\in\N$ and $\sigma=\{1,\dots,R\}\sqcup\{1,\dots,R\}$. There are $E,W\in\N$ such that: for every
$e\in\N^\sigma$ with $e_{(1,a)}\le s$ and $e_{(2,a)}\ge E$ for all $a$, every $M\ge E$, every $v\le V$, and all
$f,D_d,g\in\baselatoneO[X_q]$ with $f-D_dg$ in the ideal generated by the $X_{(2,a)}$ and the $(1+X_{(1,a)})^{\ell^M}-1$,
$\varpi^{v+1}\nmid g(0)$, and $D_d\equiv u\prod_aX_{(1,a)}^{n_a}$ modulo $\varpi$ and the $X_q^{\ell^M}$ with $u(0)\not\equiv0$
and $n_a\le n_{\max}$, some coefficient of $S_e(f)$ of degree $\le W$ is not divisible by $\varpi^{v+1}$.
\formallabel{GlobalTheory.OrderEval.exists_low_coeff_of_factor}
\end{proposition}

\begin{proof}
Let $m_0=Rn_{\max}\ell^s$, $E=(V+1)+(V+1)m_0$ and $W=(V+1)m_0$. The polynomials whose coefficients in degrees $\le W$ are
divisible by $\varpi^{V+1}$ form an ideal, and $S_e$ maps the generators $X_{(2,a)}$ and $(1+X_{(1,a)})^{\ell^M}-1$ into it,
by \cref{pf:baselat1:subst}(ii), since $S_e(X_{(2,a)})=(1+X)^{\ell^{e_{(2,a)}}}-1$ with $e_{(2,a)}\ge E$ and
$S_e((1+X_{(1,a)})^{\ell^M}-1)=(1+X)^{\ell^{e_{(1,a)}+M}}-1$ with $M\ge E$. Hence $\varpi^{V+1}$ divides the coefficients of
$S_e(f)-S_e(D_d)S_e(g)$ in degrees $\le W$. The residue field has characteristic $\ell$ (as $\varpi\mid\ell$), and the
weight $w=\sum_an_a\ell^{e_{(1,a)}}$ satisfies $w\le m_0<\ell^M$. By \cref{pf:baselat1:subst}(iii),(iv) the coefficient of
$X^w$ in $S_e(D_d)$ is a unit. \Cref{pf:baselat1:congmul}, applied to $S_e(D_d),S_e(g),S_e(f)$ with $m_0$, $v$,
$n_T=V+1$ and $M_T=W$ (the constant term of $S_e(g)$ is $g(0)$), gives a coefficient of $S_e(f)$ of degree
$\le(v+1)m_0\le W$ not divisible by $\varpi^{v+1}$.
\end{proof}

\subsubsection{Valuations at roots of unity}

\begin{lemma}\label{pf:baselat1:rootval}
Let $\baselatoneO'$ be a discrete valuation ring in which $\ell\neq0$ is not a unit, $k\ge1$, and $\zeta\in\baselatoneO'$ a
primitive $\ell^k$-th root of unity. Then $\zeta\ne1$ and $\baselatonenv(\zeta-1)=1/\varphi(\ell^k)$. If
$\baselatoneO\to\baselatoneO'$ is an injective local homomorphism of discrete valuation rings, then
$\baselatonenv_{\baselatoneO'}(x)=\baselatonenv_{\baselatoneO}(x)$ for $x\in\baselatoneO$.
\formallabel{GlobalTheory.T.normalizedValuation_primitiveRoot_prime_pow_sub_one}
\formallabel{GlobalTheory.T.normalizedValuation_algebraMap_eq}
\end{lemma}

\begin{proof}
Since $\ell\ne0$ in the domain $\baselatoneO'$, the $\varphi(\ell^k)$ elements $\zeta^a$, $a\in(\Z/\ell^k)^\times$, are the
distinct roots of the cyclotomic polynomial $\Phi_{\ell^k}$, and $\Phi_{\ell^k}(1)=\ell$, so
$\ell=\prod_a(1-\zeta^a)$. For $a$ prime to $\ell$, $(1-\zeta^a)/(1-\zeta)=1+\zeta+\dots+\zeta^{a-1}$ is integral, and so is
its inverse (write $\zeta=(\zeta^a)^b$ with $ab\equiv1$), so all factors have the valuation of $1-\zeta$, which is therefore
$v(\ell)/\varphi(\ell^k)$ (and $\zeta\ne1$ since $\ell\ne0$); this is the argument of \cite[Lemma~1.4]{Washington1997},
given there for $k=1$. For the second statement,
the image of a uniformiser of $\baselatoneO$ is a unit times the $e$-th power of a uniformiser of $\baselatoneO'$, so
$v'(x)=e\,v(x)$ for all $x$, and the ratio $v'(x)/v'(\ell)$ is unchanged.
\end{proof}

\begin{proposition}[Valuation at $\zeta-1$]\label{pf:baselat1:valroot}
Let $h\in\baselatoneO[X]$ and $w\in\N$ with $h_w\ne0$. There are a finite set $\mathcal B\subseteq\N$ with
$|\mathcal B|\le w+v(h_w)$ and reals $c(k)$ such that for every $k\ge1$ with $k\notin\mathcal B$:
$c(k)\le\baselatonenv(h_w)+w/\varphi(\ell^k)$, and for every discrete valuation ring $\baselatoneO'$ with an injective local
homomorphism $\baselatoneO\to\baselatoneO'$ and every primitive $\ell^k$-th root of unity $\zeta\in\baselatoneO'$, one has
$h(\zeta-1)\ne0$ and $\baselatonenv(h(\zeta-1))=c(k)$.
\formallabel{GlobalTheory.PolyCycloValue.value_at_root_sub_one}\formallabel{GlobalTheory.PolyCycloValue.card_double_min}
\formallabel{GlobalTheory.T.valuation_polynomial_at_root_of_unity_sub_one}
\end{proposition}

\begin{proof}
Let $a_i=v(h_i)$ for $i$ in the support $\mathcal S$ of $h$, and $q_k=v(\ell)/\varphi(\ell^k)$; $q_k>0$ and $q_k$ is
strictly decreasing for $k\ge1$. By \cref{pf:baselat1:rootval} the term $h_i(\zeta-1)^i$ has normalised valuation
$(a_i+iq_k)/v(\ell)$. For each $k$ let $I(k)\in\mathcal S$ minimise $a_i+iq_k$, and put $c(k)=(a_{I(k)}+I(k)q_k)/v(\ell)\le
(a_w+wq_k)/v(\ell)=\baselatonenv(h_w)+w/\varphi(\ell^k)$.

Let $\mathcal B$ be the set of $k\ge1$ for which the minimum is attained at two indices $I_k<J_k$. Two observations: (1) if
$i$ is a minimiser at $k$ and $m\in\mathcal S$, $m<i$, then $a_i<a_m$ (because $a_i+iq_k\le a_m+mq_k$ and $mq_k<iq_k$);
(2) if $k<k'$, $m$ is a minimiser at $k$ and $m'$ at $k'$, then $m\le m'$ (otherwise adding the two minimality inequalities
gives $(m-m')(q_k-q_{k'})\le0$). Define $\phi(k)=I_k$ if $I_k<w$ and $\phi(k)=w+a_{J_k}$ otherwise. Then $\phi(k)<w+a_w$: in
the second case $J_k>I_k\ge w$, so $a_{J_k}<a_w$ by (1). And $\phi$ is injective on $\mathcal B$: let $k<k'$ in $\mathcal B$;
by (2) $J_k\le I_{k'}<J_{k'}$, so $a_{J_{k'}}<a_{J_k}$ by (1). If $I_k<w$ and $I_{k'}<w$ then $\phi(k)=I_k<J_k\le
I_{k'}=\phi(k')$; if $I_k<w\le I_{k'}$ then $\phi(k)<w\le\phi(k')$; if $I_k\ge w$ then $I_{k'}\ge J_k>w$ and
$\phi(k)=w+a_{J_k}>w+a_{J_{k'}}=\phi(k')$. Hence $|\mathcal B|\le w+a_w$.

For $k\notin\mathcal B$, $k\ge1$, the term $h_{I(k)}(\zeta-1)^{I(k)}$ has strictly smaller valuation than all other terms of
$h(\zeta-1)=\sum_{i\in\mathcal S}h_i(\zeta-1)^i$, so by the ultrametric inequality $h(\zeta-1)\ne0$ and
$\baselatonenv(h(\zeta-1))=c(k)$.
\end{proof}

\begin{corollary}\label{pf:baselat1:levels}
Let $J$ be a finite set, $v\colon J\to\N$ and $W\in\N$. There is $N\in\N$ such that for every family $(h_y)_{y\in J}$ in
$\baselatoneO[X]$ such that each $h_y$ has a coefficient of degree $w_y\le W$ not divisible by $\varpi^{v_y+1}$, there is a
finite set $\mathcal B$ with $|\mathcal B|\le N$ such that for every $k\ge1$ with $k\notin\mathcal B$, every $\baselatoneO'$
as in \cref{pf:baselat1:valroot}, every primitive $\ell^k$-th root of unity $\zeta\in\baselatoneO'$ and every $y$:
$h_y(\zeta-1)\neq0$ and $\baselatonenv(h_y(\zeta-1))\le v_y/v(\ell)+W/\varphi(\ell^k)$.
\formallabel{GlobalTheory.OrderEval.exists_levels}
\end{corollary}

\begin{proof}
Take $N=\sum_y(W+v_y)$ and $\mathcal B=\bigcup_y\mathcal B_y$, with $\mathcal B_y$ given by \cref{pf:baselat1:valroot} for
$h_y$ and $w_y$; since $v(h_{y,w_y})\le v_y$, $|\mathcal B_y|\le W+v_y$ and
$\baselatonenv(h_y(\zeta-1))\le v_y/v(\ell)+W/\varphi(\ell^k)$.
\end{proof}

\subsubsection{The evaluation theorem}

\begin{theorem}[Evaluation at roots of unity]\label{pf:baselat1:orderEval}
Let $\baselatoneO_0\to\baselatoneO$ be an injective local homomorphism of discrete valuation rings, $\baselatoneO_0$ with
finite residue field, and $\ell$ a prime with $\ell\ne0$ and $\ell$ not a unit in $\baselatoneO_0$. Let $R,D,n_{\max},M_0\in\N$,
$\sigma=\{1,\dots,R\}\sqcup\{1,\dots,R\}$, $I,J$ finite sets, $\mu\colon I\to\N$, $v\colon J\to\N$, and for $M\in\N$ let
$F_M(x)\in\baselatoneO_0[X_q:q\in\sigma]$ ($x\in I$) and $f_M(y),D_M(y),g_M(y)\in\baselatoneO[X_q]$ ($y\in J$) be such that for
$M\ge M_0$:
\begin{itemize}
\item $F_M(x)$ has a monomial of degree $\le D$ whose coefficient has valuation $\le\mu(x)$;
\item $f_M(y)$ satisfies the expansion conditions of \cref{pf:baselat1:lowfactor} with $D_M(y),g_M(y)$, $v(g_M(y)(0))\le v(y)$,
  $u(0)\ne0$ and exponents $\le n_{\max}$.
\end{itemize}
Then for every $\eps>0$ there are $M\ge M_0$, $1\le k\le M$ and $e\in\N^\sigma$ with $e_q<k$ and
$\ell^{e_q}/\varphi(\ell^k)\le\eps$ for all $q$, such that for every discrete valuation ring $\baselatoneO'$ with an injective
local homomorphism $\baselatoneO\to\baselatoneO'$ and every primitive $\ell^k$-th root of unity $\zeta\in\baselatoneO'$, with
$u_q=\zeta^{\ell^{e_q}}$:
\begin{gather*}
F_M(x)(u-1)\ne0,\qquad \baselatonenv(F_M(x)(u-1))\le\frac{\mu(x)}{v_{\baselatoneO_0}(\ell)}+\eps,\\
f_M(y)(u-1)\ne0,\qquad \baselatonenv(f_M(y)(u-1))\le\frac{v(y)}{v_{\baselatoneO}(\ell)}+\eps
\end{gather*}
for all $x\in I$ and $y\in J$.
\formallabel{GlobalTheory.OrderEval.order_eval}
\formallabel{GlobalTheory.T.exists_level_order_exponents_small_valuation_determinants}
\end{theorem}

\begin{proof}
The element $\ell$ is nonzero and not a unit in $\baselatoneO$ either, since the homomorphism is injective and local. Apply
\cref{pf:baselat1:lowlimit} over $\baselatoneO_0$ to the $F_M(x)$ (the hypothesis ``valuation $\le\mu(x)$'' means
``not divisible by $\varpi_0^{\mu(x)+1}$''), with the enumeration $(1,0),\dots,(1,R-1),(2,0),\dots,(2,R-1)$ of $\sigma$; let
$a$ be the resulting constant. Let $V\ge v(y)$ for all $y$, and let $E,W_2$ be given by \cref{pf:baselat1:lowfactor} over
$\baselatoneO$ for $R,n_{\max},V$ and $s=aR$. Put $e_{(1,i)}=ai$ and $e_{(2,i)}=a(i+R)+E$; then consecutive exponents in the
enumeration differ by at least $a$, $e_{(1,i)}\le aR$ and $e_{(2,i)}\ge E$. Let $W_1$ be given by \cref{pf:baselat1:lowlimit}
for this $e$, and $N_1$, $N_2$ by \cref{pf:baselat1:levels} for $(I,\mu,W_1)$ over $\baselatoneO_0$ and $(J,v,W_2)$ over
$\baselatoneO$. Let $C\ge\ell^{e_q}$ for all $q$ and $k_1\in\N$ with $k_1\ge(C+W_1+W_2)/\eps$. By \cref{pf:baselat1:lowlimit}
choose $M\ge M_0+E+k_1+C+1+N_1+N_2$ such that every $S_e(F_M(x))$ has a coefficient of degree $\le W_1$ not divisible by
$\varpi_0^{\mu(x)+1}$. By \cref{pf:baselat1:lowfactor} (applicable as $M\ge E$ and $v(y)\le V$), every $S_e(f_M(y))$ has a
coefficient of degree $\le W_2$ not divisible by $\varpi^{v(y)+1}$. Let $\mathcal B_1,\mathcal B_2$ be the exceptional sets of
\cref{pf:baselat1:levels} for the families $S_e(F_M(x))$ and $S_e(f_M(y))$; the interval $[k_1+C+1,k_1+C+1+N_1+N_2]$ has
$N_1+N_2+1$ elements, so it contains some $k\notin\mathcal B_1\cup\mathcal B_2$.

Then $1\le k\le M$; $e_q<\ell^{e_q}\le C<k$; and $\varphi(\ell^k)\ge k\ge k_1$ (as $\varphi(\ell^k)=\ell^{k-1}(\ell-1)\ge k$),
so $c/\varphi(\ell^k)\le\eps$ for every $c\le C+W_1+W_2$, in particular $\ell^{e_q}/\varphi(\ell^k)\le\eps$. For $\baselatoneO'$
and $\zeta$ as in the statement, $S_e(F_M(x))(\zeta-1)=F_M(x)(u-1)$ and $S_e(f_M(y))(\zeta-1)=f_M(y)(u-1)$, and
\cref{pf:baselat1:levels} (over $\baselatoneO_0\to\baselatoneO'$, resp.\ $\baselatoneO\to\baselatoneO'$) gives the
non-vanishing and the bounds with error terms $W_1/\varphi(\ell^k)\le\eps$ and $W_2/\varphi(\ell^k)\le\eps$.
\end{proof}

\begin{remark}\label{pf:baselat1:rem:finres}
The finiteness of the residue field of $\baselatoneO_0$ is not used in the proof of \cref{pf:baselat1:orderEval}.
\end{remark}

\begin{corollary}\label{pf:baselat1:existsLevel}
In the setting of \cref{pf:baselat1:choice}, let $i$ be a twist index, $R,D_d,n_{\max},M_0\in\N$, $I,J$ finite sets,
$\mu\in\N$, $v\colon J\to\N$, and for $M\ge M_0$ let $F_M(x)\in\baselatoneO_0[X_q]$ have a coefficient of degree $\le D_d$ and
valuation $\le\mu$, and let $f_M(y)\in\baselatoneO[X_q]$ satisfy the polynomial expansion conditions of
\cref{pf:baselat1:def:ordhyp} with $(M,R,n_{\max},v(y))$. Then for every $\eps>0$ there are $M\ge M_0$, $1\le k\le M$ and
$e$ with $e_q<k$ and $\ell^{e_q}/\varphi(\ell^k)\le\eps$, such that for every primitive $\ell^k$-th root of unity
$\zeta\in\baselatoneO_k^\times$ the values $F_M(x)((\zeta^{\ell^{e_q}}-1)_q)$ and $f_M(y)((\zeta^{\ell^{e_q}}-1)_q)$ in
$\baselatoneO_k$ are nonzero with normalised valuations at most $\mu/v_{\baselatoneO_0}(\ell)+\eps$ and
$v(y)/v_{\baselatoneO}(\ell)+\eps$.
\formallabel{GlobalTheory.TwistChoice.exists_level}
\end{corollary}

\begin{proof}
This is \cref{pf:baselat1:orderEval} for $\baselatoneO_0\to\baselatoneO\to\baselatoneO_k$ (the maps are injective and
local, the residue field of $\baselatoneO_0$ is finite, and $\ell\ne0$ is not a unit in $\baselatoneO_0$), with the sequences
extended by $0$ for $M<M_0$.
\end{proof}

%======================================================================================================
\subsection{Proofs of the quoted construction results}\label{pf:baselat1:ssec:deferred}

In this subsection we prove the main construction results which were quoted in
\cref{pf:baselat1:ssec:stages,pf:baselat1:ssec:export}. Each proof follows the formal proof step by step. The
ingredients of these proofs which are themselves long constructions are quoted, with their statements, as results
of the formal development \cite{LeanSource2026}.

\subsubsection{Gauge families}\label{pf:baselat1:sssec:gaugeproofs}

We keep the notation of \cref{pf:baselat1:def:gauge}. A gauge family $G$ is \emph{Haar} if it has property (a) of
\cref{pf:baselat1:def:gauge}, and \emph{basis invariant} if it has property (b). For $J\in M_M(E)$ and
$P\in\GL_M(E)$ we write $J^P={}^t\!\sigma(P)JP$.

\begin{definition}[Standard measures and step relations]\label{pf:baselat1:def:stdmeas}
Let $b$ be a basis of $E$ over $F$, and write $F_\infty=F\otimes_\Q\R$.
\begin{enumerate}[label=(\alph*)]
\item A \emph{family of standard local Haar measures} $s$ for $b$ consists of additive Haar measures $\lambda_\infty$ on
  $E\otimes_FF_\infty$, $\lambda_v$ on $E\otimes_FF_v$ for every finite place $v$, $\nu_\infty$ on $F_\infty$ and $\nu_v$
  on $F_v$, such that $\lambda_v$ gives volume $1$ to the $\OO_{F_v}$-span of $b$ and $\nu_v(\OO_{F_v})=1$.
\item Let $J\in M_M(E)$ be hermitian and nondegenerate and $\beta\in F^\times$. A measure $\nu$ on
  $(\A_F\otimes_FE)^M$ is \emph{standard for the quadric} $h_J(y,y)=\beta$ if there are measures $m_\infty$ and $m_v$ such that
  $m_\infty$ (resp.\ $m_v$) is the fibre measure at $\beta$ of $\lambda_\infty^M$ over $\nu_\infty$ (resp.\ of $\lambda_v^M$
  over $\nu_v$) along the map $y\mapsto h_J(y,y)$, the series $\sum_v|m_v(L_v)-1|$ converges, where $L_v$ is the
  $\OO_{F_v}$-span of $b$ in $(E\otimes_FF_v)^M$, and $\nu$ is the restricted product of $m_\infty$ and the $m_v$ with
  respect to the $L_v$.
\item \emph{Step data} $(J,x,\beta,B,J')$ of rank $M+1$ consist of a hermitian nondegenerate $J\in M_{M+1}(E)$,
  $\beta\in F^\times$, $x\in E^{M+1}$ with $h_J(x,x)=\beta$, and $B\in M_{(M+1)\times M}(E)$ with $h_J(x,B_{\cdot j})=0$ for
  all $j$, such that $J'=J^B={}^t\!\sigma(B)JB$ is hermitian and nondegenerate.
\item For step data and measures $\mu$ on $U_J$ and $\mu'$ on $U_{J'}$, the \emph{step relation}
  $\mathrm{SR}_s(\mu,\mu')$ means: for every measure $\nu$ which is standard for the quadric $h_J(y,y)=\beta$ and every
  isomorphism $e\colon U_{J,x}\to U_{J'}$ compatible with $B$ (i.e.\ $gB=B\,e(g)$), the triple
  $(\mu,e^{-1}_*\mu',\nu)$ is compatible in the sense of \cref{pf:baselat1:def:gauge}(c).
\end{enumerate}
\formal{GlobalTheory.GaugeEx.StdMeas}\formal{GlobalTheory.GaugeEx.IsStdX}%
\formal{GlobalTheory.GaugeEx.StepData}\formal{GlobalTheory.GaugeEx.StepRel}
\end{definition}

Step data are determined by $(J,x,B)$: $\beta=h_J(x,x)$ and $J'=J^B$. If $(\mu,\mu',\nu)$ is compatible and
$c,k\ge0$, then $((kc)\mu,c\mu',k\nu)$ is compatible, since both sides of the defining identity are multiplied by
$kc$; we use this scaling property below.

\begin{lemma}\label{pf:baselat1:haarinj}
Let $\mu$ be a Haar measure on a locally compact group and $c,c'\ge0$ with $c\mu=c'\mu$. Then $c=c'$.
\formallabel{GlobalTheory.GaugeEx.smul_haar_inj}
\end{lemma}

\begin{proof}
A compact neighbourhood $K$ of the identity has $0<\mu(K)<\infty$, and $c\mu(K)=c'\mu(K)$.
\end{proof}

\begin{lemma}[Step constants]\label{pf:baselat1:stepconst}
Let $E/F$ be a CM extension, $s$ a family of standard local Haar measures and $G$ a Haar gauge family. There are
constants $c(M,J,x,B)>0$ such that $\mathrm{SR}_s(c(M,J,x,B)\,G.\mu_{M+1,J},\,G.\mu_{M,J'})$ for all $M\ge2$ and all step
data $(J,x,\beta,B,J')$ of rank $M+1$.
\formallabel{GlobalTheory.GaugeEx.exists_const}
\end{lemma}

\begin{proof}
Fix $M\ge2$ and step data. By \lref{GlobalTheory.T.exists_unique_oneVectorStep_haarMeasure}, which states
that for $M\ge2$, step data of rank $M+1$ and a Haar measure $\mu'$ on $U_{J'}$ there is a Haar measure $\mu$ on $U_J$
with $\mathrm{SR}_s(\mu,\mu')$, and that any two measures step related to $\mu'$ are equal, applied to the Haar measure
$\mu'=G.\mu_{M,J'}$, there is such a $\mu$. Since $G.\mu_{M+1,J}$ is also a Haar measure, $\mu=c\,G.\mu_{M+1,J}$ with
$c>0$ (uniqueness of Haar measures). As step data are determined by $(J,x,B)$, the choice of $c$ only depends on
$(M,J,x,B)$; for $(M,J,x,B)$ not extending to step data put $c=1$.
\end{proof}

\begin{lemma}[Standard families]\label{pf:baselat1:stdfamily}
Let $E/F$ be a CM extension, $b$ a basis of $E$ over $F$ and $s$ a family of standard local Haar measures for $b$.
There are a Haar and basis invariant gauge family $G$ and constants $\alpha(M,J)>0$ such that
\begin{enumerate}[label=(\roman*)]
\item $\alpha(M,J^P)=\alpha(M,J)$ for all hermitian nondegenerate $J\in M_M(E)$ and $P\in\GL_M(E)$;
\item $\mathrm{SR}_s\bigl(\alpha(M+1,J)\,G.\mu_{M+1,J},\ \alpha(M,J')\,G.\mu_{M,J'}\bigr)$ for all $M\ge2$ and all step data
  $(J,x,\beta,B,J')$ of rank $M+1$.
\end{enumerate}
\formallabel{GlobalTheory.GaugeEx.exists_std_family}
\end{lemma}

\begin{proof}
By \lref{GlobalTheory.T.exists_basisInvariant_haar_family_adelicUnitary}, which states that for a CM
extension there is a Haar and basis invariant gauge family, we fix such a $G$, and by \cref{pf:baselat1:stepconst}
we fix step constants $c$ for $G$. We regard $c$ as a function with values in the group $\R_{>0}$ and check the two
hypotheses of \lref{GlobalTheory.T.exists_coboundary_of_invariant_unitary_step_constants}. That result states:
\emph{if a function $c$ on quadruples $(M,J,x,B)$ with values in a commutative group is (I) invariant under
compatible changes of basis, i.e.\ $c(M,J^P,x_2,B_2)=c(M,J,x,B)$ whenever $M\ge2$, $(J,x,\beta,B,J')$ are step data,
$P\in\GL_{M+1}(E)$, $Q\in\GL_M(E)$, $Px_2=x$ and $PB_2=BQ$, and (II) independent of paths, i.e.\
$c(M+1,J,x_1,B_1)\,c(M,J_1,y_1,C_1)=c(M+1,J,x_2,B_2)\,c(M,J_2,y_2,C_2)$ whenever $M\ge2$, $(J,x_i,\beta_i,B_i,J_i)$ and
$(J_i,y_i,\beta_i',C_i,J_0)$ are step data ($i=1,2$) and $B_1C_1=B_2C_2$, then there is $\alpha$ with
$\alpha(M,J^P)=\alpha(M,J)$ for hermitian nondegenerate $J$ and $\alpha(M+1,J)=c(M,J,x,B)\,\alpha(M,J')$ for all
$M\ge2$ and step data.}

(I) Let $M,(J,x,\beta,B,J'),P,Q,x_2,B_2$ be as in (I), put $J_2=J^P$ and $J_2'=J'^Q$, and write $c_1=c(M,J,x,B)$ and
$c_2=c(M,J_2,x_2,B_2)$. By basis invariance there are isomorphisms $e_P\colon U_J\to U_{J_2}$ and
$e_Q\colon U_{J'}\to U_{J_2'}$ compatible with $P$ and $Q$, and $e_{P*}G.\mu_{M+1,J}=G.\mu_{M+1,J_2}$,
$e_{Q*}G.\mu_{M,J'}=G.\mu_{M,J_2'}$. By
\lref{GlobalTheory.T.unitaryOneVectorStep_compatible_changeOfBasis}, which states that under these
assumptions $(J_2,x_2,\beta,B_2,J_2')$ are step data and that $\mathrm{SR}_s(\mu,\mu')$ implies
$\mathrm{SR}_s(e_{P*}\mu,e_{Q*}\mu')$, the relation $\mathrm{SR}_s(c_1G.\mu_{M+1,J},G.\mu_{M,J'})$ gives
$\mathrm{SR}_s(c_1G.\mu_{M+1,J_2},G.\mu_{M,J_2'})$. On the other hand $\mathrm{SR}_s(c_2G.\mu_{M+1,J_2},G.\mu_{M,J_2'})$
by the choice of $c$. By the uniqueness statement of
\lref{GlobalTheory.T.exists_unique_oneVectorStep_haarMeasure} the two measures on $U_{J_2}$ coincide, and
\cref{pf:baselat1:haarinj} gives $c_2=c_1$.

(II) Let the data be as in (II), with $J$ of size $M+2$ and $J_0$ of size $M$. The measures
$\tau_i=c(M+1,J,x_i,B_i)\,c(M,J_i,y_i,C_i)\,G.\mu_{M+2,J}$ and $\mu_i=c(M,J_i,y_i,C_i)\,G.\mu_{M+1,J_i}$ satisfy
$\mathrm{SR}_s(\mu_i,G.\mu_{M,J_0})$ and $\mathrm{SR}_s(\tau_i,\mu_i)$ (the step relation is preserved when both measures
are multiplied by the same constant, by the scaling property). By
\lref{GlobalTheory.T.unitaryStabilizer_measure_independent_of_path}, which states that in this situation, for
a Haar measure $\mu_0$ on $U_{J_0}$ and measures with $\mathrm{SR}_s(\mu_i,\mu_0)$ and $\mathrm{SR}_s(\tau_i,\mu_i)$ for
$i=1,2$, one has $\tau_1=\tau_2$, we get $\tau_1=\tau_2$, and \cref{pf:baselat1:haarinj} gives (II).

Hence there is $\alpha$ as in the quoted result. Property (i) is its first conclusion. For (ii), multiply the
relation $\mathrm{SR}_s(c(M,J,x,B)\,G.\mu_{M+1,J},G.\mu_{M,J'})$ by $\alpha(M,J')$ and use
$\alpha(M+1,J)=c(M,J,x,B)\,\alpha(M,J')$.
\end{proof}

\begin{theorem}[Gauges for CM extensions]\label{pf:baselat1:gaugeCM}
Let $F$ be totally real and $E/F$ a totally complex quadratic extension, with an $F$-automorphism $\sigma$ of $E$
inducing complex conjugation under every complex embedding of $E$. Then there is a gauge family $G$ with
$\mathrm{Gauge}(G)$.
\formallabel{GlobalTheory.GaugeEx.exists_isGauge_of_isCMExt}
\end{theorem}

\begin{proof}
Choose a basis $b$ of $E$ over $F$ and a family $s$ of standard local Haar measures for $b$; such a family exists
(\lref{GlobalTheory.T.nonempty_standardLocalHaarMeasures}). Let $G,\alpha$ be as in
\cref{pf:baselat1:stdfamily}. We use three further facts.
\begin{itemize}
\item There are constants $\kappa(n)>0$ ($n\in\N$) such that for $M\ge1$, $J\in M_M(E)$ hermitian nondegenerate,
  $\beta\in F^\times$, a nontrivial continuous unitary character $\psi$ of $\A_F$ trivial on $F$, and a measure $\nu$
  standard for the quadric $h_J(y,y)=\beta$, a measure $\nu_X$ is a $\psi$-measure for this quadric if and only if
  $\nu_X=\kappa(M)\,\nu$ (\lref{GlobalTheory.T.quadricFibreMeasure_eq_smul_normalisedProduct}).
\item Since $F$ is totally real, a nontrivial continuous unitary character $\psi$ of $\A_F$ trivial on $F$ exists
  (\lref{GlobalTheory.T.exists_nontrivial_adeleChar_trivial_on_totallyReal}).
\item For $M+1\ge3$, $J$ hermitian nondegenerate and $\beta\ne0$ there is a unique measure $\nu$ standard for the
  quadric $h_J(y,y)=\beta$ (\lref{GlobalTheory.T.existsUnique_stdHermitianFibreProduct_of_rank_ge_three}); and for
  step data $(J,x,\beta,B,J')$ there is an isomorphism $U_{J,x}\to U_{J'}$ compatible with $B$
  (\lref{GlobalTheory.T.unitary_stabilizer_anisotropic_vector_equiv_complement}).
\end{itemize}
Put $K(M)=\prod_{i=0}^{M}\kappa(i)$, so that $K(M)>0$ and $K(M+1)=\kappa(M+1)K(M)$, and define
$G'.\mu_{M,J}=K(M)\,\alpha(M,J)\,G.\mu_{M,J}$. We verify (a)--(c) of \cref{pf:baselat1:def:gauge} for $G'$.

(a) A positive multiple of a Haar measure is a Haar measure.

(b) Let $J'=J^P$. By basis invariance of $G$ there is an isomorphism compatible with $P$, and every such $e$ satisfies
$e_*G.\mu_{M,J}=G.\mu_{M,J'}$; multiplying by $K(M)\alpha(M,J)=K(M)\alpha(M,J')$ (\cref{pf:baselat1:stdfamily}(i)) gives
$e_*G'.\mu_{M,J}=G'.\mu_{M,J'}$.

(c) Let $M+1\ge3$ and $(J,x,\beta,B,J')$ be as in \cref{pf:baselat1:def:gauge}(c), i.e.\ step data of rank $M+1$; then
$M\ge2$. Let $\nu$ be the standard measure for the quadric, and let $\psi$ be a character as above. For every
character $\psi'$ of this kind, $\nu_X=\kappa(M+1)\nu$ is a $\psi'$-measure for the quadric, and every
$\psi'$-measure for the quadric equals $\kappa(M+1)\nu$. An isomorphism $e\colon U_{J,x}\to U_{J'}$ compatible with $B$
exists. Finally let $e'$ be any isomorphism compatible with $B$. By \cref{pf:baselat1:stdfamily}(ii) the triple
$\bigl(\alpha(M+1,J)G.\mu_{M+1,J},\ e'^{-1}_*(\alpha(M,J')G.\mu_{M,J'}),\ \nu\bigr)$ is compatible, and the scaling property
with $c=K(M)$ and $k=\kappa(M+1)$ shows that
$\bigl(G'.\mu_{M+1,J},\ e'^{-1}_*G'.\mu_{M,J'},\ \kappa(M+1)\nu\bigr)$ is compatible, because
$\kappa(M+1)K(M)=K(M+1)$.
\end{proof}

\subsubsection{Lattice layers and kernel data}\label{pf:baselat1:sssec:latproofs}

\begin{definition}[Lattice instances and kernel records]\label{pf:baselat1:def:latinst}
For a column witness $w$, a level $Z$ and column data $t$ at split primes, a \emph{served place} is a triple
$(u,v,h)$ of a word $u$, a place $v$ and an index $h$ at the residue characteristic of $v$ for the type of $u$; to it
the datum $w.D(Z)$ attaches a local group datum $\Gamma(u,v,h)$. A \emph{lattice instance} $G_y$ for $w.D(Z)$
consists of a valuation map, fillings $\mathrm{Fl}(u,v,h)$ of the local group data, levels of the fillings and the
maps of the column comparison of $t$; the \emph{lattice-layer condition} $\mathrm{Lay}(K,w.D(Z),t,G_y)$ consists of
the condition on the subspaces $V_{\mathrm{ram}}$, the identity
$\mathrm{Fl}(u,v,h).w(x,j,s)=w_j(\text{position of }(x,s))-\kappa_J(j)$ for the weights, and a lattice condition for
all levels, chains and submodules. A \emph{kernel record} $k_t\in\mathrm{Rec}_{\ker}(t,Y,G_y,\mathrm{ker})$ consists of
lifted theta objects $k_t.O$ which are standard and unfolded, kernel entries $k_t.E$ which are admissible for $k_t.O$,
and four conditions relating $\mathrm{ker}$ to $k_t.O$ at the served places.
\formal{GlobalTheory.Mem.Top.GyLay}\formal{GlobalTheory.Mem.Top.Bl.KerT}
\end{definition}

\begin{proposition}[Lattice layers and kernel data]\label{pf:baselat1:latlayer}
Let $L$ be the statement: $\mathrm{Step}^{\mathrm w}(\mathrm{Pin},\cdot)$ of the property of $w$ that for every level $Z$,
all column data $t$ at split primes, every set of places $W_t$, every norm choice $\mathrm{nch}$ with a norm-choice datum
and every $Y\in\mathcal D(w.D(Z))$ such that $W_t$ and the characters $w.I.\xi$ satisfy the $\mu$-condition of cores and
$Y.Q_p$ is the set determined by $w$, there are a lattice instance $G_y$ and a kernel family $\mathrm{ker}$ with
$\mathrm{Lay}(K,w.D(Z),t,G_y)$, a kernel record $k_t\in\mathrm{Rec}_{\ker}(t,Y,G_y,\mathrm{ker})$, entry places $P_l$, place
matrices $B_s$ and a kernel-entry hypothesis $h$ with $k_t.E=E^{\mathrm{can}}(P_l,k_t.O,B_s,h)$ (the canonical
entries). Then $L$ holds.
\formallabel{GlobalTheory.T.exists_latticeLayer_kernelDatum_trueColumns}
\formallabel{GlobalTheory.Mem.Top.Bl.gyLay_kerT_entC}
\end{proposition}

\begin{proof}
\emph{Cloud bound.} By \lref{GlobalTheory.T.columnChoice_cloudBound_le_cG},
$\mathrm{Step}^{\mathrm w}(\mathrm{Pin},q)$, where $q(w)$ states that for every level $Z$ there is $c\in\N$ such that all
clouds of $w.D(Z)$ are bounded by $c$ and $c\le c_G(L')$ for every size datum $L'$ of $w.D(Z)$ at the scale $K.s$ which is
valid for the constants of $K$. By \cref{pf:baselat1:stepcalc}(v), (iii) and (iv) it suffices to show, for fixed
$K,B,f,T$ and a pinned $w$ over $K$ with $\mathrm{Fd}(K,B,T)$, that $q(w)$ implies the property in $L$. Fix
$Z,t,W_t,\mathrm{nch}$, a norm-choice datum and $Y$ as in $L$, and a column pin of $w$.

\emph{Objects and entries.} By \lref{GlobalTheory.Mem.Col2.TC.lift_objects}, for the column input data of $w$,
its laws, the level $Z$ and the place sets $Y.Q_p$ there are lifted theta objects $O$ which are standard and
unfolded in all indices $i<m$. Choose entry places $P_l$ (they exist,
\lref{Formalism.BasicNotions.Examples₂.entPl_nonempty}). By
\lref{GlobalTheory.T.trueColumnChoices_kernelDatum_adaptedBases}, which states that for standard unfolded lifted
objects, entry places and the jets of a column pin there are place matrices $B_s$ and a kernel-entry hypothesis $h$
such that the canonical entries $E^{\mathrm{can}}(P_l,O,B_s,h)$ are admissible for $O$, we obtain $B_s$ and $h$.

\emph{Fillings.} By \lref{GlobalTheory.T.latticeLayer_trueColumns_intrinsic}, which states that the cloud bound
yields a valuation map and fillings $\mathrm{Fl}(u,v,h)$ of the local group data at all served places, with the
dimensions $(t,m)$ of the datum, with weights $w_j-\kappa_J(j)$, and satisfying the lattice condition of
$\mathrm{Lay}$ for all levels, chains and submodules, we obtain these fillings.

\emph{Levels and kernel data.} By \lref{GlobalTheory.T.levels_kernelData_at_served_place}, which states that for
standard unfolded objects $O$, admissible entries and fillings as above there are, at every served place, levels of
the filling and kernel data satisfying the four conditions of kernel records, we choose levels $\mathrm{Lv}(u,v,h)$ and
kernel data $\mathrm{KD}(u,v,h)$.

\emph{Assembly.} Let $Y_0$ be the lattice instance with the valuation map, the fillings, the levels $\mathrm{Lv}$ and
trivial remaining components, and let $G_y$ be $Y_0$ completed by the comparison maps of the kernel data $\mathrm{KD}$,
and $\mathrm{ker}$ the kernel family formed by $\mathrm{KD}$. The condition on $V_{\mathrm{ram}}$ holds for such a
completion (\lref{Formalism.L2.G.Cmp.withFirst_vram}); with the weight identity and the lattice condition this
is $\mathrm{Lay}(K,w.D(Z),t,G_y)$. The kernel record $k_t$ is formed from $O$, its standardness and unfoldedness, the
canonical entries $E^{\mathrm{can}}(P_l,O,B_s,h)$ with their admissibility, and the four conditions on
$\mathrm{KD}$; by construction $k_t.E=E^{\mathrm{can}}(P_l,k_t.O,B_s,h)$.
\end{proof}

\subsubsection{Division data}\label{pf:baselat1:sssec:divproofs}

\begin{definition}[Division records]\label{pf:baselat1:def:divrec}
A \emph{division record} $r\in\mathrm{Rec}_{\mathrm{div}}(n,Y)$, for a norm-choice datum $n$ and division data $Y$,
consists of: the condition that the set of places $Y.Q_p$ is the one determined by the column input data of $w$; a
compatibility of the divisor characters of $Y$ with the character frame of $w$ at level $Z$ and the tuple tie $n.tp$
of $n$; a local tie of this compatibility at the selected places of $w$; a predicate $I_K$ on the levels of
$Y$; and a geometric tie of $Y$ relative to $I_K$.
\formal{GlobalTheory.Mem.Top.Bl.DivT}
\end{definition}

\begin{lemma}[Existence of division records]\label{pf:baselat1:divT}
Let $w$ be a column witness with a column pin, $Z$ a level of $w$ and $l\in\mathrm{Lv}_4(\baselatoneS,K,w,Z)$. There
are division data $Y\in\mathcal D(w.D(Z))$ with $\mathrm{Rec}_{\mathrm{div}}(l.n,Y)\ne\emptyset$,
$\mathrm{Mj}(l.\mathrm{what}.r.Z,Y)$, $Y.\kappa_R(\sigma)=\kappa^\sharp(w.D(Z))$ for all $\sigma$, and
$\mathrm{La}(l.\mathrm{what}.r.Z,Y)$.
\formallabel{GlobalTheory.Mem.Top.DivRec.exists_divT}
\end{lemma}

\begin{proof}
\emph{Pins.} The level data $l$ contain the condition that the places $l.W_t$ and the characters $w.I.\xi$ satisfy the
$\mu$-condition of cores; by \lref{GlobalTheory.Mem.Top.DivRec.tupPins_of_muIsW}, which states that under this
condition every tuple tie whose characters are $w.I.\xi$ is pinned, the tuple tie $l.n.tp$ is pinned.

\emph{Construction of $Y$.} Let $X_1$ be the division datum attached to the reading $l.\mathrm{what}.r.Z$ of the row
data (built from its pins and from the base change of its intersection numbers), and let $X_2$ be the \emph{reduced}
division datum obtained from $X_1$, the pinned tuple tie $l.n.tp$, the selection of $w$ and the places above the
pin. By \lref{GlobalTheory.Mem.Top.DivRec.geoOver_stages}, which states that every division datum over the
stages $w.\mathrm{st}(Z)$ admits a geometric datum over these stages for the true geometric predicate, there are a
geometric datum $X_g$, an identification $i_M$ of its multiplicity data with those of $X_2$, a predicate
$I_K$ and a geometric tie of the \emph{graft} $Y$ of $X_g$ onto $X_2$ along $i_M$ relative to $I_K$.

\emph{The record.} The reduced divisor character of the pinned tuple tie gives a compatibility of the divisor
characters of $Y$ with the character frame and $l.n.tp$, and by \lref{GlobalTheory.Mem.HD.ReQ.reLoc}, which
states that this reduced character is locally tied at the selected places, the local tie of the record holds (the
remaining local conditions---equality of the $q_v$, degree one, ramification, the off-diagonal entries, the
$\nu$-values, the type dimensions, the sets $Z_Q$ and the first local quotient values---are components of the same
local tie). The place condition $Y.Q_p$ holds by \lref{GlobalTheory.Mem.Top.DivRec.qpIs_qpSet}, which states
that the place set defined from a pinned tuple tie of the column characters is the one determined by $w$. Together
with $I_K$ and the geometric tie this is an element of $\mathrm{Rec}_{\mathrm{div}}(l.n,Y)$.

\emph{The remaining conditions.} The multiplicity data of $Y$ are those of the reading, so $\mathrm{Mj}$ holds by
\lref{GlobalTheory.Mem.R3.exists_mjHTie}, which states that division data with the multiplicity data of a reading
are tied to it; $Y.\kappa_R(\sigma)=\kappa^\sharp(w.D(Z))$ holds by construction; and since the logarithmic
$\lambda$-values of $Y$ are those of the reading, $\mathrm{La}$ holds by
\lref{GlobalTheory.Mem.Top.DivRec.lamAre_of_eq}.
\end{proof}

\begin{lemma}[Division step for the trivial object condition]\label{pf:baselat1:divprov}
Let $\baselatoneS$ be a slot family with $O_x=\mathrm{Ob}(\top)$. Then $\mathrm{Div}(\baselatoneS)$. In particular
(A1), (A2) and (A3) imply $\mathrm{Div}(\baselatoneS)$; these three hypotheses are not used.
\formallabel{GlobalTheory.Mem.Top.DivRec.divEx_prov}
\formallabel{GlobalTheory.Mem.Top.DivRec.divEx_of_ox}
\end{lemma}

\begin{proof}
By \cref{pf:baselat1:stepcalc}(iv) it suffices to prove the defining implication for fixed $K,B,f,T$ and a pinned
column witness $w$. Let $Z$ be a level and $l\in\mathrm{Lv}_4(\baselatoneS,K,w,Z)$, and let $Y$ be as in
\cref{pf:baselat1:divT}, applied with a column pin of $w$. The first four conditions of $\mathrm{Div}$ hold for $Y$.
Let $O$ be an object with $O.X=Y$ and $O.p_0=l.p_0$ and $t$ row data equal to $l.\mathrm{what}$. Then
$\mathrm{Ob}(\top)(O,t)$ is the type of pairs consisting of a proof of the true proposition and a proof of
$\mathrm{La}(t.r.Z,O.X)$; the second component is $\mathrm{La}(l.\mathrm{what}.r.Z,Y)$ transported along the
equalities $t=l.\mathrm{what}$ and $O.X=Y$. The second statement is the first one, the hypotheses being discarded.
\end{proof}

\subsubsection{Theta data}\label{pf:baselat1:sssec:thetaproofs}

\begin{definition}[Theta data tied to all data]\label{pf:baselat1:def:alltied}
Let $D$ be a global datum, $C$ a character frame for $D$, $\xi$ a map from the tuples of $D$ to Hecke characters of
$E$, and $U$ a set of tuples.
\begin{enumerate}[label=(\alph*)]
\item An \emph{infinitesimal family} $\theta$ for $D$ is a family of archimedean data satisfying the
  infinitesimal-family condition; a \emph{standard character of sign $-1$} is the additive character $\psi$ of $\A_F$ fixed in the
  development with this sign; a \emph{Weil function} $W_f$ for $\psi$ is a realisation of the Weil representation
  attached to $\psi$; it has \emph{equal rank} if its spaces have the rank prescribed by $D$.
\item \emph{Row-function data} $R_c$ for $(\theta,\psi)$ consist of row functions $\mathrm{row}(\mu,j,x)$ on the
  groups $Y_j$, for tuples $\mu$, $j\le m$ and $x$ in the multiplicity spaces $\mathrm{Mult}(\mu,j)$, of constants
  $m_Y(\mu,j)$ and Petersson constants $c_{\mathrm{pet}}(Y_j)$; $R_c$ has the characters $(C,\xi)$ and is $\mathrm{io}$-tied
  to $W_f$ if the corresponding conditions of the development hold.
\item The \emph{row-function hypothesis} $(\mathrm{RF})(C,\xi,U)$: for every infinitesimal family $\theta$, the standard
  character $\psi$ of sign $-1$ and every Weil function $W_f$ for $\psi$, there are row-function data $R_c$ with the
  characters $(C,\xi)$, $\mathrm{io}$-tied to $W_f$, such that for $\mu\in U$, $j\le m$ and $x,y\in\mathrm{Mult}(\mu,j)$ the
  function $h\mapsto\overline{\mathrm{row}(\mu,j,x)(h)}\,\mathrm{row}(\mu,j,y)(h)$ has quotient integral $\langle x,y\rangle$
  over $Y_j$ with respect to the constant $m_Y(\mu,j)\,c_{\mathrm{pet}}(Y_j)$.
\item Theta data $X\in\mathcal Th(D)$ carry automorphic row forms $\mathrm{row}^{\mathrm{aut}}(\mu,j,x)$ and Petersson products
  $\mathrm{pet}_Y(\mu,j)$. An \emph{all-tied record} $a$ for $X$ and a provision $A$ consists of a theta tie $a.T$, a Weil
  tie $a.W$, a Bessel tie, the automorphic, type and unramified Weil ties, a finiteness constant, the
  $\varphi$-upgrade, the real ties of $D$ and the provision compatibility of $A$; \emph{Weil--theta compatibility}
  is a condition on $a.W$. A \emph{theta character tie} $H$ for a norm choice $\mathrm{nch}$ and a tuple tie $T_p$ relates
  the characters of $X$ to those of $T_p$.
\end{enumerate}
\formal{Formalism.L2.Theta.AllTied}\formal{GlobalTheory.Mem.Th.RecFun}\formal{GlobalTheory.Mem.Th.WeilFun}
\end{definition}

\begin{theorem}[Theta data tied to all data]\label{pf:baselat1:alltied}
Assume (A5), (A7), (A17), (A18) and (A19). Let $D,C,\xi,U$ be as in \cref{pf:baselat1:def:alltied}, assume that $D$ is
real tied and that $(\mathrm{RF})(C,\xi,U)$ holds. Then there are theta data $X\in\mathcal Th(D)$ which are
$\mathrm{io}$-tied, such that
\[
\mathrm{pet}_Y(\mu,j)\bigl(\mathrm{row}^{\mathrm{aut}}(\mu,j,x),\mathrm{row}^{\mathrm{aut}}(\mu,j,y)\bigr)=\langle x,y\rangle\qquad
(\mu\in U,\ j\le m,\ x,y\in\mathrm{Mult}(\mu,j)),
\]
a provision $A$ and an all-tied record $a$ for $X$ and $A$ with Weil--theta compatibility, and, for every norm
choice $\mathrm{nch}$ and every tuple tie $T_p$ with $T_p.\xi=\xi$, a theta character tie $H$ whose characters agree
with those of $a.T$.
\formallabel{GlobalTheory.Mem.Th.exists_allTied_eqTied_on}
\formallabel{GlobalTheory.Mem.Th.exists_allTied_eqTied_on_of_realised}
\end{theorem}

\begin{proof}
\emph{The inputs.} The following statements hold.
\begin{itemize}
\item The types of $D$ are additive (\lref{GlobalTheory.TypeSum.typAdd_of_discreteSeries}, from (A17) and (A18)).
\item For every standard character $\psi$ of sign $-1$ there is a Weil function of equal rank
  (\lref{GlobalTheory.Mem.WeilEqRank.exists_weilFun_eqRank}, from an adelic Weil representation with its rules
  (\lref{GlobalTheory.T.adelicWeilRepresentation_with_rules_exists}, from (A7)), the Weil representation on the
  spaces of equal rank (\lref{GlobalTheory.Mem.WeilEqRank.weilOnEqEx_of_weilRepresentation}, from (A7)) and the
  twisting identity of Weil functions (\lref{GlobalTheory.Mem.WeilTwist.weilTwEq})).
\item There is a tower of spaces for $D$ (\lref{GlobalTheory.TowerSpaces.exists_tow}).
\item \emph{Realised theta data exist} (\lref{GlobalTheory.Mem.TrueRecordEq.exists_realisedE}): for every
  infinitesimal family $\theta$, standard $\psi$ of sign $-1$, Weil function $W_f$ of equal rank, tower $T_w$ and
  row-function data $R_c$ there are theta data $X$ with a theta tie $T$, a Weil tie $W$ with $W.\psi=\psi$ and a Bessel tie
  $B$ which realise $(W_f,T_w,R_c)$. This uses the integrability of Petersson products, of the top integrand and of
  Bessel functions and the lifting of automorphic forms
  (\lref{GlobalTheory.Mem.Th.petInt_of_cuspRapidDecay}, \lref{GlobalTheory.Mem.TopIntegrand.topInt},
  \lref{GlobalTheory.Mem.BesselIntegrable.besInt}, \lref{GlobalTheory.Mem.Lift.liftAut},
  \lref{GlobalTheory.Mem.Lift.liftAutEq}), all of which follow from (A5).
\item There is a provision $A$ compatible with $D$ (\lref{GlobalTheory.exists_realPlace_record_tied}, from
  (A19)).
\end{itemize}

\emph{Assembly.} Fix an infinitesimal family $\theta$ (one exists,
\lref{GlobalTheory.Member.exists_isInfFam_family}), the standard character $\psi$ of sign $-1$ (it exists and is
unique, \lref{GlobalTheory.StdAdelicChar.existsUnique_isStdChar}), a Weil function $W_f$ of equal rank for $\psi$
and a tower $T_w$. By $(\mathrm{RF})$ choose row-function data $R_c$ with the characters $(C,\xi)$, $\mathrm{io}$-tied to
$W_f$ and with the quotient-integral identities on $U$. Let $X,T,W,B$ be realised theta data for $(W_f,T_w,R_c)$, with
the realisation property $h$. The record $a$ is formed from $T$, the automorphic tie, the type tie (from the
additivity of types), $W$, the unramified Weil tie, a finiteness constant, the $\varphi$-upgrade and $B$ (all
components of $h$), the real ties of $D$ and the compatibility of $A$. The $\mathrm{io}$-tie of $X$ is obtained from that
of $R_c$, and Weil--theta compatibility is a component of $h$.

\emph{Petersson products.} Let $\mu\in U$, $j\le m$ and $x,y\in\mathrm{Mult}(\mu,j)$. By the Petersson formula of the
theta tie $T$, $\mathrm{pet}_Y(\mu,j)(\mathrm{row}^{\mathrm{aut}}(x),\mathrm{row}^{\mathrm{aut}}(y))$ is the quotient integral over $Y_j$ of
$\overline{\mathrm{row}^{\mathrm{aut}}(x)}\,\mathrm{row}^{\mathrm{aut}}(y)$ with respect to $m_Y\,c_{\mathrm{pet}}$. The realisation property
identifies $\mathrm{row}^{\mathrm{aut}}$ with the row functions of $R_c$, the constant $m_Y$ with that of $R_c$ and the Petersson
constants of $T$ with those of $R_c$; the quotient-integral identity of $R_c$ then gives $\langle x,y\rangle$.

\emph{Characters.} For a norm choice $\mathrm{nch}$ and a tuple tie $T_p$ with $T_p.\xi=\xi$, the realisation property
shows that $T$ reads the characters $(C,\xi)$, and \lref{GlobalTheory.TowerCharRead.exists_thetaChrTied}, which
states that a theta tie reading the characters of a tuple tie yields a theta character tie with the same characters,
gives $H$.
\end{proof}

\subsubsection{Spectral objects over replaced stages}\label{pf:baselat1:sssec:satproofs}

We use the notation of \cref{pf:baselat1:def:satnew}(a).

\begin{definition}[Row packages and new stages]\label{pf:baselat1:def:rowpack}
A \emph{row package} for a datum $D$ consists of a finite group
$G'$, index types, a family $A$ of inputs, a number field $K_s$, real data $R$ (with their models and Hilbert spaces), a
face $F$ and a linear map $e_E\colon E_E\to R.H_E$ from the row space $E_E$ of $D$. The \emph{row hypothesis} on it is the
conjunction of: the level character of $R$ is $A.\mu_0$; $R$ has group levels at $F$ and a group pin; the models of $R$
hold; the image of $e_E$ is the preimage of the level part of $R$ at $F$, and $\|e_E(x)\|^2=n_P(F)\|x\|^2$; and the
dimension tie at $F$. For a spectral package $P\in\mathrm{SPk}(D,F_0)$ with Hida setting $P.X$ and good operator $P.G$,
let $E_C(P.X,P.G)$ be its space of conjugated good rows and $\mathrm{HH}(P.X)$ its space of spectral forms; $D$ prescribes
the dimension $D_{\mathrm{all}}$. A pair $(h_D,j)$ of the equality $\dim\mathrm{HH}(P.X)=D_{\mathrm{all}}$ and a linear
isometry $j\colon E_E\to E_C(P.X,P.G)$ defines a new $S$-stage $P.\mathrm{new}(h_D,j)$, and for a division update
$(N,h)$ the \emph{replaced stages} $(P.\mathrm{new}(h_D,j),N,h)$.
\formal{GlobalTheory.Mem.SRe.SNew}\formal{GlobalTheory.Mem.SRe.SAtL}\formal{GlobalTheory.Mem.SRe.SRest}%
\formal{GlobalTheory.Mem.SRe.SeeAt}\formal{GlobalTheory.Mem.SRe.RowJ}
\end{definition}

\begin{proposition}[Spectral-level objects]\label{pf:baselat1:speclevel}
Assume (A5) and (A9)--(A12). Let $D$ be a datum with a row package satisfying the row hypothesis, $F_0$ forms and
$P\in\mathrm{SPk}(D,F_0)$ with integrable spectral function, and assume $\mathrm{Gap}(D)$, $\mathrm{LT}(D)$,
$\mathrm{RTd}(D)$, $(\mathrm{IR})(D)$ and that $D^*$ has the standard value. Then there are $h_D$, an isometry
$j\colon E_E\to E_C(P.X,P.G)$, a finite adelic point $t$ and $c>0$ such that the row level is seen at $t$, the
evaluation tie with constant $c$ holds at $t$, the row isometry condition holds for $j$ at $t$, and for every
division update $(N,h)$ and all forms $F'$ on the replaced stages $(P.\mathrm{new}(h_D,j),N,h)$ with the same
degree-zero levels and values as $F_0$, the redivided full datum is finite and there are restricted spectral data
$S$ over the replaced stages with $S.j$ equal to $j$ followed by the canonical identification of conjugated good
rows.
\formallabel{GlobalTheory.T.exists_spectralLevel_objects_replacedStages}
\formallabel{GlobalTheory.Mem.SRe.exists_sRest_top}
\end{proposition}

\begin{proof}
\emph{The point $t$.} By \lref{GlobalTheory.T.exists_finiteAdelicPoint_carries_rowLevel}, which states that if
the level character of $R$ is $A.\mu_0$, $R$ has group levels and a group pin, and $(\mathrm{IR})(D)$ holds, then for
every Hida setting with the instance of $D$ there is a finite adelic point at which the row level of $(R,F)$ is seen,
we choose such a point $t$ for $P.X$.

\emph{Evaluation.} By \lref{GlobalTheory.T.lowestVector_evaluation_identifies_levelRows_conj} (from (A10),
(A5), (A11), (A12), the gap hypothesis, the level character and the models of $R$ at the archimedean places and at
$\nu_0$), there are a family $v_0$ of unit lowest vectors and $c>0$ such that evaluation at $t$ against $v_0$, scaled
by $c$, maps the multiplicity space of $(R,F)$ to the conjugated rows of $P.X$. By
\lref{GlobalTheory.T.evaluation_goodRows_to_conjugate_goodRows} the same evaluation maps the good rows $E_P(R,F)$
to the conjugated good rows $E_C(P.X,P.G)$, and the evaluation tie with constant $c$ holds at $t$. By
\lref{GlobalTheory.T.levelColumnsGoodRows_le_principal_and_finiteDimensional} (from (A9) and (A5)), the level
part of $R$ at $F$ is contained in its trivial part and the multiplicity space is finite dimensional.

\emph{The isometry.} By \lref{GlobalTheory.T.exists_rowIsometry_conjugatedGoodRows_finrank}, which states that
under the conditions on $e_E$ in the row hypothesis and the two evaluation statements, the dimension tie implies
$\dim\mathrm{HH}(P.X)=D_{\mathrm{all}}$ and there is a linear isometry $j\colon E_E\to E_C(P.X,P.G)$ satisfying the row
isometry condition at $t$, we obtain $h_D$ and $j$.

\emph{Replaced stages.} Let $(N,h)$ and $F'$ be as in the statement. The redivided full datum is finite by
\lref{GlobalTheory.T.finX_spectralRedivisionFullChoices}, which states that finiteness of a full datum is
preserved under redivision, and restricted spectral data $S$ with the stated $S.j$ exist by
\lref{GlobalTheory.Mem.SRe.exists_sRest_re}, which uses the integrability of the spectral function,
$\mathrm{LT}(D)$, $\mathrm{RTd}(D)$ and the standard value of $D^*$.
\end{proof}

\begin{theorem}[Spectral saturation]\label{pf:baselat1:satnewthm}
(A5) and (A9)--(A12) imply the spectral saturation statement $(\mathrm{SN})$ of \cref{pf:baselat1:def:satnew}.
\formallabel{GlobalTheory.Mem.SRe.UpS.satOfDiscreteSeries}
\formallabel{GlobalTheory.T.exists_spectralObjects_over_replacedStages}
\formallabel{GlobalTheory.Mem.SRe.satOfDiscreteSeries}
\end{theorem}

\begin{proof}
Let the data and hypotheses of $(\mathrm{SN})$ be given, with $D=a.D$. Apply \cref{pf:baselat1:speclevel} and let
$h_D,j,t,c$ be as there; the new $S$-stage is $Z_s=P.\mathrm{new}(h_D,j)$. Let $(N,h)$ be a division update and $F_1$
forms on the replaced stages which agree with $F_0$. \Cref{pf:baselat1:speclevel} gives the finiteness of the
redivided full datum and restricted spectral data $S$ with $S.j$ equal to $j$ followed by the canonical
identification. By \lref{GlobalTheory.Mem.SRe.nonempty_sAtL}, which assembles the spectral data at the
replaced stages from $P$, the agreement of the forms, this finiteness, $R$, $F$, $e_E$, the point $t$ with the constant
$c$, the seeing condition, the evaluation tie, $j$ with the row isometry condition, and such $S$, the spectral data at
the replaced stages exist. The three formal statements listed are the same statement (two of them are restatements
of the third).
\end{proof}

\subsubsection{Hida settings and spectral data with gap}\label{pf:baselat1:sssec:hidaproofs}

Let $D$ be a global datum with host $(F,E)$, hermitian space $V$ over $E$ with form $h$ and lattice $L_0$, ranks $r$, $r_S$ and
$m$, words $u$ and forms $F_0(u)$ of degree zero. Let $I=\mathrm{inst}(D)$ be the automorphic instance of the unitary
group of $(V,h)$ attached to $D$, with adelic group $G_{\A}$, finite part $G_f$, archimedean part $G_\infty$, finite
lattice $L_f$ and the constant $g_D=g_0g_1$.

\begin{definition}[Hida settings]\label{pf:baselat1:def:hida}
A \emph{Hida setting} $X$ for $D$ is a record assembled from the following blocks.
\begin{enumerate}[label=(\alph*)]
\item An instance $A$ of the unitary group: an automorphic instance with an identification $\iota_E$ of its field with
  $E$ (compatible with conjugation and with $F$), a basis of $V$ in which its matrix $J$ is the Gram matrix of $h$ and
  its lattice is $L_0$, and identifications of its archimedean places and signatures with those of $D$.
\item Lattice data $L$: vectors $e_a,f_a$ ($a<r_S$) of $V$ with $h(e_a,e_b)=h(f_a,f_b)=0$ and $h(e_a,f_b)=\delta_{ab}$, a
  lattice $L_b$ orthogonal to them such that $L_0$ is the sum of $L_b$ and of the spans of suitable multiples of the
  $e_a$ and $f_a$, vectors $z_u\in L_b$ orthogonal to the $e_a,f_a$ with $h(z_u,z_u)=-T(u)$, and flags $F_u$ with
  $e$-part $(e_a)$ and $z$-vector $z_u$ satisfying the fixed-point hypothesis.
\item Additive character data $P$: a character $\psi_{E,f}$ of $E\otimes\A_{F,f}$, trivial on the integral
  elements, with $\psi_{E,f}(x)\psi_{E,\infty}(x)=1$ for $x\in E$.
\item A reference frame $R_f$: adelic vectors $e_a^0,f_a^0,z^0$, weight vectors $b_w,b'_w$ and an idempotent $e_w$ in
  $L_f$ approximating the global vectors modulo $g_D$, with the corresponding congruences for their hermitian
  products, spanning $L_f$ modulo $g_D$.
\item Transport data $X_u$: elements $x_u$ of the level group in $G_f$ fixing the $e_a,f_a$, moving $z_u$ to
  $z^0$ modulo $g_D$, and satisfying the one-off condition relative to $g_D$.
\item An automorphic bundle $B_u$: a Haar measure $\mu_f$ on $G_f$, a fundamental domain, discrete series
  representations $\sigma_v$ with lowest unit vectors at the archimedean places and their tensor product, the right
  regular representation, an admissible idempotent built from $\psi_{E,f}$ and the reference frame, the
  corresponding convolution operator, and the integrability of the archimedean spectral function $X.\varphi$
  (built from the matrix coefficients of the $\sigma_v$ and their formal degrees) with respect to the measure
  $\mu_\infty$.
\item Representation data $\Pi$: for every word $u$ finite-dimensional representations of $G_\infty$ which are
  tensor products of local representations of Harish-Chandra type prescribed by $D$, with orthonormal frames.
\item Measure blocks $M_u$: Haar measures $\mu_N(u),\mu_H(u)$ on the unipotent radical and on the stabiliser of the
  frame of $u$, and the quotient measure $\nu_Q(u)$.
\item A function block $W_b$: for every word $u$ a continuous function $W_u$ on $G_\infty$ which is a tensor product of
  local functions, equivariant under the parabolic subgroup, reproduced by convolution with $X.\varphi$, and
  square-integrable of norm $1$ for $\nu_Q(u)$.
\item A measure block $M_b$: measure data $\mathrm{meas}_G$ and $\mathrm{meas}_W(u)$ (global, archimedean, finite and
  local Haar measures for the group, the frames and their stabilisers) in textbook normalisation.
\end{enumerate}
The following conditions on $X$ appear below. An \emph{adelic bridge} is an isomorphism of topological groups from
$G_{\A}$ onto the adelic points of the unitary group with the bridge property. The \emph{measure tie} for a bridge
states that the measures of $\mathrm{meas}_G$ and the archimedean measures of $\mathrm{meas}_W(u)$ are the images of
the corresponding measures of $X$ under the bridge. A character $\psi$ of $\A_F$ is \emph{admissible} if it is
continuous, unitary, trivial on $F$ and nontrivial, and \emph{standard} if
$\psi_\infty(x)=\exp(2\pi i\sum_w\Re x_w)$. \emph{Measure compatibility} $\mathrm{MC}(X,\psi)$ states that for every
word $u$ the stabiliser measure of the frame of $u$, the group measure and a $\psi$-measure on the corresponding
quadric are compatible and the unipotent radical has covolume one, and that the analogous compatibility holds for all
tuples of $2r_S+2$ vectors with nondegenerate Gram matrix; \emph{local compatibility} $\mathrm{LC}(X,\psi,u)$ states
that at every finite place the local measures of $\mathrm{meas}_G$ and $\mathrm{meas}_W(u)$ are tied with respect to
$\psi$. The \emph{level subgroup character} condition states that the level set $K_2$ of $X$ is a group and the
character $\chi_2$ is a unitary character on it.
\formal{Formalism.L2.SG.HidaSetting.HidaSetting}\formal{Formalism.L2.SG.SplitPrimeNormLemmas.mk}%
\formal{Formalism.L2.SG.AdBridge.AdBridgeIso}\formal{Formalism.L2.SG.HidaSetting.HidaSetting.MeasTie}%
\formal{Formalism.L2.SG.MassDataConc.MeasCompatψ}\formal{Formalism.L2.SG.LocVol.LocCompatX}%
\formal{Formalism.L2.SG.HidaSetting.HidaSetting.LevelSubgroupCharacter}
\end{definition}

\begin{definition}[Transports, full data and the gap hypothesis]\label{pf:baselat1:def:transport}
Let $X$ be a Hida setting for $D$.
\begin{enumerate}[label=(\alph*)]
\item A \emph{transport} of $F_0(u)$ into $X$ consists of an isometry $w_t$ of $(V,h)$ with $w_t(e_a)$, $w_t(f_a)$ the
  hyperbolic basis of $D$ and $w_t(z_u)=z_m(u)$; a ring isomorphism $\iota_A$ of finite ad\`eles compatible with
  $\iota_E$ and with the integral elements; an element $h_0$ of the finite-adelic unitary group of the Gram matrix of
  degree zero; for every class $c$ of degree zero an element $h_c$ of $G_f$ representing $y(c)h_0^{-1}$; and a map $J$
  from the representation space of $u$ to functions on the degree-zero frame which is equivariant, has values of the
  prescribed type and carries the Petersson product of $F_0(u)$ to the inner product. It \emph{matches levels},
  $\mathrm{LM}$, if $h_0^{-1}$ carries the degree-zero lattice of $F_0(u)$ to the preimage of $L_f$, and conjugation by $h_0$
  carries the level subgroup $K_u$ of $X$ onto the level group of $F_0(u)$, matching the characters.
\item The \emph{column dictionary} of a transport is a linear map $\mathrm{cd}$ from the column space of $X$ at $u$ to the
  functions on the degree-zero frame; the conditions \emph{into}, \emph{bijective} and \emph{isometric} state that its
  image lies in the space $\mathrm{Mu}_0(u)$ of $D$, that it is a bijection onto $\mathrm{Mu}_0(u)$, and that it carries the
  inner product of column data to the degree-zero Petersson product.
\item A \emph{full datum} $Y$ for $X$ consists of the level subgroup character condition, a typed reduction $T_R$ with
  $T_R.\psi_0(v)\ne1$ for all $v$, for every word $u$ a finite set of row representatives (split rows) with a column
  system admitting column representatives, sets of representatives $D_t(u)\subseteq G_{\A}$ and fundamental sets
  $F_N(u)\subseteq G_\infty$, and local elements $\delta_p$ with models $M_p$ at all places $p$ of $F$ which agree, with
  $\delta_p=g_0g_1/g_i$ when $g_i\in p$. It is \emph{finite}, $\mathrm{Fin}(Y)$, if the sharp maps of its column systems
  are bijective, the spectral spaces satisfy the inclusion and linearity conditions, the space $\mathrm{HH}(X)$
  is finite dimensional and all fibres of the row transitions are finite.
\item The \emph{gap hypothesis} $\mathrm{Gap}(D)$: $r\ge1801$ and $b_{j+1}-b_j\ge16r^{19}$ for the Harish-Chandra
  parameters $b$ of $D$ at all archimedean places $\tau$ and all $j+1<m$. The \emph{local ties} $\mathrm{LT}(D)$: at every
  place $p$ of $F$ the local datum of $D$ is isomorphic to the canonical local datum.
\end{enumerate}
\formal{Formalism.L2.SG.ColDict.Transport}\formal{Formalism.L2.SG.ColDict.LevelMatch}%
\formal{Formalism.L2.SG.ColDict.ColDictInto}\formal{Formalism.L2.SG.ColDict.ColDictBij}%
\formal{Formalism.L2.SG.ColDict.ColDictIsom}\formal{Formalism.L2.SG.Full.Full}\formal{Formalism.L2.SG.PackageTie2.FinX}%
\formal{Formalism.L2.SG.ConcreteSettingExists.GapHyp}\formal{Formalism.L2.SG.Pf.LocCP.LocTied}
\end{definition}

\begin{lemma}[Transports]\label{pf:baselat1:transport}
Let $X$ be a Hida setting for $D$ satisfying the level subgroup character condition, and assume that there is a finite
adele isomorphism $\iota_A$ compatible with $\iota_E$ and the integral elements. Then for every word $u$ there is a
transport of $F_0(u)$ into $X$ which matches levels.
\formallabel{GlobalTheory.St.SpectralData.exists_transport}
\end{lemma}

\begin{proof}
By \lref{GlobalTheory.St.WittT.wittExists} there is an isometry $w_t$ of $(V,h)$ with the required values on
$e_a,f_a,z_u$. By
\lref{GlobalTheory.St.G.genusMatch_holds} there is $h_0$ for which the lattice condition of $\mathrm{LM}$ holds,
and by \lref{GlobalTheory.St.G.levelConj_holds} there is $h_0'$ for which both the lattice condition and the
level condition of $\mathrm{LM}$ hold (the level subgroup of $X$ and the level group of $F_0(u)$ are conjugate by
$h_0'$, with matching characters). By \lref{GlobalTheory.St.G.hcExists_all} there are representatives $h_c$ for
$h_0'$, and by \lref{GlobalTheory.jExists_all} there is a map $J$ with the required properties. These data
form a transport which matches levels.
\end{proof}

\begin{lemma}[Rescaling of measures]\label{pf:baselat1:rescale}
Assume (A13). Let $X$ be a Hida setting for $D$ with an adelic bridge $h_b$ and the measure tie. There are measure
data $\mathrm{meas}'_G$, $\mathrm{meas}'_W(u)$ in textbook normalisation, with the same archimedean components as those
of $X$, and an admissible standard character $\psi$ of $\A_F$, such that the setting $X'$ obtained from $X$ by
replacing its measure data satisfies the measure tie for $h_b$, $\mathrm{MC}(X',\psi)$ and $\mathrm{LC}(X',\psi,u)$
for all words $u$.
\formallabel{GlobalTheory.St.Rescale.rescaleStep3}
\end{lemma}

\begin{proof}
By \lref{GlobalTheory.CmHermitianRankFrames.cm_herm_rank_frames} the instance of $X$ is a CM extension,
$J$ is hermitian and nondegenerate, $4r_S+5$ is at most the rank, and for every word $u$ the frame of $X$ at $u$ is
standard for $(J,T_u)$ for some $T_u\in F$. By \lref{GlobalTheory.St.siegelWeil_volumes} the convergent
Siegel--Weil formula (A13) gives, for every CM extension and hermitian nondegenerate matrix, the two volume statements
for adelic unitary groups and their stabilisers which are used below. By
\lref{Formalism.L2.SG.Pf.standard_character_exists} there is an admissible standard character $\psi$.

\emph{The group measure data.} For every tuple $X_0$ of $2r_S+2$ vectors with nondegenerate Gram matrix,
\lref{GlobalTheory.T.exists_adelicTorusMeasure_stabiliserHaar}, which states that for an admissible $\psi$ and the
global Haar measure $\mu_G$ there are a Haar measure $\mu'$ on the stabiliser of $X_0$ and a $\psi$-measure $\nu$ on the
corresponding variety such that $(\mu_G,\mu',\nu)$ are compatible, gives $\mu'_B(X_0)$; for degenerate $X_0$ we keep the
measure of $X$. Let $\mathrm{meas}'_G$ be $\mathrm{meas}_G$ with these stabiliser measures; it is in textbook
normalisation.

\emph{The frame measure data.} Fix a word $u$. The Gram matrix of the frame of $u$ is nondegenerate and hermitian
(\lref{GlobalTheory.FrameGram.std_frame_gram}). As before there are a $\psi$-measure $\nu_H$ and a Haar measure $\mu_H$
on the stabiliser of the frame, compatible with $\mu_G$; there is a Haar measure $\mu_N$ on the unipotent radical of
covolume one (\lref{GlobalTheory.T.exists_haarMeasure_unipotentRadical_covolume_one}); and at every finite place
$v$ there is a Haar measure $\mu_{S,v}$ on the local stabiliser which is tied to the local group measure for $\psi$
(\lref{GlobalTheory.T.exists_frameStabilizerHaar_normalised_fibreMeasure}). By
\lref{GlobalTheory.St.MeasWord.exists_measW_prescribed} there is measure data $\mathrm{meas}'_W(u)$ in textbook
normalisation with the archimedean components of $\mathrm{meas}_W(u)$ and with the components $\mu_N$, $\mu_H$ and
$\mu_{S,v}$.

\emph{Verification.} The measure tie for $X'$ follows from that of $X$, because the archimedean components are
unchanged. $\mathrm{MC}(X',\psi)$ consists of the compatibilities of the frames of all words, the covolume one of
$\mu_N$, and the compatibilities for all nondegenerate tuples. $\mathrm{LC}(X',\psi,u)$ consists of the local ties of
the $\mu_{S,v}$.
\end{proof}

\begin{theorem}[Existence of Hida settings]\label{pf:baselat1:setting}
Assume (A13). Let $D$ be a datum with $\mathrm{Gap}(D)$ and $F_0$ forms of degree zero for all words. There are a Hida
setting $X$ for $D$ whose instance is $\mathrm{inst}(D)$ and whose identification $\iota_E$ is the identity, an adelic bridge
$h_b$ and an admissible standard character $\psi$ such that $X$ satisfies the measure tie for $h_b$, $\mathrm{MC}(X,\psi)$
and $\mathrm{LC}(X,\psi,u)$ for all $u$, and such that, if $X$ satisfies the level subgroup character condition, for every word $u$
there is a transport of $F_0(u)$ into $X$ which matches levels.
\formallabel{GlobalTheory.St.SpectralData.exists_setting}
\end{theorem}

\begin{proof}
Two archimedean statements are used: the nonvanishing of the Bessel integrals of the relevant discrete series
(\lref{Formalism.L2.SG.Closed.BesselNVGap_holds}) and the weighted integrability of their matrix coefficients
(\lref{Formalism.L2.SG.Closed.CoefWeightedLoc_holds}). We construct the blocks of
\cref{pf:baselat1:def:hida} one after the other.
\begin{enumerate}[label=(\arabic*)]
\item $A$ is the unitary group instance attached to $D$.
\item Lattice data $L$ exist by \lref{GlobalTheory.St.Lat.latExists_of_forms0}, applied to the forms $F_0(u_0)$ of
  the base word.
\item Additive character data $P$ exist for every instance (\lref{Formalism.L2.SG.Closed.PsiExists_holds}).
\item A reference frame $R_f$ exists for every $(A,L)$ (\lref{Formalism.L2.SG.Closed.RefExists_holds}).
\item Transport data $X_u$ exist for every $(A,L,R_f)$ (\lref{Formalism.L2.SG.Closed.XuExists_holds}).
\item An automorphic bundle $B_u$ exists for every $(A,L,P,R_f)$, given the weighted integrability of matrix
  coefficients and $\mathrm{Gap}(D)$ (\lref{Formalism.L2.SG.Closed.BunExists_holds}).
\item Representation data $\Pi$ and measure blocks $M_u$ exist for every $(A,L)$
  (\lref{Formalism.L2.SG.Closed.PiExists_holds}, \lref{Formalism.L2.SG.Closed.MuExists_holds}).
\item A function block $W_b$ exists, given the Bessel nonvanishing, the weighted integrability and $\mathrm{Gap}(D)$
  (\lref{GlobalTheory.WBlockExists.realPlaceFunctionBlockExists_of_bessel_coef}).
\item There are a measure block $M_b$ and an adelic bridge $h_b$ for which the setting assembled from these blocks
  satisfies the measure tie (\lref{GlobalTheory.St.mExists_holds}).
\end{enumerate}
Let $X_1$ be the setting assembled from $A,L,P,R_f,X_u,B_u,\Pi,M_u,W_b,M_b$. By \cref{pf:baselat1:rescale} there are new
measure data and an admissible standard character $\psi$; let $X$ be $X_1$ with the new measure data. Then $X$ has the
measure tie for $h_b$, $\mathrm{MC}(X,\psi)$ and $\mathrm{LC}(X,\psi,u)$; its instance is $\mathrm{inst}(D)$ and $\iota_E$ is the
identity, by construction. Finally, if $X$ satisfies the level subgroup character condition and $u$ is a word,
\cref{pf:baselat1:transport}, applied with $\iota_A$ the identity (which is compatible with $\iota_E$ and the integral
elements), gives a transport which matches levels.
\end{proof}

\begin{theorem}[Full data]\label{pf:baselat1:fulldata}
Let $X$ be a Hida setting for $D$ and assume $\mathrm{LT}(D)$. Then there is a full datum for $X$.
\formallabel{GlobalTheory.fullExists_of_locTied}
\formallabel{Formalism.L2.SG.Pf.fullExists_tied}
\end{theorem}

\begin{proof}
The formal proof derives the statement from three general statements, which are proved in the development, and
from the level subgroup character condition, which holds for every Hida setting
(\lref{Formalism.L2.SG.Closed.level_subgroup_is_group}).
\begin{itemize}
\item Typed reductions: for every Hida setting with the level subgroup character condition there is a typed
  reduction $T_R$ with $T_R.\psi_0(v)\ne1$ for all $v$ (\lref{Formalism.L2.SG.Pf.trExistsC_holds}).
\item Lattice rows: for every word $u$ the lattice-row, surjectivity and injectivity conditions and the finiteness of
  classes hold (\lref{Formalism.L2.SG.Pf.lattice_steps_hold}). Together with the inclusion conditions which hold
  for every Hida setting (\lref{Formalism.L2.SG.Pf.latSub}, \lref{Formalism.L2.SG.Pf.inclInj_holds},
  \lref{Formalism.L2.SG.Pf.gamIsom_of}) they give the row steps of all words
  (\lref{Formalism.L2.SG.Pf.rowStepsC_of}, \lref{Formalism.L2.SG.Pf.latRow_of},
  \lref{Formalism.L2.SG.Pf.classFin_of}).
\item Local dual agreement: under $\mathrm{LT}(D)$ there are local elements $\delta_p$ and models $M_p$ at all places which
  agree and satisfy $\delta_p=g_0g_1/g_i$ when $g_i\in p$ (\lref{Formalism.L2.SG.Closed.LocCTied_holds}).
\end{itemize}
For every word $u$, the row steps give a finite set of split row representatives with a column system admitting column
representatives (\lref{Formalism.L2.SG.Pf.rowsCols_of}); there are sets of representatives $D_t(u)$
(\lref{Formalism.L2.SG.Pf.repsExist}) and fundamental sets $F_N(u)$ (\lref{Formalism.L2.SG.Pf.fundNC_holds}).
These data form a full datum.
\end{proof}

\begin{theorem}[Spectral data with gap]\label{pf:baselat1:specgap}
Assume (A12) and (A13). Let $D$ be a datum with $\mathrm{Gap}(D)$ and $\mathrm{LT}(D)$, and $F_0$ forms of degree zero.
There are a Hida setting $X$ for $D$ with instance $\mathrm{inst}(D)$ and $\iota_E$ the identity, an adelic bridge $h_b$, an
admissible standard character $\psi$, a finite full datum $Y$ for $X$, and transports $\mathrm{tr}(u)$ of $F_0(u)$ into $X$,
such that $X$ has the measure tie, $\mathrm{MC}(X,\psi)$ and $\mathrm{LC}(X,\psi,u)$ for all $u$, the spectral function
$X.\varphi$ is integrable for $\mu_\infty$, and for every word $u$ the transport $\mathrm{tr}(u)$ matches levels and
its column dictionary is into, bijective and isometric for the column system of $Y$.
\formallabel{UpperTheory.Lit.exists_spectralData_gap_alt}
\formallabel{GlobalTheory.T.local_bessel_uniqueness_holds}
\end{theorem}

\begin{proof}
Let $X,h_b,\psi$ be as in \cref{pf:baselat1:setting} and $Y$ a full datum for $X$ (\cref{pf:baselat1:fulldata}). By
\lref{CompanionTheory.Lime.BM.wuDecayOfGap_holds'}, which states that under $\mathrm{Gap}(D)$ every Hida setting
for $D$ has Wu decay (the reproducing and weighted integrability properties of the functions $W_u$) and an integrable
spectral function, $X.\varphi$ is integrable and $X$ has Wu decay. By
\lref{GlobalTheory.FinitenessBundle.finX_of_decay}, which states that (A12), Wu decay and the integrability of
$X.\varphi$ imply finiteness of every full datum, $Y$ is finite. Since $Y$ contains the level subgroup character condition, we may choose for every word $u$ a
transport $\mathrm{tr}(u)$ which matches levels (\cref{pf:baselat1:setting}). Since the stabiliser of the column
level is trivial (\lref{GlobalTheory.columnLevel_stabTriv}), \lref{GlobalTheory.St.colDictCore_holds},
which states that for a transport matching levels and a trivial column stabiliser the column dictionary is into and
bijective, and isometric for every column system with column representatives whose sharp map is injective, applies
to $\mathrm{tr}(u)$ and the column system of $Y$ (whose sharp map is injective by finiteness). The second formal
statement is identical to the first.
\end{proof}

\begin{remark}\label{pf:baselat1:rem:bessel}
The second formal statement listed in \cref{pf:baselat1:specgap} has a name referring to local Bessel uniqueness,
but it is the existence statement of \cref{pf:baselat1:specgap}; no Bessel model and no uniqueness statement occurs in
it. In the proof of \cref{pf:baselat1:specgap}, Bessel integrals enter only through the archimedean nonvanishing
statement used in the proof of \cref{pf:baselat1:setting}.
\end{remark}

\subsubsection{Inputs of the twist theorem}\label{pf:baselat1:sssec:inputproofs}

\begin{theorem}[Inputs of the twist theorem]\label{pf:baselat1:inputs}
Assume (E14), (A20) and (E15). Let $D$ be a global datum, $X_r$ a real side of $D$, $E$ twisted-form data for $D$, $Y_s$ a
series side over $X_r$ for $E$, and $G\in\mathcal G(D)$. Assume the rationality $\mathrm{RR}$ of the rows of $Y_s$ at all
nonzero index places, the boundedness $\mathrm{XB}$ of the $\xi$-values of the readings of $Y_s$, the Hecke sup-lattice
condition $\mathrm{HS}$ of $X_r$ at all $j$, the Hodge condition $\mathrm{Hdg}(D,G.\mathrm{wd})$, and that the field
$\mathit{LL}_0$ contains every complex root of unity of order $g_k(g_k-1)$, $k=0,1$. Then there are $N_0$ and an input
record $I\in\mathcal I(D,E,N_0)$, i.e.\ all the statements listed as hypotheses in \cref{pf:baselat1:twistthm} together
with $\mathrm{SH}(i)$ at every nonzero index place $i$.
\formallabel{GlobalTheory.T.exists_inputs_of_galHodgeN_roots}
\formallabel{GlobalTheory.TwInputs.exists_inputs_of_hkSup_galHodge}
\end{theorem}

\begin{proof}
\emph{Galois inputs.} By \lref{GlobalTheory.TwGalInputs.galois_inputs_alt}, which derives from (E14), (A20), (E15)
and $\mathrm{Hdg}(D,G.\mathrm{wd})$ the admissibility statements $\mathrm{adm}_1$, $\mathrm{adm}_2$, $\mathrm{adm}_8$, goodness
of all signatures, the signature condition, a function $N_0$ with $\mathrm{adm}_9(N_0)$, $\mathrm{adm}_{10}$ and
$\mathrm{SH}(i)$ at every nonzero index place, we obtain these components and $N_0$.

\emph{Laws of the series side.} Let $(\rho,\mathrm{good},\mathrm{rd})$ be the setup, its goodness and its readings,
which are components of $Y_s$. By \lref{GlobalTheory.TwInputs.laws_of_serSide}, $\mathrm{RR}$ implies the laws of
the readings (among them their faithfulness and the reading of the action). The readings have $\xi$-type
(\lref{GlobalTheory.XiTypeProof.xiType}, from faithfulness, the action reading and the level of the setup,
\lref{GlobalTheory.TwInputs.levelOf_of_serSide}). By \lref{GlobalTheory.TwInputs.supLaws_of_sides},
$\mathrm{HS}$ implies the sup-lattice conditions for the group, the full one and the one for the Hecke operators,
and the commutation of the group action and of the Hecke operators.

\emph{The record.} The universal character condition and the range condition on the Hecke indices are components of
$Y_s$. The recursion $\mathrm{UR}$ of the universal twisted series follows from the full tie of $Y_s$ and the laws
(\lref{GlobalTheory.T.universalTwistedSeriesDistribution_of_compatibility}); the specialisation of twisted series
likewise (\lref{GlobalTheory.SeriesTie.thetaLatticeCondition_of_fullSeriesComparison}); $\mathrm{adm}_{T4}(E)$
follows from the full tie and the roots of unity in $\mathit{LL}_0$ (\lref{GlobalTheory.T.admT4_of_roots_in_field});
the type-keeping condition for $E$ from the full tie and the action reading
(\lref{GlobalTheory.ExportType.typeKept_of_fullTie}); the universal sup-lattice condition from the full tie, the
laws, the $\xi$-type and $\mathrm{XB}$ (\lref{GlobalTheory.TwInputs.univSup_of_xi}); and the semisimplicity of the
Hecke action from the two commutation statements and the full sup-lattice condition
(\lref{GlobalTheory.TwHkSS.hkSS_of_sides}). Together with the Galois inputs these are the fields of $I$.
\end{proof}

\subsubsection{The coefficient bound for positive exponent}\label{pf:baselat1:sssec:boundproofs}

Fix $i\in X.\mathrm{Tw}$ with $\ell=\ell_i$, $\nu=\nu(i)$, let $\baselatoneO_0$ be the coefficient ring at $i$ with maximal
ideal $\mathfrak m$, and put $\mathfrak a=\ell^{\nu}\mathfrak m$ and $e=v_{\baselatoneO_0}(\ell^\nu)$. Let $T_s$ be the
representation attached to a signature $s$, $\rho^\vee_s$ its dual (\lref{GlobalTheory.NewPrime.exists_isDual}), and
$c$ the complex conjugation acting on $G_E$. We use the notation of \cref{pf:baselat1:def:aux,pf:baselat1:def:galside};
families of places are indexed by $\iota_{R,n}=\{1,\dots,R\}\sqcup\{1,\dots,n\}$, and we write $(1,a)$ and $(2,a)$ for the
two kinds of indices.

\begin{definition}[Induction data]\label{pf:baselat1:def:induction}
Let $s$ be a signature, $M\ge1$, $d$ an assignment of characters $d_P\colon G_E\to\Z/\ell^M$ to places, $A$ a finite set of
places, $L\le G_E$ a subgroup and $z\in G_E$.
\begin{enumerate}[label=(\alph*)]
\item $L$ is \emph{small} for $(s,z)$ if it is open, normal, stable under conjugation by $c$, contains all commutators
  $gzg^{-1}z^{-1}$, and $T_s(g)$ and the cyclotomic unit of $g$ are $\equiv1$ modulo $\mathfrak m$ for $g\in L$. It is
  \emph{$2$-small} if moreover $L$ fixes the roots of unity of order $2\ell^M$ and $T_s(g)\equiv1$ modulo $\mathfrak a$ for
  $g\in L$.
\item A family $P\colon\iota_{R,n}\to\{\text{places}\}$ is \emph{localised}, $\mathrm{PL}(M,d,A,L,z,P)$, if
  $\mathrm{Aux}(i,M,A,P,d\circ P)$; for every $a\le R$ the inertia groups at the places over $P_{(1,a)}$ lie in $L$ and every
  Frobenius element at $P_{(1,a)}$ is conjugate to $z$ modulo $L$; and for every $a\le n$ the decomposition group at
  $P_{(2,a)}$ lies in $L$ and $d_{P_{(1,b)}}(\mathrm{Frob}_{P_{(2,a)}})=\delta_{ab}$ for all $b\le R$.
\item $\mathrm{NF}_s(P,c)$: for $Q=Q(P)$, every square matrix $\Phi$ presenting the Selmer modules for $(s,d\circ P,Q)$ and
  every $F$ with $j_\iota(F)=\det\Phi$, some monomial of degree $<c$ has a coefficient in $F$ which is not in $\mathfrak a$.
\item $\mathrm{NF}^{\mathrm{adj}}_s(P,c)$: for $Q=Q(P)$ and every $\Phi$ as in (c), either $\Phi$ is empty or there are indices $b,b'$
  such that every $F$ with $j_\iota(F)=\operatorname{adj}(\Phi)_{bb'}$ has a unit coefficient at a monomial of degree $<c$
  in which no variable $X_{(1,j)}$ with $n\le j<R$ occurs.
\item A family $P'$ on $\iota_{R,n+1}$ \emph{extends} $P$ if it agrees with $P$ on $\iota_{R,n}$.
\item $\mathrm{Corank}_s(\varphi,\varphi',k)$, for $\varphi_a,\varphi'_a\in G_E$, states: for every discrete field $K$ which is an
  $\baselatoneO_0$-algebra killed by $\mathfrak m$ and onto which $\baselatoneO_0$ maps, there are $c_b\in K$ such that no $k+1$
  classes $x_a$ of the Selmer group of $T_s\otimes K$ which satisfy $\sum_bc_b\,\phi(x_a,y)(\varphi_b,\varphi'_b)=0$ for all
  classes $y$ of the Selmer group of $\rho^\vee_s\otimes K$ are linearly independent, where $\phi$ is the cup-product form of
  cocycles.
\item $\mathrm{Start}(s)$: for all $M,R,n,d,A,L,z$, $P$ on $\iota_{R,n}$ and $\varphi,\varphi'$ with $n\le R$, $4rn+2<\ell^M$, $z$ a tame
  inertia element at a twist prime, $L$ small, $\mathrm{PL}(P)$, $\varphi_a$ a Frobenius at $P_{(2,a)}$ with conjugate
  $\varphi'_a$, and $\mathrm{Corank}_s(\varphi,\varphi',1)$: $\mathrm{NF}^{\mathrm{adj}}_s(P,4rn+1)$.
\item $\mathrm{Grow}(s)$: if $\nu>0$, then for all $M,R,n,d,A,L,z$, $P$, an extension $P'$ of $P$ with $\mathrm{PL}(P')$, and $c$
  with $n<R$, $c+2r\le\ell^{M-\nu-1}$, $\nu+2\le M$, $z$ tame and $L$ $2$-small: $\mathrm{NF}_s(P,c)\Rightarrow\mathrm{NF}_s(P',c+2r)$
  and $\mathrm{NF}^{\mathrm{adj}}_s(P,c)\Rightarrow\mathrm{NF}^{\mathrm{adj}}_s(P',c+2r)$.
\item $\mathrm{Step}(s)$: if $\nu>0$, then for all $M,R,n$, admissible families $(H,S,d)$, $A,L,z,P,w$ with $n<R$,
  $2(e+1)w+2+2r\le\ell^{M-\nu-1}$, $\nu+2\le M$, $L\le H$, $z$ tame, $L$ $2$-small, $\mathrm{PL}(P)$ and $\mathrm{NF}^{\mathrm{adj}}_s(P,w+1)$,
  there is an extension $P'$ of $P$ with $\mathrm{PL}(P')$ and $\mathrm{NF}_s(P',2(e+1)w+2+2r)$.
\end{enumerate}
\formal{GlobalTheory.ResSel.Small}\formal{GlobalTheory.StepTwo.Small2}\formal{GlobalTheory.StepTwo.PiLocalization₂}%
\formal{GlobalTheory.StepTwo.NF}\formal{GlobalTheory.StepTwo.NFit}\formal{GlobalTheory.StepTwo.Ext}%
\formal{GlobalTheory.ResSel.Corank}\formal{GlobalTheory.StepTwo.Start}\formal{GlobalTheory.StepTwo.Grow}%
\formal{GlobalTheory.StepTwo.Step}
\end{definition}

Both $\mathrm{NF}_s(P,c)$ and $\mathrm{NF}^{\mathrm{adj}}_s(P,c)$ are monotone in $c$. Define $b_w(0)=w$ and
$b_w(k+1)=2(e+1)b_w(k)+2+2r$; $b_w$ is increasing in $k$.

\begin{lemma}[Iterated steps]\label{pf:baselat1:induct}
Let $\nu>0$, $t\in\N$ and $s_1,\dots,s_t$ signatures with $\mathrm{Grow}(s_j)$ and $\mathrm{Step}(s_j)$. Let $M,w$ with
$b_w(t)\le\ell^{M-\nu-1}$ and $\nu+2\le M$, $(H,S,d)$ admissible, $L\le H$ $2$-small for all $s_j$, $z$ tame, and $P_0$ a family on
$\iota_{R_0+t,R_0}$ with $\mathrm{PL}(P_0)$ and $\mathrm{NF}^{\mathrm{adj}}_{s_j}(P_0,w+1)$ for all $j$. Then for every $k\le t$ there is a family $P$ on
$\iota_{R_0+t,R_0+k}$ which agrees with $P_0$ on the indices of the first kind, with $\mathrm{PL}(P)$, $\mathrm{NF}^{\mathrm{adj}}_{s_j}(P,b_w(k)+1)$
for all $j$, and $\mathrm{NF}_{s_j}(P,b_w(k)+1)$ for all $j\le k$.
\formallabel{GlobalTheory.StepTwo.TwoBound.induct}
\end{lemma}

\begin{proof}
Induction on $k$; for $k=0$ take $P=P_0$. Let $P$ be as in the statement for $k<t$. Since $R_0+k<R_0+t$,
$2(e+1)b_w(k)+2+2r=b_w(k+1)\le b_w(t)\le\ell^{M-\nu-1}$ and $\mathrm{NF}^{\mathrm{adj}}_{s_{k+1}}(P,b_w(k)+1)$, $\mathrm{Step}(s_{k+1})$ gives an
extension $P'$ with $\mathrm{PL}(P')$ and $\mathrm{NF}_{s_{k+1}}(P',b_w(k+1))$. For every $j$, $\mathrm{Grow}(s_j)$ with $c=b_w(k)+1$ (note
$c+2r\le b_w(k+1)\le\ell^{M-\nu-1}$) carries $\mathrm{NF}^{\mathrm{adj}}_{s_j}(P,c)$ and, for $j\le k$, $\mathrm{NF}_{s_j}(P,c)$ to $P'$ with the bound
$c+2r\le b_w(k+1)+1$. By monotonicity $P'$ has the required properties for $k+1$.
\end{proof}

\begin{proposition}[Assembly of the coefficient bound]\label{pf:baselat1:thmCparts}
Assume that for every $i$ the statements $\mathrm{AP}(i)$ and $\mathrm{CH}(i)$ below hold, and that $\mathrm{Start}(s)$,
$\mathrm{Grow}(s)$ and $\mathrm{Step}(s)$ hold for every signature $s$ with $\mathrm{GH}(i,s)$. Then $(\mathrm{S}_4)$.
\begin{itemize}
\item $\mathrm{AP}(i)$: for all $M\ge1$, admissible $(H,S,d)$, tame $z$, open normal $U$, finite $A$, $R$ and
  $h_1,\dots,h_R\in H$ fixing the roots of unity of order $2\ell^M$, there is a family $P$ on $\iota_{R,R}$ with
  $\mathrm{Aux}(i,M,A,P,d\circ P)$, $\mathrm{Adp}(i,M,U,z,P,d\circ P)$, with the inertia groups at all places over all $P_q$
  contained in $U$, and with a Frobenius element $\sigma_a$ at $P_{(2,a)}$ conjugate to $h_a$ modulo $U$ for every $a$.
\item $\mathrm{CH}(i)$: for every set $\mathcal T$ of signatures with $\mathrm{GH}$ there is $R_0$ such that for every
  tame $z$ and every open normal $L$ which is small for all $s\in\mathcal T$, there are $h_1,\dots,h_{R_0}\in L$ and an open
  normal $U\le L$ such that $\mathrm{Corank}_s(\varphi,\varphi',k_\nu)$ for all $s\in\mathcal T$ whenever each $\varphi_a$ has
  conjugate $\varphi'_a$ and is conjugate to $h_a$ modulo $U$; here $k_\nu=0$ if $\nu=0$ and $k_\nu=1$ otherwise.
\end{itemize}
\formallabel{GlobalTheory.StepTwo.TwoBound.thmC_of_parts}
\end{proposition}

\begin{proof}
Let $\nu>0$ and $\mathcal T$ a set of signatures with $\mathrm{GH}$; let $R_0$ be given by $\mathrm{CH}(i)$. The set of
signatures is finite; enumerate $\mathcal T=\{s_1,\dots,s_t\}$. As $\ell\ne0$ in $\baselatoneO_0$, $e$ is finite. We prove
$\mathrm{CB}(i,e)$ with $R=R_0+t$, $D_d=b_w(t)$ for $w=4rR_0$, and $M_0=b_w(t)+\nu+3$. Let $M\ge M_0$, $(H,S,d)$ admissible, $z$
tame, $H'$ open normal and $A$ finite; then $4rR_0+2<\ell^M$ and $b_w(t)\le\ell^{M-\nu-1}$.

By \lref{GlobalTheory.ResSel.SmallSub.exists_small}, which constructs an open normal subgroup $L\le H$ which is
$2$-small for all $s\in\mathcal T$, we fix such an $L$. $\mathrm{CH}(i)$ gives $h_1,\dots,h_{R_0}\in L$ and $U\le L$. Extend
$h$ by $1$ to $R_0+t$ indices; these elements lie in $L\le H$ and fix the roots of unity of order $2\ell^M$.
$\mathrm{AP}(i)$, applied with $U'=U\cap H'$, gives a family $P$ on $\iota_{R,R}$ and Frobenius elements $\varphi_a$ at
$P_{(2,a)}$ with $g_a\varphi_ag_a^{-1}\in U'h_a\subseteq L$ for some $g_a$; as $L$ is normal, $\varphi_a\in L$. The restriction $P_0$ of
$P$ to $\iota_{R,R_0}$ satisfies $\mathrm{PL}(M,d,A,L,z,P_0)$: the auxiliary conditions are inherited, inertia at the places over
$P_q$ lies in $U'\le L$, the Frobenius condition at $P_{(1,a)}$ follows from $\mathrm{Adp}$, the decomposition group at
$P_{(2,a)}$ lies in $L$ because inertia lies in $L$ and a Frobenius $\varphi_a$ lies in the open subgroup $L$, and the values
$d_{P_{(1,b)}}(\varphi_a)$ are those prescribed by $\mathrm{Adp}$. By the choice of $U$, $\mathrm{Corank}_s(\varphi,\varphi',1)$ for
$s\in\mathcal T$, where $\varphi'_a$ is the conjugate of $\varphi_a$, and $\mathrm{Start}(s)$ gives $\mathrm{NF}^{\mathrm{adj}}_s(P_0,4rR_0+1)$.
\Cref{pf:baselat1:induct} with $k=t$ gives a family $P'$ on $\iota_{R,R}$ which agrees with $P$ on the indices of the first
kind, with $\mathrm{PL}(P')$ and $\mathrm{NF}_s(P',b_w(t)+1)$ for all $s\in\mathcal T$. Then $\mathrm{Aux}(i,M,A,P',d\circ P')$ holds,
$\mathrm{Adp}(i,M,H',z,P',d\circ P')$ holds (the Frobenius condition at $P'_{(1,a)}=P_{(1,a)}$ modulo $U'\le H'$, and the values of
$d$ at the Frobenius elements at the $P'_{(2,a)}$ by $\mathrm{PL}(P')$), and for every $s\in\mathcal T$, $\Phi$ presenting the
Selmer modules and $F$ with $j_\iota(F)=\det\Phi$ there is a monomial of degree $\le b_w(t)$ whose coefficient $x$ is not in
$\mathfrak a=\ell^\nu\mathfrak m$. In a discrete valuation ring this gives $v(x)\le v(\ell^\nu)=e$: otherwise $x$ would be
divisible by $\ell^\nu\varpi$ with $\varpi$ a uniformiser
(\lref{GlobalTheory.StepTwo.TwoBound.addVal_le_of_notMem}).
\end{proof}

\begin{lemma}[The new prime]\label{pf:baselat1:newprime}
Assume (E1)--(E11) and $(\mathrm{S}_1)$. Then $\mathrm{Step}(s)$ holds for every $i$ and every signature $s$ with
$\mathrm{GH}(i,s)$. The hypotheses (E3) and (E10) are not used.
\formallabel{GlobalTheory.NewPrime.step}
\end{lemma}

\begin{proof}
Let $\nu>0$ and the data be as in $\mathrm{Step}(s)$; put $c_0=2(e+1)w+2$, so $c_0+2r\le\ell^{M-\nu-1}$, and let
$\Lambda=\baselatoneLam_{M,\iota_{R,n}}$.

\emph{Case 1: $\mathrm{NF}_s(P,c_0)$ holds.} By
\lref{GlobalTheory.T.exists_newPlace_frobeniusConjugate_modOpenNormalSubgroup}, which states that under (E1), for
$2$-small $L$, $\mathrm{PL}(P)$, $n<R$, an admissible family with $L\le H$, an open normal $U\le L$ and $h\in L$ (with the
conditions on $d$ at $h$), there is an extension $P'$ of $P$ with $\mathrm{PL}(P')$ whose new place carries a Frobenius element
conjugate to $h$ modulo $U$, applied with $U=L$ and $h=1$, we obtain $P'$. Then $\mathrm{Grow}(s)$
(\lref{GlobalTheory.T.chebotarevInductionGrowBound_of_presentation}, from (E4) and $(\mathrm{S}_1)$) gives
$\mathrm{NF}_s(P',c_0+2r)$.

\emph{Case 2.} Otherwise there are $Q=Q(P)$, a square matrix $\Phi$ of size $a$ presenting the Selmer modules and $F_0$ with
$j_\iota(F_0)=\det\Phi$ such that all coefficients of $F_0$ in degrees $<c_0$ lie in $\mathfrak a$. If $a=0$ then
$\det\Phi=1$; the evaluation $\Lambda\to\baselatoneO_0/\mathfrak a$ with all variables $\mapsto0$
(\lref{GlobalTheory.T.exists_lambda_algebra_hom_evaluation}, which provides, for an $\baselatoneO_0$-algebra $B$ killed by
$\mathfrak a$ and elements $\xi_q\in B$ with $\xi_q^{\ell^{M-\nu-1}}=0$, an algebra map $\varphi\colon\Lambda\to B$ with
$\varphi(j_\iota(F))=F(\xi)$ and, for characters $d$, a character $\kappa$ of $G_E$ with values in $B^\times$ which is the image of the
universal character) maps $\det\Phi$ to $1$ and $j_\iota(F_0)$ to the constant term of $F_0$ modulo $\mathfrak a$, which is $0$;
hence $1\in\mathfrak a\subseteq\mathfrak m$, a contradiction. Hence by $\mathrm{NF}^{\mathrm{adj}}_s(P,w+1)$ there are $b,b'$ such that every lift of
$\operatorname{adj}(\Phi)_{bb'}$ has a unit coefficient at a monomial of degree $\le w$ avoiding the variables $X_{(1,j)}$, $n\le j<R$;
let $\mathcal K$ be the set of these indices and $F_1$ a lift of $\operatorname{adj}(\Phi)_{bb'}$ ($j_\iota$ is surjective,
\lref{GlobalTheory.TwSys.toLam_surjective}).

\emph{A finite test ring.} Since the residue field of $\baselatoneO_0$ is finite,
\lref{GlobalTheory.T.exists_finiteLocalRing_poly_divides_var_pow} gives a finite local $\baselatoneO_0$-algebra $A$ killed by
$\mathfrak a$, $t\in A$ and $a_q\in A$ with $t^{c_0}=0$, $\ell^\nu t^{c_0-1}\ne0$, $a_q=0$ for $q\in\mathcal K$, $a_q\in A^\times$ for
$q\notin\mathcal K$, and $F_1((a_qt)_q)\mid t^{(e+1)w}$. Let $\varphi\colon\Lambda\to A$ and $\kappa\colon G_E\to A^\times$ be the
evaluation at $\xi_q=a_qt$ and the corresponding character (the hypothesis $\xi_q^{\ell^{M-\nu-1}}=0$ holds as $c_0\le
\ell^{M-\nu-1}$); $\kappa$ is trivial on every $g$ with $d_{P_q}(g)=0$ for all $q\notin\mathcal K$. By
\lref{GlobalTheory.T.setup_and_antiHom_of_kills} the twisted setup for $(s,\kappa,Q)$ holds, $\kappa$ is anti-invariant
under the conjugation $c$, and $\kappa(z)=1$; by \lref{GlobalTheory.T.exists_universalSelmerCocycle_baseChange} the matrix $\Psi=\varphi(\Phi)$
presents the Selmer module of the twist by $\kappa$, with universal cocycle data $Z$.

\emph{The determinant and the adjugate.} $\det\Psi=\varphi(j_\iota(F_0))=F_0((a_qt)_q)=0$ by
\lref{GlobalTheory.T.aeval_eq_zero_of_low_coeff_of_nilpotent_multiple} (all $\xi_q$ are multiples of $t$, $t^{c_0}=0$, and
the coefficients of degree $<c_0$ vanish in $A$), and $\operatorname{adj}(\Psi)_{bb'}=F_1((a_qt)_q)$ divides $t^{(e+1)w}$. Hence
the column $x_0=\operatorname{adj}(\Psi)_{\cdot b'}$ satisfies $\Psi x_0=0$.

\emph{A dual Selmer class.} Let $\rho^\vee_s$ be the dual representation. By
\lref{GlobalTheory.T.perfect_pairing_finite_module_rows_criterion} (applied with $n=\ell^{\nu+1}$, which vanishes in $A$)
there is a finite $A$-module $Y$ with a perfect pairing $\beta$ into $\Z/\ell^{\nu+1}$ and the row criterion. The twisted dual has no
invariants (\lref{GlobalTheory.T.noInvDual_of_setup_of_isDual_of_tame_inertia}). Put $c_1=\ell^\nu t^{(e+1)w}$. The row
$y_0=\operatorname{adj}(\Psi)_{b\cdot}$ satisfies $y_0\Psi=0$ (as $\operatorname{adj}(\Psi)\Psi=\det\Psi=0$), and
$c_1\,t\,\operatorname{adj}(\Psi)_{bb'}\ne0$: otherwise, writing $t^{(e+1)w}=\operatorname{adj}(\Psi)_{bb'}\,u$, we would get
$\ell^\nu t^{2(e+1)w+1}=0$. The element $t$ kills the local invariants of the twisted dual at the split places and at the
places of $Q$ (\lref{GlobalTheory.T.local_invariants_killed_alt}, from (E4)). Hence
\lref{GlobalTheory.T.exists_dual_selmer_class_not_killed}, which also uses the order formula for the dual Selmer group
(\lref{GlobalTheory.T.selmer_order_formula_dual_selmer}, from (E2), (E5)--(E9) and (E11)), gives a dual Selmer cocycle
$y$ for $\rho^\vee_s$ twisted by $\kappa^{-1}$ such that $c_1y$ is not a coboundary.

\emph{Cocycle form and the new prime.} Let $L'=L\cap\bigcap_{q\notin\mathcal K}\ker d_{P_q}$; it is $2$-small
(\lref{GlobalTheory.NewPrime.small_inf}) and $\kappa$ is trivial on $L'$. Let $x$ be the cocycle of the Selmer module
corresponding to $x_0$. Since $x_0$ has the entry $\operatorname{adj}(\Psi)_{bb'}$, which divides $t^{(e+1)w}$, and $c_1y$ is not a coboundary,
\lref{GlobalTheory.T.exists_smallSubgroup_cocycleForm_ne_zero} gives $h\in L'$ at which the cup-product form of $x$
and $y$ does not vanish. By \lref{GlobalTheory.T.exists_open_normal_subgroup_two_cocycles} there is an open normal
$U\le L'$ which contains every $g\in L'$ at which $x$, $y$ and their conjugates by $c$ vanish, and such that the form of $x$
and $y$ does not vanish at any element conjugate to $h$ modulo $U$.
The Chebotarev statement of Case 1, applied with $U$ and $h$ (its conditions on the inertia groups at the places $P_{(1,c)}$
follow from \lref{GlobalTheory.T.inertiaVanishing_noInvariants_firstHalfPlaces}), gives an extension $P'$ with
$\mathrm{PL}(P')$ whose new place $p'=P'_{(2,n)}$ carries a Frobenius element $\sigma$ conjugate to $h$ modulo $U$ with
$d_{p'}(\sigma)=0$; then $\sigma\in L'$, and the decomposition group at $p'$ lies in $L'$
(\lref{GlobalTheory.ChosenPrimes.dec_le}).

\emph{$\mathrm{NF}_s(P',c_0+2r)$.} Suppose not: there are $Q'=Q(P')$, $\Phi'$ presenting the Selmer modules for $P'$ and $F'$ with
$j(F')=\det\Phi'$ whose coefficients of degree $<c_0+2r$ lie in $\mathfrak a$. By
\lref{GlobalTheory.T.exists_ring_sq_zero_nilpotent_augmented} there is an $A$-algebra $B$ with a retraction $\pi\colon B\to A$
and $x_1,x_2\in\ker\pi$ with $x_1^2=0$, $x_2^{2r}=0$, $\ker\pi\subseteq(x_1,x_2)$, and $\pi(f)=0$ whenever $x_1x_2^{2r-1}f=0$.
Evaluate at $\xi_{(2,n)}=x_1$ (the new variable), $\xi_{(1,n)}=x_2$ and $\xi_q=a_qt$ otherwise; this gives $\varphi'$, a character
$\kappa'$ and base-changed universal cocycle data $Z'$ for $\Psi'=\varphi'(\Phi')$. The character $\kappa'$ reduces to $\kappa$ under
$\pi$ (a direct computation with the values $\xi_q$), is anti-invariant, its values on the inertia group at $p'$ are
determined by $x_1$, and $\kappa'(\sigma)=1+x_2$; moreover $\kappa$ is trivial on the decomposition group at $p'$. These are the
hypotheses of \lref{GlobalTheory.NewPrime.det_ne_zero}, which, using (E2), (E4), the dual cocycle $y$, the
non-vanishing of the form of $x$ and $y$ at $\sigma$, the pairing $\beta$ and the retraction $\pi$, shows $\det\Psi'\ne0$. On the other hand
$\det\Psi'=F'(\xi)=0$ by \lref{GlobalTheory.T.aeval_eq_zero_low_coeff_three_nilpotents}, because every $\xi_q$ is a
multiple of $t$ or equal to $x_1$ or $x_2$, $t^{c_0}=0$, $x_1^2=0$, $x_2^{2r}=0$, and the coefficients of degree $<c_0+2r$ vanish in
$B$. This contradiction proves $\mathrm{NF}_s(P',c_0+2r)$.
\end{proof}

\begin{theorem}[The coefficient bound for positive exponent]\label{pf:baselat1:boundtwo}
(E1)--(E11) imply $(\mathrm{S}_4)$.
\formallabel{GlobalTheory.T.boundModTwoUniformizerAtTwo_of_chebotarev}
\formallabel{GlobalTheory.StepTwo.Dyadic.bound_mod_two_of_hypothesis}
\end{theorem}

\begin{proof}
By \lref{GlobalTheory.SelFun.oneMatrix_of_galoisCohomology} $(\mathrm{S}_1)$ holds. For $i$ and $s$ with $\mathrm{GH}(i,s)$, the
residual Selmer groups of $T_s$ and of its dual have the same dimension over every residue field $K$
(\lref{GlobalTheory.T.residual_selmer_finrank_eq_dual}), and the local cohomology at the places of the first kind of
a positioned family vanishes (\lref{GlobalTheory.ScalarFrob.firstHalf}, from (E4)--(E7)). We apply
\cref{pf:baselat1:thmCparts}:
\begin{itemize}
\item $\mathrm{AP}(i)$ holds by Chebotarev's theorem (E1) (\lref{GlobalTheory.ChosenPrimes.stepAB});
\item $\mathrm{CH}(i)$ holds by \lref{GlobalTheory.ChooseForms.chooseH}, from (E8) and, for every $s$ with
  $\mathrm{GH}(i,s)$, the two spanning statements for cup-product forms of residual cocycles
  (\lref{GlobalTheory.ResSel.SpanOddPf.spanOdd}, \lref{GlobalTheory.ResSel.AltSpan.spanTwo_of_galHyp}) and the
  equality of dimensions;
\item $\mathrm{Start}(s)$ holds by \lref{GlobalTheory.StepTwo.StartInd.start_of_resBound}, from the vanishing of local
  cohomology at a first place (\lref{GlobalTheory.ScalarFrob.firstPlace}, from (E5)--(E7)) and the residual bound
  (\lref{GlobalTheory.ResOrder.resBound}, from (E4), the vanishing at the first half of the places, the reciprocity
  statement \lref{GlobalTheory.ResRecip.reciprocity} and the equality of dimensions);
\item $\mathrm{Grow}(s)$ holds by \lref{GlobalTheory.Growth.grow} (from (E4) and $(\mathrm{S}_1)$);
\item $\mathrm{Step}(s)$ holds by \cref{pf:baselat1:newprime}.
\end{itemize}
Hence $(\mathrm{S}_4)$. The first formal statement listed is a restatement of the second.
\end{proof}


%======================================================================================================
\subsection{Proof of the main theorem}\label{pf:baselat1:ssec:main}

\begin{theorem}\label{pf:baselat1:main}
Let $h_L$ be a proof of $\sfname{RecTExists}$, and assume the hypothesis packages $\mathcal H_A$, $\mathcal H_E$,
$\mathcal H_M$ and $\mathcal H_L(h_L)$ (\cref{pf:baselat1:def:packages}). Let $N,N_g$ be norm-choice extensions such
that for every configuration and every norm-choice core $c$:
\begin{enumerate}[label=(\alph*)]
\item an element of $N(\dots,c)$ yields an element of the law extension $N_{\mathrm{law}}(\dots,c)$;
\item an element of $N_g(\dots,c)$ yields an element of $N(\dots,c)$ and the goodness condition $\mathrm{Gd}_0$ for
  $(w.D(Z),\mathrm{nch},x_d)$.
\end{enumerate}
Assume moreover:
\begin{enumerate}[label=(\alph*),resume]
\item starting witnesses eventually have level three for $\baselatoneS_N$:
  $\mathrm{Step}^{\mathrm w}(\mathrm{St},\mathrm H_3(\baselatoneS_N))$;
\item the level-three movement conditions $\mathrm{Mv}_3(\baselatoneS_N)$ and $\mathrm{Mv}_3(\baselatoneS_{N_g})$;
\item the position movement condition for $\baselatoneS_N$;
\item the second pass from built cores to starting witnesses of level three for $\baselatoneS_{N_g}$:
  $\mathrm{Pass}\bigl(\kappa_0,\widehat{\mathrm{Bu}}(\baselatoneS_N),\ (K,B,f,T)\mapsto\exists w,\ \mathrm{St}(w)\wedge
  \mathrm H_3(\baselatoneS_{N_g})(w)\bigr)$;
\item the multiplicity condition for $\baselatoneP_N$.
\end{enumerate}
Then $\sfname{MemberE}(\baselatoneS_N)$ and
$\sfname{ExBL}\bigl(\kappa_0,\widehat{\mathrm{Bu}}(\baselatoneS_N),\widehat{\mathrm{Gd}}(\baselatoneS_N),
\mathrm{Lk}(\baselatoneP_N)\bigr)$.
\formallabel{GlobalTheory.Mem.Top.Bl.exists_baseLattice_of_pair_alt}
\end{theorem}

\begin{proof}
\emph{Export.} \Cref{pf:baselat1:export}, applied to the members of $\mathcal H_E$, to rapid decay of cusp forms and
global Artin reciprocity from $\mathcal H_A$, to (L1), (L2) and (L6) of
$\mathcal H_L(h_L)$, to $N$ with $S$ trivial and $\Gamma$ the closure of the true geometric predicate, and to (a), gives
(M3) for $\baselatoneS_N$.

\emph{Member existence.} \Cref{pf:baselat1:closing} applies: its published hypotheses are members of $\mathcal H_A$
(ray class fields, principal ideals, CM curves, discrete-series statements, Harish-Chandra finiteness, rapid decay (twice),
uniqueness of discrete series, finite-dimensionality of rows, Siegel--Weil, formal degrees, the mass ratio, the Weil
representation, Casselman--Wallach and Harish-Chandra modules); its hypotheses (a)--(e) are (c), (d), the export statement
just proved, (f) and (b). Hence $\sfname{MemberE}(\baselatoneS_N)$.

\emph{Bundles.} By \lref{GlobalTheory.T.geometric_parameter_uniform_alt}, which states that bundles linked through
weighted points exist for the parameter of every norm-choice extension whose data yield law data, provided the position
movement and multiplicity conditions hold (and Green's cuspidal characters, the discrete series cover and published
automorphic and class-field inputs), applied with the members of $\mathcal H_M$ and $\mathcal H_A$, to $N$, (a), (e) and
(g), we obtain $\sfname{ExBL}(\kappa_0,\widehat{\mathrm{Bu}}(\baselatoneS_N),\widehat{\mathrm{Gd}}(\baselatoneS_N),
\mathrm{Lk}(\baselatoneP_N))$.
\end{proof}

%======================================================================================================
\subsection{Remarks on the formal statements}\label{pf:baselat1:ssec:remarks}

\begin{remark}\label{pf:baselat1:rem:summary}
\begin{enumerate}[label=(\arabic*)]
\item \emph{Names.} The conclusion of \cref{pf:baselat1:main} is member existence and bundle existence for the parameter
  $\baselatoneP_N$; no lattice and no prime occurs in it. The theorem proved in \cref{pf:baselat1:export} is an export
  statement; no automorphy statement occurs. The existence statement of \cref{pf:baselat1:specgap}, on which \cref{pf:baselat1:spack} rests, has a
  formal name referring to local Bessel uniqueness; no Bessel model occurs in it (\cref{pf:baselat1:rem:bessel}).
\item \emph{Unused hypotheses.} The following hypotheses are not used in the proofs: the mass ratio in
  \cref{pf:baselat1:core} and in \cref{pf:baselat1:has4} (where also $h_L$ is unused); Casselman--Wallach and
  Harish-Chandra modules in \cref{pf:baselat1:theta}; ray class fields, principal ideals and CM curves in
  \cref{pf:baselat1:div}; the Greenberg--Wiles formula in the proof of the twist theorem
  (\cref{pf:baselat1:rem:gw}); cohomological dimension two (E3) and unramified orthogonality (E10) in
  \cref{pf:baselat1:newprime}; the finiteness of the residue field in \cref{pf:baselat1:orderEval}. In
  \cref{pf:baselat1:main}, of the local Langlands package $\mathcal H_L(h_L)$ only (L1), (L2) and (L6) are used
  ((L3)--(L5) are not), and of $\mathcal H_M$ only the discrete series cover (H1) and, among the
  theta-condition theorems (H2), Green's description of cuspidal characters (both are passed to the cited bundle
  theorem); (H3)--(H5) are not used.
\item \emph{True components.} The standing hypotheses $\mathrm{SH}(i)$ of \cref{pf:baselat1:def:exportAt} and the
  Galois hypotheses of \cref{pf:baselat1:def:galside} each contain the true proposition as a conjunct. The spectral
  extension of $\baselatoneP_N$ is the one-element type and its geometric predicate holds identically, so these two
  components of the slot family impose no condition beyond the $\lambda$-condition in $O_x$.
\item \emph{Duplicated hypothesis.} Rapid decay of cusp forms occurs twice among the hypotheses of
  \cref{pf:baselat1:closing}.
\item \emph{Strength.} All statements of \cref{pf:baselat1:ssec:calculus,pf:baselat1:ssec:stages} are ``eventual''
  statements over scale contexts and ``large'' data; their content depends on the existence of contexts, bases and
  twisting data with the required properties, which is not addressed in this section.
\end{enumerate}
\end{remark}

